\documentclass[12pt,a4paper]{article}

% ---------------------------------------------------
% GRUNDSPRACHE & SCHRIFT
% ---------------------------------------------------
\usepackage[ngerman]{babel}

\usepackage{fontspec}
\setmainfont{Times New Roman}
\usepackage{enumitem}

% ---------------------------------------------------
% SEITENLAYOUT
% ---------------------------------------------------
\usepackage{geometry}
\geometry{
    left=4cm,
    right=2cm,
    top=4cm,
    bottom=2cm,
    headsep=1cm,
    footskip=2cm
}

% ---------------------------------------------------
% FORMATIERUNG
% ---------------------------------------------------
\usepackage{setspace}
\onehalfspacing
\setlength{\parskip}{6pt}
\setlength{\parindent}{0pt}
% Zeilenumbruch-Notreserve gegen Overfull-Boxen bei langen Komposita
\emergencystretch=1em

% ---------------------------------------------------
% FUSSNOTENFORMATIERUNG (einzeilig, 10pt)
% ---------------------------------------------------
\usepackage{etoolbox}
\appto{\footnotesize}{\fontsize{10pt}{12pt}\selectfont}
\renewcommand{\footnotesize}{\fontsize{10pt}{12pt}\selectfont}

\makeatletter
\renewcommand\@makefntext[1]{%
  \noindent\makebox[1.8em][l]{\@makefnmark}%
  \begin{spacing}{1}{\fontsize{10pt}{12pt}\selectfont #1}\end{spacing}}
\makeatother

\usepackage{titlesec}
\titlespacing*{\section}{0pt}{12pt}{6pt}
\titlespacing*{\subsection}{0pt}{12pt}{6pt}
\titlespacing*{\subsubsection}{0pt}{12pt}{6pt}

\titleformat{\section}
  {\normalfont\bfseries\normalsize}{\thesection}{1em}{}

\titleformat{\subsection}
  {\normalfont\bfseries\normalsize}{\thesubsection}{1em}{}

\titleformat{\subsubsection}
  {\normalfont\bfseries\normalsize}{\thesubsubsection}{1em}{}

% Gliederungstiefe: bis Ebene 3 (x.y.z) nummerieren und ins Inhaltsverzeichnis
\setcounter{secnumdepth}{3}
\setcounter{tocdepth}{3}

\pretocmd{\section}{\clearpage}{}{}

% ---------------------------------------------------
% GRAFIKEN, TABELLEN, VERZEICHNISSE
% ---------------------------------------------------
% amsmath/amssymb: Voraussetzung fuer equation, align, \text, \operatorname.
% Das Formelverzeichnis (\formeleintrag) war eingerichtet, das Paket zum
% Formelschreiben fehlte -- ergaenzt am 18.08.2026.
\usepackage{amsmath}
\usepackage{amssymb}
\usepackage{graphicx}
\usepackage{booktabs}
\usepackage{listings}
\usepackage{xcolor}
\definecolor{codegruen}{rgb}{0.0,0.45,0.0}
\definecolor{codegrau}{rgb}{0.45,0.45,0.45}
\definecolor{codeblau}{rgb}{0.10,0.20,0.65}
\definecolor{codehintergrund}{rgb}{0.97,0.97,0.97}
\lstdefinestyle{pythonstil}{
   language=Python,
   escapeinside={(*@}{@*)},
   backgroundcolor=\color{codehintergrund},
   basicstyle=\ttfamily\footnotesize,
   keywordstyle=\color{codeblau}\bfseries,
   commentstyle=\color{codegruen}\itshape,
   stringstyle=\color{codegrau},
   numbers=left, numberstyle=\tiny\color{codegrau}, numbersep=8pt,
   frame=single, rulecolor=\color{codegrau},
   breaklines=true, showstringspaces=false, aboveskip=0pt, belowskip=0pt,
   literate={ä}{{\"a}}1 {ö}{{\"o}}1 {ü}{{\"u}}1 {ß}{{\ss}}1
            {Ä}{{\"A}}1 {Ö}{{\"O}}1 {Ü}{{\"U}}1
}
\lstset{style=pythonstil}
% Quellcode-Beschriftung: siehe \floatname im Block QUELLCODE weiter unten
\renewcommand{\topfraction}{0.9}
\renewcommand{\bottomfraction}{0.7}
\renewcommand{\textfraction}{0.07}
\renewcommand{\floatpagefraction}{0.75}
\setcounter{topnumber}{3}\setcounter{bottomnumber}{2}\setcounter{totalnumber}{4}
\usepackage{placeins}
\pretocmd{\subsection}{\FloatBarrier}{}{}
% ---------------------------------------------------
% VERZEICHNISFORMATIERUNG (tocloft)
% ---------------------------------------------------
\usepackage{tocloft}
\tocloftpagestyle{fancy}

% Überschriften der Verzeichnisse: fett, 12pt, schwarz
\renewcommand{\cfttoctitlefont}{\bfseries\normalsize}
\renewcommand{\cftlottitlefont}{\bfseries\normalsize}
\renewcommand{\cftloftitlefont}{\bfseries\normalsize}

\setlength{\cftafterlottitleskip}{0pt}
\setlength{\cftafterloftitleskip}{0pt}

\renewcommand{\cftsecfont}{\normalfont}
\renewcommand{\cftsecpagefont}{\normalfont}
\renewcommand{\cftsubsecfont}{\normalfont}
\renewcommand{\cftsubsecpagefont}{\normalfont}
\renewcommand{\cftsubsubsecfont}{\normalfont}
\renewcommand{\cftsubsubsecpagefont}{\normalfont}

\renewcommand{\cftsecdotsep}{\cftdotsep}
\renewcommand{\cftsubsecdotsep}{\cftdotsep}
\renewcommand{\cftsubsubsecdotsep}{\cftdotsep}

\setlength{\cfttabindent}{0pt}
\setlength{\cftfigindent}{0pt}

\renewcommand{\cfttabdotsep}{\cftdotsep}
\renewcommand{\cftfigdotsep}{\cftdotsep}
\renewcommand{\cftdotsep}{1}

% Präfix "Tabelle X:" bzw. "Abbildung X:" in den Verzeichnissen
\renewcommand{\cfttabpresnum}{Tabelle }
\renewcommand{\cfttabaftersnum}{:}
\setlength{\cfttabnumwidth}{5em}
\renewcommand{\cfttabfont}{\normalfont}
\renewcommand{\cfttabpagefont}{\normalfont}

\renewcommand{\cftfigpresnum}{Abbildung }
\renewcommand{\cftfigaftersnum}{:}
\setlength{\cftfignumwidth}{6em}
\renewcommand{\cftfigfont}{\normalfont}
\renewcommand{\cftfigpagefont}{\normalfont}

\renewcommand{\cftbeforetabskip}{12pt}
\renewcommand{\cftbeforefigskip}{12pt}

% ---- FORMELVERZEICHNIS -----------------------------
% Eigenes Verzeichnis via tocloft. Nutzung im Text:
%   \begin{equation} ... \end{equation}
%   \formeleintrag{Zielfunktion der Ridge Regression}
\newlistof{formeln}{equ}{Formelverzeichnis}
\renewcommand{\cftequtitlefont}{\bfseries\normalsize}
\setlength{\cftbeforeequtitleskip}{0pt}
\setlength{\cftafterequtitleskip}{0pt}
\renewcommand{\cftformelnpresnum}{Formel }
\renewcommand{\cftformelnaftersnum}{:}
\setlength{\cftformelnnumwidth}{5em}
\setlength{\cftformelnindent}{0pt}
\renewcommand{\cftformelnfont}{\normalfont}
\renewcommand{\cftformelnpagefont}{\normalfont}
\renewcommand{\cftformelndotsep}{\cftdotsep}
\renewcommand{\cftbeforeformelnskip}{12pt}
\newcommand{\formeleintrag}[1]{%
  \addcontentsline{equ}{formeln}{\protect\numberline{\theequation}#1}}

\usepackage[nottoc]{tocbibind}
\usepackage{float}
\usepackage{tabularx}
\usepackage{longtable}
\usepackage[justification=raggedright,singlelinecheck=false]{caption}
\captionsetup{labelfont=bf, textfont=bf}
\captionsetup[figure]{skip=6pt}   % Abstand Abbildung  -> Unterschrift

% ---------------------------------------------------
% EINHEITLICHE TABELLENFORMATIERUNG
% ---------------------------------------------------
% Beschriftung: fester Abstand zur Tabelle, immer 12pt fett, unabhaengig
% vom Schriftgrad des Tabellenkoerpers.
\captionsetup[table]{skip=4pt, font=normalsize}

% Spaltentypen: Y = dehnbar linksbuendig, L{Breite} = fest linksbuendig.
\newcolumntype{Y}{>{\raggedright\arraybackslash}X}
\newcolumntype{L}[1]{>{\raggedright\arraybackslash}p{#1}}

% Rahmen jeder Tabelle -- ersetzt \centering, \small/\footnotesize,
% \singlespacing und alle lokalen \setlength-Zeilen.
\newcommand{\tabformat}{%
  \centering
  \singlespacing
  \footnotesize
  \setlength{\tabcolsep}{5pt}%
  \renewcommand{\arraystretch}{1.15}%
}

% Quellenzeile: 9pt, einzeilig, linksbuendig, 3pt unter \bottomrule.
% Aufruf: \tabquelle   bzw.   \tabquelle[Eigene Darstellung nach ...]
\newcommand{\tabquelle}[1][Eigene Darstellung]{%
  \par\setlength{\parskip}{0pt}\vspace{3pt}%
  {\setstretch{1}\fontsize{9pt}{11pt}\selectfont\raggedright Quelle: #1\par}%
}
% ---------------------------------------------------
% QUELLCODE ALS EIGENES GLEITOBJEKT
% ---------------------------------------------------
% Beschriftung oben, Code, Quellenzeile unten. Mit lstlisting allein nicht
% moeglich: unter einen gleitenden lstlisting-Block laesst sich nichts setzen.
\floatstyle{plaintop}
\newfloat{quellcode}{htbp}{lol}
\floatname{quellcode}{Code}
\captionsetup[quellcode]{skip=4pt, font=normalsize}
\newlistof{quellcodes}{lol}{Codeverzeichnis}
\makeatletter\let\l@quellcode\l@quellcodes\makeatother
\renewcommand{\cftloltitlefont}{\bfseries\normalsize}
\setlength{\cftbeforeloltitleskip}{0pt}
\setlength{\cftafterloltitleskip}{0pt}
\renewcommand{\cftquellcodespresnum}{Code }
\renewcommand{\cftquellcodesaftersnum}{:}
\setlength{\cftquellcodesnumwidth}{7em}
\setlength{\cftquellcodesindent}{0pt}
\renewcommand{\cftquellcodesfont}{\normalfont}
\renewcommand{\cftquellcodespagefont}{\normalfont}
\renewcommand{\cftquellcodesdotsep}{\cftdotsep}
\setlength{\cftbeforequellcodesskip}{12pt}

% ---------------------------------------------------
% LINKS
% ---------------------------------------------------
\usepackage[hidelinks]{hyperref}
\usepackage{csquotes}
% xstring liefert \IfBeginWith, das der @online-Treiber von
% biblatex-chicago benutzt. Ohne dieses Paket bricht das
% Literaturverzeichnis mit "Undefined control sequence".
\usepackage{xstring}

% ---------------------------------------------------
% FUSSNOTEN-FORMATIERUNG
% ---------------------------------------------------
\usepackage[hang,flushmargin]{footmisc}
\setlength{\footnotemargin}{1.5em}
\setlength{\footnotesep}{10pt}

% ---------------------------------------------------
% KOPF-/FUSSZEILEN
% ---------------------------------------------------
\usepackage{fancyhdr}
\pagestyle{fancy}
\fancyhf{}
\renewcommand{\headrulewidth}{0pt}
\fancyhead[C]{\thepage}
\setlength{\headheight}{15pt}
\setlength{\headsep}{20pt}

% ---------------------------------------------------
% LITERATURDATENBANK
% ---------------------------------------------------
% literatur.bib ist der Master -- fuer die Eintraege wie fuer das
% Register. Ueber jedem Eintrag stehen %%-Zeilen mit Verwendungsstelle,
% Zugang und der geprueften Seitenangabe; Biber ignoriert alles
% ausserhalb der @-Eintraege, die Angaben stoeren also nicht.
%
% Die Volltexte liegen im Projektordner unter Quellen/, benannt nach dem
% Bib-Key.
%
% REGEL FUER SEITENANGABEN: Zitiert wird nur eine Seite, die am Volltext
% nachgeschlagen wurde. Die %%-Zeile haelt Datum und Belegsatz fest,
% damit die Angabe spaeter nachvollziehbar bleibt, ohne die Quelle
% erneut lesen zu muessen. Eine maschinelle Pruefung gibt es NICHT --
% aeltere Fassungen dieses Kommentars nannten tools/quellen.py und
% docs/10_QUELLEN.md; beides existiert nicht.
%
% Seit dem 18.08.2026 ist die Datenbank eine eigene Datei statt eines
% filecontents-Blocks, sonst zwingt jede Quellenaenderung zum Hochladen
% der kompletten main.tex.

% ---------------------------------------------------
% BIBLATEX-KONFIGURATION (Chicago, einheitlich)
% ---------------------------------------------------
\usepackage[
    backend=biber,
    style=chicago-notes,
    doi=false,
    isbn=false,
    url=false
]{biblatex}
\addbibresource{literatur.bib}

\DefineBibliographyStrings{ngerman}{
    andothers = {et\addabbrvspace al\adddot},
}
\renewcommand*{\finalnamedelim}{\addspace\&\addspace}

\setlength{\bibitemsep}{0.5\baselineskip}
\setlength{\bibhang}{2em}
\renewcommand*{\bibfont}{\normalfont\singlespacing}

\DeclareFieldFormat{postnote}{S.~#1}
\DeclareFieldFormat{multipostnote}{S.~#1}

\renewcommand{\thefootnote}{\arabic{footnote}}

\makeatletter
\renewcommand\@makefntext[1]{%
  \noindent\makebox[1.8em][l]{\@makefnmark}#1%
}

\DeclareNameAlias{sortname}{family-given}
\ExecuteBibliographyOptions{maxnames=2, minnames=1}

\DeclareCiteCommand{\footcite}
  {}
  {%
    \footnote{%
      \printfield{prenote}%
      \iffieldundef{prenote}{}{\addspace}%
      \printnames[labelname]{labelname}\addspace%
      \mkbibparens{\printfield{year}}%
      \iffieldundef{postnote}
        {}
        {\addcomma\space\printfield{postnote}}%
    }%
  }
  {\multicitedelim}
  {}

\makeatother

\AtBeginDocument{%
  \renewcommand{\cfttabfont}{\normalfont}%
  \renewcommand{\cfttabpagefont}{\normalfont}%
  \renewcommand{\cftfigfont}{\normalfont}%
  \renewcommand{\cftfigpagefont}{\normalfont}%
}

% ---------------------------------------------------
% DOKUMENTBEGINN
% ---------------------------------------------------
\begin{document}

% ---------------------------------------------------
% TITELBLATT
% ---------------------------------------------------
\newgeometry{left=2.5cm, right=2.5cm, top=4cm, bottom=2cm}
\begin{titlepage}
\centering

\includegraphics[width=5cm]{FOM_Logo01_2012_rgb.jpg} \\[0.5cm]

\textbf{FOM Hochschule für Oekonomie \& Management} \\
Hochschulzentrum Bonn \\[1.0cm]

Berufsbegleitender Studiengang Wirtschaftsinformatik \\
zum \\
Bachelor of Science (B.Sc.) \\[1.5cm]

Bachelorarbeit zum Thema \\[1.5cm]

{\Large \textbf{„Vorhersage von Feuerwehreinsätzen auf Basis struktureller und
sozioökonomischer Merkmale -- ein Vergleich ausgewählter Machine-Learning-Methoden
am Beispiel von Stadtteildaten“}}

\vfill

\begin{flushleft}
\begin{tabular}{@{}ll}
1. Gutachter: & Prof. Dr. Julian Schröter \\
2. Gutachter: & Oliver Bach \\
Name: & Lukas Nießen \\
Matrikelnummer: & 728085 \\
Wortzahl: & XXXXX Wörter \\
Abgabedatum: & 07.10.2026 \\
\end{tabular}
\end{flushleft}

\end{titlepage}
\restoregeometry

% ---------------------------------------------------
% DISCLAIMER
% ---------------------------------------------------
%\newpage
%\pagenumbering{Roman}
%\setcounter{page}{2}
%\phantomsection
%\section*{Disclaimer}
%\addcontentsline{toc}{section}{Disclaimer}

%\textbf{KI-Nutzung:} In dieser wissenschaftlichen Arbeit wurden zur sprachlichen Unterstützung bei der Erstellung des Textes KI-Modelle eingesetzt. Insbesondere wurden Anthropic Claude und Google Gemini zur Verbesserung der Grammatik, Struktur, Rechtschreibung und Lesbarkeit verwendet. Die inhaltliche Konzeption, Argumentation und jedwede Schlussfolgerungen liegen jedoch vollständig in der Verantwortung des Autors. Sollten Fremdwerke verwendet worden sein, sind diese entsprechend durch Fußnoten gekennzeichnet.

%\textbf{Geschlechterinklusive Sprache:} Darüber hinaus wird in der vorliegenden Arbeit das generische Maskulinum zur besseren Lesbarkeit verwendet. Die in dieser Arbeit verwendeten Personenbezeichnungen beziehen sich, sofern nicht anders kenntlich gemacht, auf alle Geschlechter.

% ---------------------------------------------------
% VERZEICHNISSE
% ---------------------------------------------------
\pagenumbering{Roman}
\newpage
\tableofcontents

\newpage
\listoffigures

\newpage
\phantomsection
{\bfseries\normalsize Abkürzungsverzeichnis}
\vspace{6pt}
\addcontentsline{toc}{section}{Abkürzungsverzeichnis}

\begin{tabular}{@{}ll}
ACS & American Community Survey \\[12pt]
AUROC & Area Under the Receiver Operating Characteristic \\[12pt]
CRISP-DM & Cross-Industry Standard Process for Data Mining \\[12pt]
EDA & Explorative Datenanalyse \\[12pt]
GLM & Generalized Linear Model \\[12pt]
ICC & Intraclass Correlation Coefficient \\[12pt]
MAE & Mean Absolute Error \\[12pt]
NFIRS & National Fire Incident Reporting System \\[12pt]
RMSE & Root Mean Square Error \\[12pt]
SFFD & San Francisco Fire Department \\[12pt]
SFPD & San Francisco Police Department \\[12pt]
SHAP & SHapley Additive exPlanations \\[12pt]
\end{tabular}


\newpage
\listoftables

\newpage
\phantomsection
\addcontentsline{toc}{section}{Codeverzeichnis}
\listofquellcodes

\newpage
\phantomsection
\addcontentsline{toc}{section}{Formelverzeichnis}
\listofformeln

% ---- SYMBOLVERZEICHNIS ------------------------------
% Angelegt 28.08.2026. Auflage aus der Sprechstunde vom 24.08.: "Falls man
% Formeln hat braucht man auf jeden Fall auch ein Symbolverzeichnis."
% Aufbau bewusst wie das Abkuerzungsverzeichnis oben -- schlichtes tabular,
% keine automatische Liste, weil die Symbole von Hand gepflegt werden.
% Stand: alle Symbole aus Kapitel 2 (sechs Formeln). Beim Schreiben weiterer
% Kapitel MITFUEHREN, nicht am Ende nachbauen.
\newpage
\phantomsection
{\bfseries\normalsize Symbolverzeichnis}
\vspace{6pt}
\addcontentsline{toc}{section}{Symbolverzeichnis}

\begin{longtable}{@{}ll}
$y$ & Zielgröße einer Beobachtung \\[12pt]
$\mathbf{y}$ & Vektor der Zielgröße über alle Beobachtungen \\[12pt]
$\mathbf{x}$ & Vektor der erklärenden Merkmale einer Beobachtung \\[12pt]
$\mathbf{X}$ & Designmatrix, zeilenweise die Merkmalsvektoren aller Beobachtungen \\[12pt]
$\boldsymbol{\beta}$ & Vektor der zu schätzenden Koeffizienten, $\hat{\boldsymbol{\beta}}$ seine Schätzung \\[12pt]
$g(\cdot)$ & Linkfunktion eines verallgemeinerten linearen Modells \\[12pt]
$\mathbb{E}[\cdot]$ & Erwartungswert \\[12pt]
$\mu$ & bedingter Erwartungswert der Zielgröße \\[12pt]
$B$ & Wohnbevölkerung eines Stadtteils, Exposition der Poisson-Regression \\[12pt]
$\Pr(\cdot)$ & Wahrscheinlichkeit \\[12pt]
$K$ & Anzahl der Klassen der nominalen Zielgröße \\[12pt]
$n$ & Anzahl der Beobachtungen \\[12pt]
$p$ & Anzahl der Merkmale; in Testergebnissen der Überschreitungswert \\[12pt]
$\lambda$ & Strafgewicht der Ridge Regression \\[12pt]
$\mathbf{I}$ & Einheitsmatrix \\[12pt]
$\mathbf{v}$ & beliebiger Vektor im Nachweis der Regularität \\[12pt]
$J$ & Anzahl der Endknoten eines Entscheidungsbaums \\[12pt]
$\mathcal{L}$ & Zielfunktion des Gradient Boosting \\[12pt]
$l(\cdot)$ & Verlustfunktion einer einzelnen Beobachtung \\[12pt]
$\Omega(\cdot)$ & Strafterm für die Komplexität eines Baums \\[12pt]
$f_k$ & $k$-ter Baum eines Boosting-Modells \\[12pt]
$P_k$, $R_k$ & Präzision und Trefferquote der Klasse $k$ \\[12pt]
$F1_k$ & F1-Maß der Klasse $k$ \\[12pt]
$\mathcal{O}(\cdot)$ & Landau-Symbol für die Ordnung des Rechenaufwands \\[12pt]
\end{longtable}
% longtable zaehlt den Tabellenzaehler auch ohne Beschriftung hoch. Ohne die
% folgende Zeile beginnt das Tabellenverzeichnis bei "Tabelle 2".
\addtocounter{table}{-1}

% ---------------------------------------------------
% TEXTTEIL
% ---------------------------------------------------
\newpage
\pagenumbering{arabic}
\setcounter{page}{1}

% ===================================================
\section{Einleitung}
% Umfang: ca. 3 Seiten
\subsection{Problemstellung und Praxisrelevanz}
% Einstieg über das Praxisproblem: Ressourcenplanung und Präventionssteuerung
% der Feuerwehr, nicht über die Algorithmen.

% ---- ÜBERNAHME AUS EXPOSÉ (Kap. 1, Abs. 1--2) -- Tempus angepasst ----
% VORHER ERGÄNZEN: 1--2 Absätze zum Praxisproblem (Ressourcenplanung,
% Präventionssteuerung). Das Exposé steigt methodenseitig ein, die Thesis darf das nicht.

Wie viel eine Feuerwehr im Ernstfall ausrichtet, hängt an der Zeit bis zum
Eintreffen der Einsatzkräfte: Für Wohnungsbrände ist der Zusammenhang zwischen Hilfsfrist und
Todesopfern nachgewiesen.\footcite[706]{Jaldell2017} Diese Zeit entscheidet
sich vor dem Alarm, in der Vorhaltung von Wachstandorten, Fahrzeugen und
Personal und wo eine belastbare Einsatzhistorie fehlt, muss die dafür
nötige Erwartung künftiger Last aus den Strukturmerkmalen eines Gebiets
entstehen.

Ob das gelingt, hängt am Verfahren. Verglichen werden Ridge Regression, Random
Forest und XGBoost, drei Algorithmen mit grundlegend unterschiedlichem
methodischem Ansatz; welcher davon auf einer konkreten Datenkonstellation am
besten trifft, lässt sich theoretisch nicht bestimmen und muss empirisch
entschieden werden. Als Anwendungsfall dient die Vorhersage von
Feuerwehreinsätzen in den Stadtteilen von San Francisco: Einsatzdaten des San
Francisco Fire Department (SFFD), sozioökonomische Indikatoren der American
Community Survey (ACS), Kriminalitätskennzahlen und bauliche Strukturmerkmale,
zu einem stadtteilbezogenen Datenbestand zusammengeführt.
% ---- ENDE ÜBERNAHME ----

\subsection{Zielsetzung und Forschungsfragen}
% Zentrale Forschungsfrage + vier Unterfragen (UF1--UF4).

% ---- ÜBERNAHME AUS EXPOSÉ (Kap. 1, Abs. 3 + Forschungsfragen) -- nahezu 1:1 ----
% ACHTUNG: Dieser Absatz war im Exposé mit Fußnote 2 als KI-unterstützt gekennzeichnet
% (Prompt 1). Der Eintrag muss ins KI-Verzeichnis der Thesis mitwandern.

Die Arbeit verfolgt zwei aufeinander aufbauende Zielsetzungen. Zunächst wird geprüft, ob die
sozioökonomischen, kriminalitätsbezogenen und baulichen Faktoren in statistisch nachweisbarem
Zusammenhang mit der Einsatzlast stehen bzw. einen prädiktiven Beitrag leisten. Aufbauend darauf
werden zwei Stränge modelliert: zum einen die Menge der Einsatzlast, gemessen an der Anzahl
Einsätze pro Stadtteil und Monat und zusätzlich an der auf die Wohnbevölkerung normierten Rate
(Regression auf aggregierten Ereignishäufigkeiten), zum anderen ihre Zusammensetzung, gemessen an der je Stadtteil und Monat überwiegenden Einsatzart als kategorialer Zielgröße (Klassifikation). Auf der Einsatzhäufigkeit werden Ridge Regression,
Random Forest und XGBoost trainiert, auf der Einsatzart Random Forest und XGBoost. Die
unterschiedliche Verfahrenszahl folgt der Skalierung der jeweiligen Zielgröße und wird in
Kapitel 6 begründet. Für alle Verfahren werden die Hyperparameter optimiert und die
Prognosegüte anhand etablierter Gütemaße im Rahmen einer stadtteilweisen
Kreuzvalidierung gegenübergestellt. Die Aufteilung trennt dabei ganze
Stadtteile statt Zeiträume. Geprüft wird damit die Generalisierung auf unbekannte Stadtteile und
nicht die Fortschreibung eines bereits bekannten Stadtteils in die Zukunft.

% AUFLAGE SCHROETER 04.08.2026 (R-17, 03_STAND.md Abschnitt 6: "offen, zugesagt"):
% Die Zielsetzung der Validierung ist WOERTLICH als "Generalisierung auf unbekannte
% Stadtteile" zu formulieren. Der Begriff ist Schroeter gegenueber als umgesetzt
% gemeldet, steht aber weder hier noch in 5.4. Ein Satz an dieser Stelle, ein
% zweiter in 5.4 -- beide muessen die Wendung enthalten.

Daraus ergibt sich folgende zentrale Forschungsfrage:

\begin{quote}
Inwiefern lassen sich die Häufigkeit und die Art von Feuerwehreinsätzen in den Stadtteilen
San Franciscos durch sozioökonomische, kriminalitätsbezogene und bauliche Merkmale statistisch
beschreiben und vorhersagen, und welches der drei Verfahren -- Ridge Regression, Random Forest
oder XGBoost -- erzielt dabei die höchste Prognosegüte?
\end{quote}

Daraus leiten sich folgende Unterfragen ab:

\begin{itemize}[leftmargin=*, itemsep=2pt]
  \item Lassen sich statistisch nachweisbare Zusammenhänge zwischen sozioökonomischen,
        kriminalitätsbezogenen und baulichen Merkmalen und der stadtteilbezogenen
        Einsatzhäufigkeit bzw. Einsatzart feststellen?
  \item Wie unterscheiden sich Ridge Regression, Random Forest und XGBoost hinsichtlich ihrer
        Prognosegüte (RMSE/MAE/R² bzw. Macro-F1/Macro-AUROC) im Vergleich zu den Baselines,
        bewertet mittels stadtteilweiser Kreuzvalidierung?
  \item Wie verhalten sich die Modelle hinsichtlich Trainings- und Inferenzaufwand?
  \item Welche Implikationen ergeben sich aus den Ergebnissen für die Modellauswahl bei
        vergleichbaren tabellarischen Prognoseaufgaben?
\end{itemize}
% ---- ENDE ÜBERNAHME ----

\subsection{Aufbau der Arbeit}
% Gliederung in Fließtext, inkl. Verweis auf CRISP-DM als Rückgrat von Kap. 3--7.
%
Bei der vorliegenden Arbeit handelt es sich um eine quantitativ-empirische Untersuchung auf Basis
vorhandener Sekundärdaten. Ihr Aufbau orientiert sich am Cross-Industry Standard Process for Data
Mining (CRISP-DM), einem vom Branchensektor und von der eingesetzten Technologie unabhängigen
Vorgehensmodell für Data-Mining-Projekte, das Abschnitt~\ref{sec:crispdm} mit seinen sechs Phasen
vorstellt.\footcite[3, 13]{Chapman2000}

Kapitel 2 steht außerhalb des Prozessmodells und legt die theoretischen und methodischen
Grundlagen; es schließt mit der Forschungslücke und den Erwartungen, die Kapitel 8 den Befunden
gegenüberstellt. Die Kapitel 3 bis 7 durchlaufen die Phasen der Reihe nach: Business
Understanding, Data Understanding, Data Preparation, Modeling und Evaluation. Kapitel 8 diskutiert die Ergebnisse, bewertet die Algorithmen
vergleichend, ordnet den Erklärungsbeitrag der untersuchten Faktoren ein und reflektiert
Limitationen sowie die Übertragbarkeit der Befunde.

% ===================================================
\section{Theoretische Grundlagen und Forschungsstand}
% Umfang: Zielwert 9 Seiten (Entscheidung 27.08.2026) -- siehe Seitenbudget in docs/09_SCHREIBPLAN.md
% ZWISCHEN 2 UND 2.1 STEHT KEIN TEXT (Blindabsatz, Auflage 27.07.). Was das
% Kapitel leistet und warum es ausserhalb des CRISP-DM-Zyklus steht, ist in
% 1.3 bereits gesagt und wird hier NICHT wiederholt -- genau die Redundanz,
% die die Auflage verhindern soll. Der aktuelle Zustand ist richtig.
\subsection{Feuerwehreinsätze als Prognosegegenstand}
% Domäne: Einsatzarten, Einsatzlast, sozioökonomische und bauliche
% Einflussfaktoren (Jennings 1999/2013).

% ---- ÜBERNAHME AUS EXPOSÉ (Kap. 2, Abs. 4, erster Teil) -- 1:1 ----


Eine Feuerwehr rückt nicht nur zu Bränden aus. Die in den Vereinigten Staaten einheitliche Erfassung
nach dem National Fire Incident Reporting System (NFIRS) ordnet jeden Einsatz einer von neun Codeserien zu.
Neben Bränden stehen dort unter anderem rettungsdienstliche Einsätze, technische Hilfeleistungen,
Gefahrstofflagen und Fehlalarme.\footcite[3-21]{USFA2015} Aus dieser Erfassung lassen sich zwei
Größen bilden. Die Einsatzlast eines Gebiets ist die Zahl der Einsätze in einem Zeitraum, unabhängig
von ihrer Art. Das Einsatzartprofil ist ihre Verteilung auf die Serien. Beide sind für die Vorhaltung
von Bedeutung: Die eine bestimmt, wie viele Kräfte bereitstehen müssen, die andere, welche Ausrüstung
sie brauchen.
Eine etablierte Forschungstradition hat gezeigt, dass sich die Häufigkeit von Wohnbränden in
städtischen Quartieren systematisch mit sozioökonomischen und baulichen Strukturmerkmalen
verknüpfen lässt. Jennings zeigt in seinem Review, dass Faktoren wie Armut, Arbeitslosigkeit und
schlechte Wohnverhältnisse konsistent mit erhöhter Brandinzidenz korrelieren, unter anderem weil
in einkommensschwachen Haushalten ältere und wartungsärmere Elektrik sowie Baumaterialien
verbreitet sind. Ergänzend identifiziert er einen ökologischen Erklärungsansatz: Stadtteile mit
hoher Bevölkerungsdichte und baulicher Verdichtung weisen strukturell höhere Brandrisiken auf,
unabhängig vom Verhalten einzelner Bewohner.\footcite[14, 27]{Jennings1999}

% ERLEDIGT 27.08.2026: Einsatzarten und Einsatzlast stehen jetzt im ersten
% Absatz, die Abgrenzung Brandinzidenz gegen Einsatzlast im dritten.
% Jennings2013 ist bewusst NICHT aufgenommen -- bei einer knappen Seite
% traegt Jennings1999 den Abschnitt allein.
% ---- ENDE ÜBERNAHME ----

Diese Forschungstradition untersucht allerdings überwiegend die Brandinzidenz, also einen Ausschnitt
des Einsatzgeschehens. Die Zielgrößen dieser Arbeit sind breiter, denn sie umfassen alle Einsatzarten.
Ob sich die für Brände dokumentierten Zusammenhänge auf die gesamte Einsatzlast übertragen lassen, ist
deshalb offen. Für die Übertragung in eine andere Richtung gibt Jennings bereits eine Warnung: Die
Nachfrage nach Feuerwehrleistungen lässt sich am ehesten innerhalb einer Stadt vorhersagen, während
der Vergleich zwischen Städten wenig belastbare Ergebnisse liefert.\footcite[14]{Jennings1999} Auf den Vergleich von Stadtteilen innerhalb einer Stadt, ist diese Arbeit
angelegt.

\subsection{Datengetriebene Prognose und Algorithmenvergleiche}
% Firebird (Madaio et al. 2016); Benchmark-Literatur (Fernández-Delgado et al.
% 2014, Grinsztajn et al. 2022, Shwartz-Ziv & Armon 2022).

% ---- Herkunft: Exposé Kap. 2, Abs. 1--3 + Abs. 4 zweiter Teil. NICHT MEHR 1:1 ----
% 27.08.2026 an den Volltexten korrigiert: Fernández-Delgado trug den Satz
% "kein Algorithmus ist überall überlegen" NICHT (S. 3175 sagt fast das
% Gegenteil); belegt ist der Abstand 82,0 zu 86,9. Die Aufzählung
% "Stichprobengröße, Signalstärke, Multikollinearität..." stand nicht bei
% Grinsztajn; dort stehen drei Dateneigenschaften, S. 6--9. Begründung im
% Registerkopf der beiden Einträge in literatur.bib.

Die Arbeit ist hinsichtlich der analytischen Verfahren dem Bereich Supervised Machine Learning und
Predictive Analytics zuzuordnen. Sie ordnet sich in die aktuelle Diskussion zum Modellvergleich auf
tabellarischen Daten ein, der für die angewandte Datenwissenschaft von zentraler Bedeutung ist:
Fernández-Delgado et al. vergleichen 179 Klassifikatoren aus 17 Familien auf 121 Datensätzen. Die
Random-Forest-Varianten schneiden im Mittel am besten ab; aufschlussreicher als diese Rangfolge ist
der Abstand, der bestehen bleibt: Die im Mittel beste Konfiguration erreicht 82,0\,\% Genauigkeit,
der jeweils datensatzspezifisch beste Klassifikator 86,9\,\%.\footcite[3133, 3175]{FernandezDelgado2014}
Ein im Mittel überlegenes Verfahren ist damit auf einem einzelnen Datensatz regelmäßig nicht das
beste. Systematische empirische Algorithmenvergleiche sind deshalb kein Methoden-Beiwerk, sondern
Voraussetzung einer begründeten Modellwahl.

Anders als bei Bild- oder Sprachdaten, bei denen Deep-Learning-Verfahren dominieren, gelten bei
strukturierten tabellarischen Daten weiterhin baumbasierte Ensemble-Methoden als Stand der Technik.
Aktuelle Benchmark-Studien zeigen, dass Modelle wie XGBoost und Random Forest auf tabellarischen
Daten mittlerer Größe gegenüber Deep-Learning-Modellen überlegen bleiben, und zwar auch nach
umfangreichem Hyperparameter-Tuning und unabhängig von der Behandlung kategorialer
Variablen.\footcite[1, 6]{Grinsztajn2022} In einer komplementären Untersuchung kommen Shwartz-Ziv und
Armon zu dem Ergebnis, dass XGBoost neuere Deep-Learning-Architekturen auf einer Vielzahl von
Datensätzen übertrifft und dabei einen deutlich geringeren Tuning-Aufwand
erfordert.\footcite[1]{ShwartzZiv2021}

Trotz dieser allgemeinen Befunde sind die Ergebnisse einzelner Vergleichsstudien nicht ohne weiteres übertragbar, denn der Vorsprung der Baumverfahren beruht auf Eigenschaften der Daten und nicht auf einer generellen Überlegenheit. Er entsteht dort, wo die Zielfunktion unregelmäßig verläuft und wo viele Merkmale nichts beitragen: Neuronale Netze bevorzugen glatte Zusammenhänge und werden von uninformativen Merkmalen stärker gestört.\footcite[6-7]{Grinsztajn2022} Entfernt man diese Merkmale, schrumpft der Abstand zu ihnen; fügt man weitere hinzu, wächst er.\footcite[7]{Grinsztajn2022}

Aktuelle Arbeiten zeigen darüber hinaus, dass datengetriebene Modelle einen substantiellen Beitrag
zur operativen Ressourcenplanung von Feuerwehren leisten können: Das Firebird-Projekt der Stadt
Atlanta berechnet mit Random Forest und Support Vector Machine Brandrisikoscores für über 5.000 Gebäude und erreicht
True-Positive-Raten von bis zu 71\,\%.\footcite[185-186]{Madaio2016} Woran sich ein Verfahren messen
lassen muss, um als geeignet zu gelten, ist damit noch nicht gesagt; dafür braucht es eine Referenz.

% AUSBAUEN: Hsu et al. (2025) -- XGBoost auf Feuerwehr-Einsatzdaten. Steht in der
% Literaturübersicht des Exposés, aber ohne Textbeleg. Fehlt bislang in literatur.bib.
% ---- ENDE ÜBERNAHME ----

% #########################################################################
% UMBAU 18.08.2026 -- WARUM DIE REFERENZMODELLE EIN EIGENES UNTERKAPITEL SIND
% #########################################################################
% Vorher stand hier ein einziger Unterabschnitt "Zaehldatenmodelle als
% Baseline" als fuenfter Punkt hinter Ridge, RF und XGBoost. Drei Probleme:
%
%  1  Die MULTINOMIALE LOGISTISCHE REGRESSION kam ueberhaupt nicht vor --
%     obwohl sie die Stufe-2-Baseline des gesamten Klassifikationsstrangs ist
%     und auf dem Hold-out beide Baumverfahren schlaegt (03_STAND.md 5.7).
%  2  Der Abschnitt nannte die NEGATIVE BINOMIAL als Baseline. Seit #45 ist
%     es das unpenalisierte Poisson-GLM. Die Negative Binomial gehoert
%     weiterhin hierher -- aber als ABGEGRENZTE ERWEITERUNG, nicht als
%     gewaehltes Modell.
%  3  Die tragende Behauptung der Arbeit -- ein GLM kann Kruemmung, aber
%     KEINE Wechselwirkungen -- wird in 5.4 und 6.2 verwendet, ohne in den
%     Grundlagen hergeleitet zu sein. Genau das ruegt das Gutachten unter R9.
%
% REIHENFOLGE: Referenzmodelle VOR den untersuchten Verfahren. Erst die
% Messlatte, dann die Kandidaten, die sie schlagen muessen. Ein
% Referenzmodell als fuenfter Punkt hinter den drei ML-Verfahren liest sich
% als Mitlaeufer -- und das Ergebnis der Arbeit ist, dass es gewinnt.
% Ausserdem ist das GLM das elementarere Objekt: Ridge ist ein
% penalisiertes lineares Modell, 2.4 wird dadurch kuerzer, nicht laenger.
%
% ABGRENZUNG NACH 5.4 -- die Regel aus dem Block vor Kapitel 5 gilt hier
% genauso: Kapitel 2 erklaert, WIE ein Verfahren funktioniert, allgemein und
% lehrbuchhaft, OHNE diesen Datensatz. Was NICHT hierher darf:
%     die konkrete Wahl, die Baselinewerte, Auflage D, der zweistufige
%     Aufbau mit den tatsaechlichen Baselines, der Dispersionsindex 62,8
% Das alles steht in 5.4. Wer es hier vorwegnimmt, schreibt 5.4 doppelt.
%
% UMFANG: rund zwei Seiten. Schroeter will den Grundlagenteil High-Level.
% #########################################################################

% #########################################################################
% GESTRICHEN 27.08.2026 -- "Grundlagen des ueberwachten Lernens"
% #########################################################################
% Hier stand ein eigener Unterabschnitt (Regression gegen Klassifikation,
% Verlustfunktion, Bias-Varianz-Zerlegung, Generalisierung). Er faellt weg,
% weil er genau das Basiswissen erklaert, das NICHT gefragt ist: "in
% Grundlagenteil sehr High-Level anfangen" (Vorgaben) und "Ruhig High Level
% einsteigen -- nicht alle Basics erklaeren" (Sprechstunde 24.08.).
% Der Inhalt ist nicht verloren, er steht jetzt dort, wo er gebraucht wird:
%   Regression gegen Klassifikation, Verlustfunktion
%       -> Eroeffnungsabsatz von 2.3.1, zwei Saetze
%   Bias-Varianz und Generalisierung
%       -> 2.5.1, unmittelbar neben Over- und Underfitting
% Damit ruecken alle folgenden Nummern um eins vor: Referenzmodelle 2.3,
% Untersuchte Verfahren 2.4, Modelloptimierung 2.5, Interpretierbarkeit 2.6,
% CRISP-DM 2.7, Forschungsluecke 2.8. Die Registerzeilen in literatur.bib
% sind bereits auf diese Zaehlung geschrieben.
% #########################################################################

\subsection{Referenzmodelle}
\label{sec:referenzmodelle}
% Rund anderthalb Seiten. Hier entstehen drei der Formeln, die das
% Formelverzeichnis: erledigt, sechs Eintraege stehen (27./28.08.2026).
%
% GLIEDERUNGSTIEFE -- bewusst flach (geaendert 18.08.2026):
% Zwei Referenzmodelle, zwei Unterabschnitte. Genauso wie 2.4 je einen
% Unterabschnitt fuer Ridge, Random Forest und XGBoost vorsieht. Eine
% fruehere Fassung hatte VIER Unterabschnitte fuer zwei Modelle -- das
% haette den Referenzmodellen mehr Gliederungsgewicht gegeben als jedem
% einzelnen Vergleichsverfahren und waere unverhaeltnismaessig gewesen.
%
% ZWISCHEN 2.3 UND 2.3.1 STEHT KEIN TEXT -- Blindabsatz, Auflage 27.07.
% Der gemeinsame Rahmen (warum eine Referenz noetig ist, die drei
% Bestandteile eines GLM nach Nelder & Wedderburn 1972, der kanonische Link,
% die Eigenschaft KEINEN freien Parameter zu haben -- darauf beruht spaeter
% Auflage D) eroeffnet deshalb 2.3.1 und wird dort am konkreten Modell
% entwickelt statt vorab abstrakt.  -> FORMEL 1: g(E[y|x]) = x'beta
% Das ist kein Notbehelf: konkret vor abstrakt liest sich besser, und die
% allgemeine Form steht damit unmittelbar neben ihrer ersten Anwendung.
%
% BELEGE -- ALLE NOCH NICHT FREIGEGEBEN.
% Kandidaten, noch zu beschaffen und zu pruefen: Nelder &
% Wedderburn 1972, McCullagh & Nelder 1989, Gourieroux et al. 1984,
% Agresti 2013, James et al. 2021, Santos Silva & Tenreyro 2010.
% Sie waren am 18.08.2026 kurzzeitig in der Datenbank und sind wieder
% entfernt worden -- ohne Freigabe kommt nichts hier hinein.
% AUS DEM EXPOSE VERFUEGBAR: Hastie2009 -- steht in literatur.bib und ist
% sofort einsetzbar. Cameron2013 steht im Expose, aber NICHT in literatur.bib
% (geprueft 24.08.2026) und ist damit noch NICHT einsetzbar; Zugang ueber die
% FOM beschaffen, dann Eintrag mit Registerzeilen anlegen. Zusammen decken
% beide den Abschnitt weitgehend ab.
% Lehrbuecher sind hier RICHTIG -- Schroeters Regel lautet Lehrbuecher fuer
% Grundlagen, Paper fuer den Forschungsstand.
% OFFEN: Der OFFSET fuer die Exposition ist in ISL nicht behandelt. Dafuer
% wird Cameron2013 gebraucht -- Zugang ueber die FOM pruefen.
%
% DIE GRENZE DER MODELLKLASSE steht als SCHLUSSABSATZ am ENDE von 2.3.2.
% Das ist KEIN Blindabsatz -- die Regel betrifft nur die Stelle zwischen
% einer Ueberschrift und ihrer ersten Unterueberschrift. Text am Ende eines
% Unterabschnitts ist unbedenklich. Sie ist kein Modell und bekaeme sonst dieselbe
% Gliederungsebene wie eines -- inhaltlich ist sie aber der Grund, warum
% dieses Unterkapitel ueberhaupt existiert: Ein GLM kann Kruemmung, aber
% keine Wechselwirkungen ohne explizite Spezifikation; genau die finden RF
% und XGBoost konstruktionsbedingt. In 5.4 und 6.2 wird das nur noch
% zitiert. Ohne diesen Absatz steht die tragende Aussage der Arbeit
% unbelegt im empirischen Teil (Gutachten R9).
% Der Absatz ist ein BLINDABSATZ im Sinne der Klaerung vom 11.08.: Er fuehrt
% ein neues Argument ein, er fasst nichts zusammen.
% RUECKFALL, falls er sich beim Lesen verliert: als 2.3.3 mit eigener
% Ueberschrift fuehren.
% Belege: Hastie2009 (Kapitel 4.4 und 9.2), McCullagh1989.

\subsubsection{Poisson-Regression für Zähldaten mit Exposition}
% ZUGEZOGEN 27.08.2026 aus dem gestrichenen Grundlagenabschnitt: Der
% Eroeffnungsabsatz dieses Unterabschnitts klaert in zwei Saetzen Regression
% gegen Klassifikation und den Begriff der Verlustfunktion -- so knapp, dass
% 6.1 die konkrete Wahl noch belegen kann (dort steht Hastie2009 S. 349).
% Log-Link, Erwartungswert gleich Varianz, Offset fuer die Exposition.
%                       -> FORMEL 2: log(mu) = log(E) + x'beta
% UEBERDISPERSION: was sie beschaedigt und was nicht. Der Poisson-Schaetzer
% bleibt bei richtig spezifiziertem bedingtem Mittelwert auch dann
% konsistent, wenn die Verteilungsannahme verletzt ist -- beschaedigt werden
% die Standardfehler (Gourieroux, Monfort & Trognon 1984, S. 701--720).
% DIE NEGATIVE BINOMIAL GEHOERT HIERHER, aber als abgegrenzte Erweiterung:
% Sie loest ein INFERENZproblem und bringt mit dem Dispersionsparameter eine
% zusaetzliche Groesse mit. Ein Modell mit reinen Punktvorhersagen braucht
% beides nicht. Damit ist der Wechsel von #45 theoretisch vorbereitet und
% wirkt in 5.4 nicht wie eine nachtraegliche Rechtfertigung.
% KEINE Zahlen aus dem Datensatz hier -- der Dispersionsindex steht in 4.2.
% Belege: Cameron2013, Gourieroux1984, McCullagh1989.
% ---- AUSGESCHRIEBEN 27.08.2026, vier Absaetze, Formel 1 und 2 ----
Als Referenz eignet sich nur ein Modell, das dieselbe Aufgabe mit geringerem Aufwand löst. Dies leisten hier zwei generalisierte lineare Modelle (GLM). Modelle dieser Klasse haben drei Bestandteile: eine Verteilung aus der Exponentialfamilie, einen
linearen Prädiktor und eine Linkfunktion $g$, die beide verbindet. Die Linkfunktion transformiert den
Erwartungswert der Zielgröße so, dass die transformierte Größe linear von den Merkmalen abhängt:
\begin{equation}
  g\bigl(\mathbb{E}[y \mid \mathbf{x}]\bigr) = \mathbf{x}^{\top}\boldsymbol{\beta}
\end{equation}
\formeleintrag{Allgemeine Form eines verallgemeinerten linearen Modells}
Dabei ist $\mathbb{E}[y \mid \mathbf{x}]$ der bedingte Erwartungswert der Zielgröße, kurz $\mu$,
$\mathbf{x}$ der Vektor der Merkmale einer Beobachtung und $\boldsymbol{\beta}$ der Vektor der zu
schätzenden Koeffizienten; das Produkt $\mathbf{x}^{\top}\boldsymbol{\beta}$ ist der lineare
Prädiktor. Für die lineare, die logistische und die Poisson-Regression lautet die Linkfunktion $g(\mu)=\mu$,
$g(\mu)=\log\bigl(\mu/(1-\mu)\bigr)$ beziehungsweise $g(\mu)=\log(\mu)$.\footcite[173]{James2023} Nelder und Wedderburn haben diese
Klasse eingeführt, um Verfahren, die zuvor getrennt behandelt wurden, unter einem gemeinsamen
Schätzverfahren zusammenzufassen.\footcite[370, 372]{Nelder1972}
Die erste Zielgröße dieser Arbeit ist die Zahl der Einsätze je Stadtteil und Monat: eine nichtnegative ganze Zahl, nach oben unbegrenzt und mit vielen kleinen Werten rechtsschief verteilt. Für Größen dieser Bauart ist die Poisson-Verteilung die Standardannahme, und als Link dient der Logarithmus. Damit sind negative Vorhersagen ausgeschlossen und die
Koeffizienten wirken multiplikativ auf den Erwartungswert. Kennzeichnend für die Poisson-Verteilung
ist, dass Erwartungswert und Varianz zusammenfallen, $\mu = \mathbb{E}(Y) = \operatorname{Var}(Y)$:
Je größer der erwartete Wert, desto größer auch seine Streuung. Ein lineares Modell mit konstanter
Fehlervarianz kann diesen Zusammenhang nicht abbilden.\footcite[168-170]{James2023}
Die Stadtteile San Franciscos unterscheiden sich erheblich in ihrer Wohnbevölkerung; ein Stadtteil mit
40.000 Einwohnern hat mehr Einsätze als einer mit 4.000, ohne dass darin ein Befund läge. Modelliert
wird deshalb nicht die Zahl selbst, sondern die auf die Wohnbevölkerung $B$ bezogene Rate. Setzt man
deren Logarithmus linear in den Merkmalen an, $\log(\mu/B) = \mathbf{x}^{\top}\boldsymbol{\beta}$, so
folgt durch Umstellen:
\begin{equation}
  \log(\mu) = \log(B) + \mathbf{x}^{\top}\boldsymbol{\beta}
\end{equation}
\formeleintrag{Poisson-Regression mit Offset für die Exposition}
Damit tritt der Term $\log(B)$ als Regressor mit fest auf eins gesetztem Koeffizienten in den linearen Prädiktor ein. Ein so vorgegebener Term wird als Offset bezeichnet. Die Bezugsgröße selbst wird als Exposition bezeichnet. Sie muss keine Zeitspanne sein, sondern kann jede Bezugsgröße sein, zu der die erwartete Zahl der Ereignisse proportional ist.\footcite[74-75]{Winkelmann2008} Sie verschiebt das Modell von einer Aussage über Häufigkeiten zu einer Aussage über Raten, ohne einen zusätzlichen Freiheitsgrad zu kosten.
Empirische Zähldaten halten die Gleichheit von Erwartungswert und Varianz selten ein. Liegt die
Varianz deutlich über dem Erwartungswert, spricht man von Überdispersion.\footcite[172]{James2023}
Ist der bedingte Erwartungswert richtig spezifiziert, bleibt der Poisson-Schätzer auch dann
konsistent, wenn die Verteilungsannahme nicht zutrifft; verzerrt werden die geschätzten Standardfehler
und damit die Inferenz.\footcite[89]{Winkelmann2008} Dieses Ergebnis geht auf Gouriéroux, Monfort
und Trognon zurück.\footcite[8]{Gourieroux1981} Die gebräuchliche Erweiterung, die negative Binomialregression, schätzt zusätzlich einen Dispersionsparameter und löst damit die Varianz vom Erwartungswert.\footcite[20-21]{Winkelmann2008} An der Schätzung des bedingten Mittelwerts ändert das nichts. Korrigiert wird die Inferenz. Ein Referenzmodell, das ausschließlich Punktvorhersagen liefert und an ihnen gemessen wird, braucht diese Korrektur nicht.

\subsubsection{Multinomiale logistische Regression}
% Das Gegenstueck zum Poisson-GLM fuer eine nominale Zielgroesse mit K
% ungeordneten Klassen: K-1 Log-Odds gegen eine Referenzkategorie,
% Softmax-Rueckrechnung, unpenalisierte Maximum-Likelihood.
%                       -> FORMEL 3: log(P(y=k)/P(y=K)) = x'beta_k
% Die Klassengewichtung ist am 27.08.2026 hier NICHT aufgenommen worden:
% James2023 behandelt sie nicht, und sie ist eine Entscheidung der Arbeit.
% Sie steht in 5.4 und 6.2, Mechanismus ueber Pedregosa2011.
% Belege, tatsaechlich verwendet: James2023 S. 143 (Referenzkategorie),
%         S. 144 (Rueckrechnung, Chancenverhaeltnis zwischen zwei Klassen).
% <- HIER der Schlussabsatz "Die Grenze der Modellklasse", ohne Ueberschrift.
%    Begruendung und Inhalt siehe Kopf von 2.3.
% ---- AUSGESCHRIEBEN 27.08.2026, vier Absaetze, Formel 3 ----
Die zweite Zielgröße dieser Arbeit ist die je Stadtteil und Monat dominierende Einsatzart, also eine
nominale Größe mit $K$ ungeordneten Klassen. Das Gegenstück zum Poisson-Modell ist hier die multinomiale
logistische Regression. Sie benötigt eine Referenzkategorie, und zwar aus einem zwingenden Grund: Die
Wahrscheinlichkeiten hängen allein von den Differenzen der Koeffizientenvektoren ab, die Niveaus
selbst sind nicht identifiziert. Setzt man die Koeffizienten einer Klasse auf null, ist das Modell
bestimmt. Welche Klasse das ist, bleibt Konvention.\footcite[140]{WinkelmannBoes2006} Gegen die Klasse $K$ als
Referenz schätzt das Modell $K-1$ Gleichungen für die logarithmierten Chancenverhältnisse:
\begin{equation}
  \log \frac{\Pr(y = k \mid \mathbf{x})}{\Pr(y = K \mid \mathbf{x})}
    = \mathbf{x}^{\top}\boldsymbol{\beta}_k ,
    \qquad k = 1, \dots, K-1
\end{equation}
\formeleintrag{Multinomiale logistische Regression gegen eine Referenzkategorie}
Aus diesen $K-1$ Gleichungen lassen sich die Wahrscheinlichkeiten aller $K$ Klassen zurückrechnen. Sie addieren sich zu eins und vorhergesagt wird die Klasse mit der größten. Das logarithmierte Chancenverhältnis zwischen zwei beliebigen Klassen ergibt sich als Differenz der zugehörigen Koeffizientenvektoren.\footcite[145-146]{James2023}
Beide Referenzmodelle werden unpenalisiert nach der Maximum-Likelihood-Methode
angepasst.\footcites[][79]{Winkelmann2008}[][141]{WinkelmannBoes2006} Damit teilen
sie eine Eigenschaft, auf der der Vergleich dieser Arbeit aufbaut: Sie besitzen keinen freien
Hyperparameter und benötigen folglich kein Suchbudget. Ihre Güte hängt nicht davon ab, wie großzügig
ein Suchraum bemessen wurde und sie ist ohne weiteres reproduzierbar. Genau das qualifiziert sie als
Messlatte und macht den Strafterm der Ridge Regression in
Abschnitt~\ref{sec:untersuchte-verfahren} zu dem einen Parameter, den die Referenz nicht kennt.
Damit ist zugleich die Grenze dieser Modellklasse benannt. Der lineare Prädiktor ist additiv in den
Merkmalen, und diese Additivität lässt sich auf zwei Weisen aufweichen. Krümmung ist unproblematisch:
Nimmt man Potenzen eines Merkmals in den Prädiktor auf, bildet das Modell nichtlineare Zusammenhänge
ab, ohne seine Struktur zu verlassen.\footcite[98-99]{James2023} Eine Wechselwirkung zweier Merkmale
dagegen entsteht nicht von selbst. Sie muss als eigener Regressor aufgenommen werden, der als Produkt
der beiden Merkmale gebildet wird.\footcite[95]{James2023} Wer eine Wechselwirkung nicht vorab
vermutet, findet sie in einem verallgemeinerten linearen Modell auch nicht.
Baumbasierte Verfahren unterliegen dieser Beschränkung nicht. Zerlegt man die zu schätzende Funktion in
ihre Bestandteile, einen Summanden je Merkmal, einen je Merkmalspaar, einen je Merkmalstripel und so
fort, so kann ein Baum mit $J$ Endknoten Wechselwirkungen bis zur Ordnung $J-1$
abbilden, ohne dass eine davon spezifiziert werden müsste.\footcite[362]{Hastie2009} Ein verallgemeinertes lineares Modell
ohne Produktterme ist demgegenüber der erste Summand allein. Der mögliche Vorsprung von Random Forest
und XGBoost gegenüber der Referenz liegt damit genau in den übrigen Summanden. Ob dort etwas liegt,
ist keine theoretische, sondern eine empirische Frage.
\subsection{Untersuchte Verfahren}
\label{sec:untersuchte-verfahren}
% ZWISCHEN 2.4 UND 2.4.1 STEHT KEIN TEXT (Blindabsatz). Die Einordnung --
% die Verfahren, die die Referenz aus 2.3 schlagen muessen, in der Ordnung
% linear, Bagging, Boosting -- eroeffnet 2.4.1.
% Der gestrichene Grundlagenabschnitt wird hier NICHT nachgeholt -- was zum
\subsubsection{Ridge Regression}
% Zielfunktion, L2-Regularisierung, Multikollinearität (Hoerl & Kennard 1970).
% Anknuepfen an 2.3: Ridge ist ein PENALISIERTES lineares Modell -- der
% Strafterm ist genau der freie Hyperparameter, den die Referenzmodelle
% nicht haben.                          -> FORMEL 4: Ridge-Zielfunktion
Die drei untersuchten Verfahren geben drei verschiedene Antworten auf dieselbe Frage: Wie lässt sich ein Modell flexibler machen, ohne dass es sich zu eng an die Trainingsdaten anpasst? Ridge bleibt beim linearen Prädiktor und schrumpft die Koeffizienten, der Random Forest mittelt über viele gleichzeitig gewachsene Bäume und senkt damit die Varianz, XGBoost setzt Bäume nacheinander auf die verbleibenden Fehler und senkt damit die Verzerrung.\footcite[588]{Hastie2009} In dieser Reihenfolge lösen sich die
Verfahren immer weiter vom linearen Ansatz.
Die Ridge Regression behält den linearen Prädiktor bei und ändert stattdessen, wie seine
Koeffizienten geschätzt werden. Sind die Merkmale untereinander stark korreliert, reagiert die
Kleinste-Quadrate-Schätzung empfindlich auf kleine Störungen in den Daten, und die geschätzten
Koeffizienten fallen betragsmäßig zu groß aus.\footcite[55-56]{Hoerl1970} Ridge ergänzt die
Residuenquadratsumme deshalb um einen Strafterm, der große Koeffizienten teuer macht:
\begin{equation}
  \hat{\boldsymbol{\beta}}^{\,\mathrm{Ridge}}
  = \arg\min_{\boldsymbol{\beta}} \left\{
      \sum_{i=1}^{n}\bigl(y_i - \mathbf{x}_i^{\top}\boldsymbol{\beta}\bigr)^{2}
      + \lambda \sum_{j=1}^{p} \beta_j^{2}
    \right\}
\end{equation}
\formeleintrag{Zielfunktion der Ridge Regression}
Der Faktor $\lambda$ steuert, wie stark bestraft wird. Bei $\lambda = 0$ entsteht wieder die
gewöhnliche Kleinste-Quadrate-Schätzung; mit wachsendem $\lambda$ werden die Koeffizienten gegen null
gezogen.\footcite[241]{James2023} Das Schrumpfen tauscht eine kleine Verzerrung gegen deutlich
weniger Streuung der Vorhersagen.\footcite[243-244]{James2023}
Das Minimum dieser Zielfunktion lässt sich direkt berechnen. Nullsetzen des Gradienten führt auf
\begin{equation}
  \hat{\boldsymbol{\beta}}^{\,\mathrm{Ridge}}
  = \bigl(\mathbf{X}^{\top}\mathbf{X} + \lambda\mathbf{I}\bigr)^{-1}
    \mathbf{X}^{\top}\mathbf{y}
\end{equation}
\formeleintrag{Geschlossene Lösung der Ridge Regression}
Der Strafterm addiert damit eine positive Konstante auf die Diagonale von
$\mathbf{X}^{\top}\mathbf{X}$. Für jeden Vektor $\mathbf{v} \neq \mathbf{0}$ ist
$\mathbf{v}^{\top}(\mathbf{X}^{\top}\mathbf{X} + \lambda\mathbf{I})\,\mathbf{v}
= \lVert\mathbf{X}\mathbf{v}\rVert^{2} + \lambda\lVert\mathbf{v}\rVert^{2} > 0$. Die
Matrix ist positiv definit und damit invertierbar, auch wenn
$\mathbf{X}^{\top}\mathbf{X}$ selbst nicht vollen Rang hat.\footcite[64]{Hastie2009}
Genau diesen Fall erzeugen stark korrelierte Merkmale: Ridge bleibt dort schätzbar, wo die Kleinste-Quadrate-Rechnung versagt.
\subsubsection{Random Forest}
% Bagging, Feature-Subsampling, OOB-Fehler (Breiman 2001).
Ein einzelner Entscheidungsbaum teilt den Merkmalsraum schrittweise auf und sagt in jedem Endknoten
einen konstanten Wert vorher. Er findet Wechselwirkungen von selbst, ist aber instabil: Schon kleine
Änderungen an den Daten können zu einem völlig anderen Baum führen. Das Bagging mildert diese
Instabilität. Es lässt viele Bäume auf Bootstrap-Stichproben wachsen und mittelt ihre Vorhersagen.
Die Beobachtungen, die ein Baum nicht gezogen hat, dienen ihm als Testdaten und liefern nebenbei eine
Fehlerschätzung, ohne dass dafür Daten zurückgehalten werden müssten.\footcite[343-345]{James2023}
Ein starkes Merkmal steht sonst in fast jedem Baum an der ersten Teilung; die Bäume ähneln einander,
und das Mitteln bringt wenig. Der Random Forest lässt deshalb an jedem Knoten nur eine zufällige
Teilmenge der Merkmale zur Teilung zu. Ihre Größe ist die Stellgröße des Verfahrens: Je kleiner sie
ist, desto unähnlicher werden die Bäume einander und desto weniger streut der Mittelwert über
sie.\footcite[346-347]{James2023}
Breiman zeigt, dass der Verallgemeinerungsfehler eines solchen Waldes mit wachsender Baumzahl gegen
einen Grenzwert strebt. Wie gut dieser Grenzwert ist, hängt von zwei Größen ab: davon, wie treffsicher die einzelnen Bäume
sind und wie stark sie miteinander korrelieren.\footcite[5]{Breiman2001}
\subsubsection{XGBoost}
% Gradient Boosting, Objective mit Regularisierung (Chen & Guestrin 2016).
%                                       -> FORMEL 5: regularisiertes Objective
XGBoost setzt ebenfalls auf viele Bäume, kehrt aber deren Bauweise um. Die Bäume wachsen nicht
unabhängig voneinander, sondern nacheinander: Jeder neue Baum korrigiert das, was die Summe der
bisherigen noch falsch vorhersagt; dieses schrittweise Nachbessern heißt Boosting. Eine solche Folge
neigt zur Überanpassung, und deshalb tritt zur Verlustfunktion ein zweiter Term, der die Komplexität
der einzelnen Bäume begrenzt:
\begin{equation}
  \mathcal{L} \;=\; \sum_{i=1}^{n} l\bigl(y_i, \hat{y}_i\bigr) \;+\; \sum_{k} \Omega(f_k)
\end{equation}
\formeleintrag{Regularisierte Zielfunktion des Gradient Boosting}
Der erste Term misst, wie weit die Vorhersagen danebenliegen. Der zweite Term $\Omega$ bestraft die
Komplexität der Bäume, also die Zahl ihrer Endknoten und die Größe der dort abgelegten
Werte.\footcite[786]{Chen2016} Eine Schrittweite verkleinert zusätzlich den Beitrag jedes neuen
Baumes und lässt den folgenden Raum zum Nachbessern.\footcite[787]{Chen2016} Die Stellgrößen sind
dabei nicht unabhängig voneinander wählbar: Eine höhere Zahl von Bäumen verlangt eine kleinere
Schrittweite und umgekehrt, und die Suche nach guten Werten wird entsprechend
aufwendiger.\footcite[17]{Probst2019}

\subsection{Modelloptimierung und Modellbewertung}
\label{sec:modellbewertung}
% ZWISCHEN 2.5 UND 2.5.1 STEHT KEIN TEXT (Blindabsatz).

\subsubsection{Hyperparameter-Tuning, Overfitting und Underfitting}
% Random Search: Bergstra2012, S. 281 (Befund und Mechanismus), S. 282
% (Gittersuche) -- am Volltext geprueft 27.08. RF-Robustheit gegenueber
% Defaults: Probst2019, Locator Abschn. 3.1 -- ACHTUNG, siehe Warnung unten.
% ZUGEZOGEN 27.08.2026 aus dem gestrichenen Grundlagenabschnitt:
% Bias-Varianz-Zerlegung und Generalisierung gehoeren hierher, unmittelbar
% neben Over- und Underfitting. Beleg liegt vor und ist geprueft --
% Hastie2009 S. 38 (Abschn. 2.9, geprueft 21.08.) oder James2023 S. 30, 33, 37.

% ---- ÜBERNAHME AUS EXPOSÉ (Kap. 3, Abs. 3, letzter Satz) -- Satzfragment, ausbauen ----

% ACHTUNG 20.08.2026 -- DIESER SATZ AUS DEM EXPOSE IST FALSCH BELEGT.
% Probst et al. (2019) schreiben das Gegenteil: RF laufe mit den Defaults meist gut und sei
% "far less tunable than other algorithms such as support vector machines" (Abschn. 3.1).
% Die Seite 166 existiert nicht, der Artikel traegt nur die Artikelnummer e1301.
% In 6.3 ist der Satz bereits umformuliert. Hier NICHT in der alten Fassung uebernehmen.
% Die besondere Sensitivität von Random Forest gegenüber Hyperparameter-Einstellungen ist gut
% dokumentiert.\footcite[166]{Probst2019}
% -> Als Beleg brauchbar, trägt aber keinen eigenen Abschnitt. Hier gehört die
%    Funktionsweise von Random Search hin (Bergstra & Bengio 2012) sowie
%    Overfitting/Underfitting -- beides liefert das Exposé nicht.
% ---- ENDE ÜBERNAHME ----

Jedes der drei Verfahren verfügt über Einstellungen, die nicht aus den Daten geschätzt, sondern vorab gewählt werden. Sie steuern die Flexibilität des Modells, die eine Abwägung darstellt.
Ein zu starres Modell bildet die Struktur der Daten nicht ab und liegt schon bei den Trainingsdaten daneben. Man spricht in diesem Fall von Underfitting. Ein zu flexibles Modell folgt dem Rauschen der Trainingsdaten und verliert dadurch an Genauigkeit bei neuen Beobachtungen , ein Phänomen, das als Overfitting bezeichnet wird.\footcite[124]{Murphy2022} Formal zerlegt sich
der erwartete Fehler in einen Verzerrungs- und einen Varianzanteil. Beide verhalten sich gegenläufig
zur Flexibilität, sodass das Minimum des Gesamtfehlers zwischen den Extremen
liegt.\footcite[38, 223]{Hastie2009}

Die passenden Einstellungen zu suchen heißt, den Modellfehler über einem mehrdimensionalen Raum zu
minimieren. Die Zielgröße dieser Suche ist der Fehler auf einer Validierungsmenge, und sie liegt nicht in
geschlossener Form vor.\footcite[319]{Murphy2022} Eine Gittersuche prüft jede Kombination vorgegebener
Werte, und die Zahl der Versuche wächst exponentiell mit der Zahl der Parameter. Eine Zufallssuche
zieht stattdessen eine feste Zahl von Punkten und findet bei gleichem Rechenbudget die besseren
Modelle, weil in den meisten Datensätzen nur wenige Parameter wirklich wichtig sind aber je nach
Datensatz verschiedene.\footcite[281-282]{Bergstra2012} Eine begutachtete Übersicht bestätigt diesen
Befund.\footcite[7-8]{Bischl2023} Unabhängig vom Suchverfahren gilt dabei, dass der beim Tuning
erreichte Wert nicht zugleich die Prognosegüte belegen kann: Wer auf denselben Daten optimiert und
bewertet, erhält eine optimistisch verzerrte Schätzung. Die Bewertung verlangt deshalb eine zusätzliche,
äußere Aufteilung.\footcite[18-19]{Bischl2023}

\subsubsection{Gütemaße für Regression und Klassifikation}
% RMSE, MAE, R^2; Macro-F1, AUROC. Hier gehoeren die Formeln hin.
% Belege, alle vorhanden und am Volltext geprueft:
%   James2023 S. 150-153  Konfusionsmatrix, Sensitivitaet, Spezifitaet,
%                         ROC, AUC -- aber KEIN F-Mass
%   Manning2008           Praezision, Trefferquote, F1, Makro-Mittelung
%   Hand2001              Multiclass-AUROC, passend zu macro_auroc im Code
%   Opitz2021 S. 1, 6     welches der ZWEI Macro-F1 gemeint ist. Die Arbeit
%                         rechnet das arithmetische Mittel der klassenweisen
%                         F1-Werte (m03_struktur.py:292, average="macro"),
%                         also Gl. 2 -- die von Opitz empfohlene Fassung.
%                         S. 6 verlangt ausdruecklich, die Formel zu nennen;
%                         ein Satz erledigt das.

Die beiden Zielgrößen verlangen verschiedene Gütemaße. Die Einsatzhäufigkeit wird als Regressionsaufgabe bewertet. Leitmaß ist die Wurzel des mittleren quadratischen Fehlers (RMSE), weil sie in der Einheit der Zielgröße bleibt. Daneben stehen der mittlere absolute Fehler (MAE), der große Abweichungen schwächer gewichtet und das Bestimmtheitsmaß als Anteil der erklärten Streuung.\footcite[29, 79]{James2023}

Die überwiegende Einsatzart wird als Klassifikationsaufgabe bewertet. Dort ist der Anteil richtig zugeordneter Fälle, die Accuracy, allein irreführend, sobald die Klassen ungleich besetzt sind: Ein Modell, das stets die häufigste Klasse vorhersagt, erreicht bereits einen hohen Wert. Ausgangspunkt ist deshalb die Konfusionsmatrix. Aus ihr ergibt sich über verschiedene Entscheidungsschwellen hinweg die ROC-Kurve mit der Fläche unter ihr (AUROC).\footcite[152-155]{James2023} Über die Klassen gemittelt ist dieses Maß auch bei mehr als zwei Klassen anwendbar.\footcite[176]{Hand2001} Für mehrere Klassen wird je Klasse aus Präzision und
Trefferquote das F1-Maß als deren harmonisches Mittel gebildet. Dieses fällt
deutlich ab, sobald einer der beiden Werte klein ist und verlangt damit hohe Präzision und hohe
Trefferquote zugleich.\footcite[175]{Murphy2022} Das Macro-F1 ist das arithmetische
Mittel dieser klassenweisen Werte:\footcite[1]{Opitz2021}

\begin{equation}
  F1_k = \frac{2\,P_k\,R_k}{P_k + R_k} ,
  \qquad
  \text{Macro-}F1 = \frac{1}{K}\sum_{k=1}^{K} F1_k
\end{equation}
\formeleintrag{Klassenweises F1-Maß und makro-gemitteltes F1}

Diese Mittelung gewichtet jede Klasse gleich, unabhängig davon, wie häufig sie auftritt. Sie ist das Leitmaß des Strukturstrangs aber nicht die einzige gebräuchliche Definition: Die verbreiteten Formeln können auseinanderfallen und zu unterschiedlichen Rangfolgen führen, weshalb hier ausdrücklich die oben angegebene verwendet wird.\footcite[6]{Opitz2021}

Ein Gütewert aus der Kreuzvalidierung ist selbst eine Zufallsgröße: Dieselben Daten liefern bei
anderer Zuteilung der Blöcke einen anderen Wert. Der Abstand zweier Verfahren ist deshalb ohne Prüfung
nicht von dieser Zuteilung zu trennen. Berichtet werden neben dem Punktwert die Streuung und eine
Testentscheidung.\footcite[434]{Fahrmeir2023}

Geprüft wird gepaart und verteilungsfrei. Gepaart, weil jedes Verfahren dieselben Blöcke sieht und
sich die Streuung zwischen den Blöcken dadurch heraushebt. Verteilungsfrei, weil über die Verteilung
der Differenzen nichts bekannt ist und sie sich an wenigen Werten auch nicht prüfen ließe. Der
Vorzeichen-Rang-Test nach Wilcoxon rangiert dazu die absoluten Differenzen und summiert die Ränge der
positiven. Gleiche Beträge erhalten Durchschnittsränge.\footcite[447-450]{Fahrmeir2023} Für den
Vergleich von Lernverfahren gilt er gegenüber dem $t$-Test als das sicherere Verfahren, weil er dessen
Verteilungsannahme nicht benötigt.\footcite[8]{Demsar2006}

Mit der Zahl der Vergleiche wächst die Wahrscheinlichkeit, mindestens einmal fälschlich einen
Unterschied zu finden: Bei zehn Tests zum Niveau 0,05 beträgt sie 0,401.\footcite[435]{Fahrmeir2023}
Die Bonferroni-Korrektur senkt dafür jedes Einzelniveau auf $\alpha/k$. Das Verfahren von Holm
erreicht dieselbe Absicherung schrittweise, indem es die $p$-Werte aufsteigend sortiert und den
$i$-ten gegen $\alpha/(k-i+1)$ prüft; es ist dabei trennschärfer.\footcite[12-13]{Demsar2006}

Das Niveau $\alpha = 0{,}05$ ist Konvention. Ein inhaltlicher Grund für gerade diesen Wert besteht
nicht. Ein nicht signifikantes Ergebnis belegt deshalb keine Gleichheit, sondern das Fehlen eines
Nachweises.\footcite[485]{Altman1995}

\subsubsection{Validierung bei abhängigen Daten}
% UMBENANNT 27.08.2026, vorher "Zeitbewusste Validierung". Die alte
% Ueberschrift stammt aus der urspruenglichen Fassung von UF2 und passt
% nicht mehr: seit #29 wird STADTTEILWEISE kreuzvalidiert, nicht
% zeitreihengerecht. Inhalt in dieser Reihenfolge:
%   1  Warum die unabhaengige Zufallsaufteilung hier nicht traegt -- die
%      Beobachtungen sind in Stadtteilen gruppiert und in der Zeit geordnet
%   2  Block- bzw. Leave-Group-Out-Kreuzvalidierung als Antwort darauf
%      -> Roberts2017, S. 913 (Empfehlung), S. 925 ("Step 3"), geprueft 24.08.
%   3  Leakage und der Verfuegbarkeitsgrundsatz
%      -> Kaufman2011, S. 559 ("no-time-machine requirement"), S. 560
%         ("learn-predict separation"), geprueft 24.08.
%   4  Die Zeitkomponente als ZWEITER Punkt, nicht als erster: Lags und
%      rollierende Fenster enden im Vormonat. Bergmeir2012 nur, falls die
%      Rolling-/Expanding-Window-Verfahren wirklich dargestellt werden --
%      der Volltext liegt hinter einer Elsevier-Schranke, der Abschnitt
%      kommt ohne ihn aus.
% KEINE Beschreibung des eigenen Vorgehens hier -- die steht in 5.4 und 6.1.
% Das Exposé-Fragment ("Die Evaluation erfolgt zeit-bewusst mittels
% Time-Series-Cross-Validation", Kap. 3 Abs. 4) ist damit ueberholt und
% NICHT zu uebernehmen: es beschreibt ein Vorgehen, das so nicht mehr gilt.

Die Güte eines Modells wird an Daten gemessen, die es nicht gesehen hat. Die übliche $k$-fache
Kreuzvalidierung teilt die Beobachtungen dafür zufällig in $k$ Blöcke, hält jeden einmal zurück und
mittelt die Ergebnisse.\footcite[201-208]{James2023} Diese zufällige Aufteilung setzt voraus, dass die
Beobachtungen voneinander unabhängig sind. Bei räumlich oder zeitlich strukturierten Daten trifft das
nicht zu. Liegen Beobachtungen desselben Gebiets gleichzeitig in der Trainings- und in der Testmenge,
misst die Kreuzvalidierung nicht mehr, wie gut ein Modell auf unbekanntes Gebiet verallgemeinert,
sondern nur noch, wie gut es bekanntes Gebiet wiedererkennt.

Die Antwort darauf ist eine blockweise Aufteilung. Die Gruppen, entlang derer die Abhängigkeit
verläuft, werden dabei vollständig einer Seite zugeschlagen, sodass keine Gruppe geteilt wird. Roberts
et al. empfehlen, die Blöcke so zu schneiden, dass sie derjenigen Struktur entsprechen, auf die das
Modell später angewendet werden soll.\footcite[913, 925]{Roberts2017}

Eine zweite, davon unabhängige Fehlerquelle ist die Datenleckage. Von ihr spricht man, wenn
Information in die Merkmale gelangt, die zum Zeitpunkt der Vorhersage noch nicht verfügbar gewesen
wäre. Kaufman et al. fassen die Anforderung so, dass ein zulässiges Modell nur auf Information beruhen
darf, die zeitlich vor der Zielgröße liegt. Sie halten zugleich fest, dass Leckage bei der Aufbereitung
der Daten zu vermeiden ist und nicht erst bei der Modellierung.\footcite[559-560]{Kaufman2011}

\subsection{Interpretierbarkeit von Modellen}
% SHAP (Lundberg & Lee 2017), Zielkonflikt Güte vs. Interpretierbarkeit
% (Rudin 2019, Molnar 2022), erklärend vs. prognostisch (Shmueli 2010).

% ---- Herkunft: Exposé Kap. 2, Abs. 3 zweiter Teil + Kap. 3, Abs. 4. NICHT MEHR 1:1 ----
% 27.08.2026 korrigiert: Der Satzteil, Post-hoc-Verfahren wie SHAP könnten
% den Zielkonflikt "teilweise entschärfen", wird von Rudin nicht gestützt,
% sondern bestritten (S. 1, S. 4). Ersetzt durch das, was Rudin auf S. 2--3
% tatsächlich sagt; SHAP hat jetzt Lundberg2017 S. 4 als eigenen Beleg.

Gütemaße und Validierung sagen, wie gut ein Modell trifft. Sie sagen nicht, welche Merkmale das
Ergebnis tragen. Diese zweite Frage berührt einen Zielkonflikt zwischen Prognosegüte und
Interpretierbarkeit, den Rudin auf eine Setzung zurückführt, die nicht allgemein gilt: Der Glaube an
einen solchen Zielkonflikt halte viele Forschende davon ab, ein interpretierbares Modell überhaupt zu
versuchen. Ist der Leistungsunterschied zwischen den Verfahren klein, könne er sich sogar umkehren,
weil ein interpretierbares Modell die Daten- und Modellarbeit im nächsten Durchgang erleichtert und
so die Genauigkeit insgesamt erhöht.\footcite[2-3]{Rudin2019}Weil sich eine nachträgliche Erklärung auf andere Merkmale stützen kann als das erklärte Modell und dessen Rechenweg dann nicht getreu abbildet, fordert Rudin für Entscheidungen mit hohem Risiko inhärent interpretierbare Modelle.\footcite[1, 4]{Rudin2019}

Die vorliegende Arbeit trägt diesem Einwand Rechnung, statt ihn zu übergehen. Ein inhärent
interpretierbares Modell steht als Referenz im Vergleich selbst und nicht nur als nachträgliche
Erklärung der übrigen. Die Methode SHapley Additive exPlanations (SHAP), die Merkmalsbeiträge als
Shapley-Werte einer bedingten Erwartungsfunktion bestimmt,\footcite[4]{Lundberg2017} wird auf die
Black-Box-Verfahren angewandt, also genau dort, wo Rudins Vorbehalt greift, und mit ihm im Blick.

Davon zu trennen ist eine zweite Unterscheidung, die zwischen erklärender
und prognostischer Modellierung, auf die Shmueli hinweist. Erklärende Modelle prüfen theoretisch begründete Zusammenhänge,
prognostische Modelle sagen neue Beobachtungen vorher; in der Praxis fällt beides häufig zusammen,
obwohl es verschiedene Ziele sind.\footcite[291]{Shmueli2010} Daraus folgt eine Möglichkeit, die für
diese Arbeit unmittelbar von Bedeutung ist: Ein im Sinne der Theorie falsches Modell kann unter
Umständen besser vorhersagen als das richtige.\footcite[293]{Shmueli2010}

% AUSBAUEN: Funktionsweise der Shapley-Werte (Lundberg & Lee 2017, S. 17), Molnar (2022)
% als Überblick. Das Exposé liefert nur die Einordnung, nicht die Methodik.
% ---- ENDE ÜBERNAHME ----

\subsection{CRISP-DM als Vorgehensmodell}
\label{sec:crispdm}
% Sechs Phasen (Wirth & Hipp 2000); Deployment als dokumentierte Lücke
% (Schröer et al. 2021) und Begründung der Abgrenzung.

% ---- Herkunft: Exposé Kap. 3, Abs. 1, ohne den Satz zur Klassenverteilung (-> 4.2). NICHT MEHR 1:1 ----
% 27.08.2026 korrigiert: Die Gewichtung von Modeling und Evaluation ist eine
% Entscheidung dieser Arbeit und stand fälschlich unter einer Schröer-Fußnote.
% Die Fußnote trägt jetzt die Deployment-Abgrenzung, die sie belegen kann
% (S. 531--532) -- damit ist zugleich die ERGÄNZEN-Notiz unten erledigt.

Methodisch folgt die Arbeit dem CRISP-DM-Schema, das nach dem Anspruch seiner Urheber industrie-,
werkzeug- und anwendungsneutral sein sollte.\footcite[3]{Chapman2000} Es gliedert ein Projekt in
sechs iterativ durchlaufene Phasen: Business Understanding, Data Understanding, Data Preparation,
Modeling, Evaluation und Deployment.\footcite[13]{Chapman2000}

Da der Schwerpunkt dieser Arbeit auf dem Algorithmenvergleich liegt, erhalten die Phasen Data Preparation, Modeling und
Evaluation überproportionales Gewicht. Die Phase Deployment ist nicht Gegenstand der Arbeit, weil kein
produktiver Einsatz vorgesehen ist und die Ergebnisse als Entscheidungsgrundlage übergeben werden,
nicht als lauffähiges System. Diese Grenze ist keine Eigenheit dieser Arbeit: Die Deployment-Phase fehlt in einigen untersuchten Studien, obwohl sie Teil
des Prozessmodells ist, dass werter Schröer als Lücke des Feldes.\footcite[531-532]{Schroeer2021}

% ERLEDIGT 27.08.2026 -- Begründung, warum die Deployment-Phase nicht Gegenstand ist. Schröer et al.
% identifizieren genau diese Phase als dokumentierte Lücke -- lässt sich elegant zitieren.
% ---- ENDE ÜBERNAHME ----

\subsection{Forschungsstand und Forschungslücke}
\label{sec:forschungsluecke}
% Scharnier zwischen Theorie und Empirie. Ableitung der Erwartungen,
% die in 8.1 und 8.2 gegen die Befunde gestellt werden.

Der Forschungsstand lässt sich in fünf Strängen bündeln, die dem Aufbau dieses Kapitels folgen. Tabelle~\ref{tab:forschungsstand} stellt die einschlägigen Arbeiten nebeneinander und gibt je Quelle die tragende Aussage wieder. Aufgenommen sind nur empirische und methodische Originalarbeiten, ausgewählt nach thematischer Nähe, Aktualität der Untersuchungsanlage und Prüfbarkeit am Volltext; die Lehrbuchliteratur der Abschnitte 2.3 bis 2.5 bleibt unberücksichtigt.

% ---------------------------------------------------------------
% 2.8  Quertabelle Forschungsstand
% braucht: longtable, booktabs, array  (array kommt mit tabularx)
% ---------------------------------------------------------------
\begingroup
\singlespacing
\footnotesize
\setlength{\tabcolsep}{4pt}
\begin{longtable}{@{}
  >{\raggedright\arraybackslash}p{0.22\textwidth}
  >{\raggedright\arraybackslash}p{0.06\textwidth}
  >{\raggedright\arraybackslash}p{0.66\textwidth}@{}}
\caption{Quertabelle zum Forschungsstand}\label{tab:forschungsstand}\\
\toprule
\textbf{Autor} & \textbf{Jahr} & \textbf{Verdichtung}\\
\midrule
\endfirsthead
\toprule
\textbf{Autor} & \textbf{Jahr} & \textbf{Verdichtung}\\
\midrule
\endhead
\midrule
\multicolumn{3}{r@{}}{\itshape Fortsetzung auf der nächsten Seite}\\
\endfoot
\bottomrule
\endlastfoot

\multicolumn{3}{@{}l}{\textbf{Einsatzprognose und urbanes Risiko}}\\[2pt]

Brantingham \& Brantingham & 1997
& Zählung, Rate und Standortquotient desselben Geschehens ergeben verschiedene
Bilder; erst der relative Deliktmix zeigt, worin sich ein Gebiet unterscheidet.\\

Jennings & 1999
& Sozioökonomische und bauliche Merkmale sind die wichtigsten Bestimmungsgrößen
des Brandgeschehens; innerhalb einer Stadt sind die Zusammenhänge konsistenter
als zwischen Städten.\\

Lum \& Isaac & 2016
& Modelle auf polizeilichen Erfassungsdaten reproduzieren deren Verzerrung und
verstärken sie über eine Rückkopplung.\\

Madaio et~al. & 2016
& Brandrisiko ist aus sozioökonomischen, kriminalitätsbezogenen und baulichen
Merkmalen prognostizierbar -- auf Gebäudeebene, nur für Gewerbeobjekte.\\

Jaldell & 2017
& Kürzere Hilfsfristen senken die Zahl der Brandtoten, mit dem größten Effekt
im unteren Bereich der Fristen.\\

Dey et~al. & 2023
& Standortplanung stützt sich auf eine erwartete statt auf die beobachtete
Einsatzlast; Gegenstand sind Standorte, nicht die monatliche Last.\\
\addlinespace

\multicolumn{3}{@{}l}{\textbf{Verfahrensvergleiche auf tabellarischen Daten}}\\[2pt]

Fernández-Delgado et~al. & 2014
& Über 121 Datensätze führt Random Forest im Mittel, ohne auf dem einzelnen
Datensatz durchgängig das beste Verfahren zu sein.\\

Shwartz-Ziv \& Armon & 2021
& XGBoost übertrifft neuere Deep-Learning-Architekturen bei deutlich geringerem
Tuningaufwand.\\

Grinsztajn et~al. & 2022
& Baumverfahren bleiben über jedes Suchbudget überlegen; Ursache sind
Glättungsneigung, uninformative Merkmale und Rotationsinvarianz der Netze.\\
\addlinespace

\multicolumn{3}{@{}l}{\textbf{Tuning, Gütemaße und Vergleichspraxis}}\\[2pt]

Bergstra \& Bengio & 2012
& Auf jedem Datensatz wirken nur wenige Hyperparameter, aber jeweils andere;
daran scheitert die Gittersuche, die Zufallssuche nicht.\\

Probst et~al. & 2019
& Wie viel Tuning bringt, hängt am Verfahren: beim Random Forest am wenigsten,
bei der regularisierten Regression am meisten.\\

Opitz \& Burst & 2021
& Für das Macro-F1 sind zwei Formeln im Umlauf, die zu deutlich
unterschiedlichen Werten und Rangfolgen führen können.\\

Schröer et~al. & 2021
& CRISP-DM ist De-facto-Standard; die Deployment-Phase fehlt in der Mehrzahl
der ausgewerteten Studien.\\

Bischl et~al. & 2023
& Optimierung und Bewertung müssen getrennt bleiben; erst geschachteltes
Resampling schätzt die künftige Güte unverzerrt.\\
\addlinespace

\multicolumn{3}{@{}l}{\textbf{Validierung bei abhängigen Daten}}\\[2pt]

Hurlbert & 1984
& Wer abhängige Beobachtungen wie unabhängige behandelt, testet mit einem
Fehlerterm, der nicht zur Hypothese passt.\\

Nadeau \& Bengio & 2003
& Wer die Streuung aus der Wahl der Trainingsmenge unterschlägt, unterschätzt
die Varianz grob und weist Verfahren fälschlich als besser aus.\\

Kaufman et~al. & 2011
& Leakage ist ein Entwurfsproblem, kein Einzelfehler; vermieden wird es durch
Trennung von Lernen und Vorhersage entlang der Zeitachse.\\

Roberts et~al. & 2017
& Zufällige Kreuzvalidierung unterschätzt den Fehler strukturierter Daten;
Blocken ist angemessener, erzwingt aber selbst Extrapolation.\\
\addlinespace

\multicolumn{3}{@{}l}{\textbf{Interpretierbarkeit und Zielkonflikt}}\\[2pt]

Strobl et~al. & 2008
& Die Permutationswichtigkeit bevorzugt korrelierte Prädiktoren; eine bedingte
Permutation korrigiert das.\\

Shmueli & 2010
& Erklärung und Prognose sind verschiedene Ziele; ein falsches Modell kann
besser prognostizieren als das richtige.\\

Rudin & 2019
& Nachträgliche Erklärungen sind der Rechnung der Black Box nicht treu; der
Zielkonflikt kann sich umkehren.\\

Hooker et~al. & 2021
& Permutieren erzwingt Extrapolation in dünn besetzte Bereiche und überzeichnet
korrelierte Merkmale.\\

Aas et~al. & 2021
& Kernel-SHAP unterstellt unabhängige Merkmale; bei Korrelation werden die
Erklärungen irreführend.\\

\end{longtable}

\noindent{\footnotesize Quelle: Eigene Darstellung\par}
\endgroup
\nocite{Brantingham1997,Jennings1999,LumIsaac2016,Madaio2016,Jaldell2017,Dey2023,
FernandezDelgado2014,ShwartzZiv2021,Grinsztajn2022,Bergstra2012,Probst2019,
Opitz2021,Schroeer2021,Bischl2023,Hurlbert1984,Nadeau2003,Kaufman2011,
Roberts2017,Strobl2008,Shmueli2010,Rudin2019,Hooker2021,Aas2021}

Dass Baumverfahren auf tabellarischen Daten die aussichtsreichere Wahl sind, ist breit belegt, ohne dass daraus für einen einzelnen Datensatz ein bestes Verfahren folgt.\footcite[3175]{FernandezDelgado2014}\footcite[1]{Grinsztajn2022} Ebenso belegt ist, dass sozioökonomische und bauliche Merkmale das Brandgeschehen bestimmen, am deutlichsten innerhalb einer Stadt.\footcite[27]{Jennings1999} Wie belastbar ein Vergleich ist, entscheidet aber der Versuchsplan: gleiches Suchbudget, Trennung von Optimierung und Bewertung\footcite[Abschn. 4.4]{Bischl2023} und eine Aufteilung entlang der Abhängigkeitsstruktur.\footcite[913]{Roberts2017}

Die Vergleichsstudien arbeiten auf unabhängigen Datensätzen und prüfen nicht, ob ihre Rangfolge auch für gruppierte Beobachtungen und unbekannte Gruppen gilt. Die Feuerwehrarbeiten prognostizieren Risiko im Querschnitt, für einzelne Gebäude\footcite[185]{Madaio2016} oder Standorte\footcite[34]{Dey2023}, nicht als laufende Einsatzlast eines Gebiets. Keine misst sich an einem begründeten Zähldatenmodell. Die Validierungsliteratur entwickelt das Blocken an ökologischen und zeitreihenbezogenen Beispielen, nicht an einem Panel aus Gebietseinheiten. Der Rechenaufwand kommt als eigenes Kriterium nirgends vor.

Das Problem dieser Arbeit ist damit benannt: für einen Stadtteil ohne eigene Einsatzhistorie die monatliche Einsatzlast aus allgemein verfügbaren Merkmalen vorherzusagen — und zu entscheiden, welches Verfahren das leistet, gemessen an Güte, Interpretierbarkeit und Aufwand zugleich. Entscheidungseinheit ist deshalb der Stadtteil im Monat, Zielgröße eine Zählgröße mit Exposition; Übertragbarkeit wird nur für unbekannte Stadtteile derselben Stadt beansprucht, zwischen Städten trägt sie nicht.\footcite[14]{Jennings1999} Das Kriminalitätsmerkmal geht als relativer Anteil ein\footcite[264]{Brantingham1997}, mit der Verzerrung polizeilicher Daten als Einschränkung.\footcite[15]{LumIsaac2016} Referenz ist ein begründetes Zähldatenmodell statt eines naiven; der Aufwand wird als drittes Kriterium geführt.
% NICHT übernehmen: Absatz 1 Satz 1 („Ziel der Arbeit ist es aufzuzeigen...") und der
% Schlussabsatz („Als Ergebnis soll eine ... Blaupause ... vorliegen") -- reine Exposé-Rhetorik.
% ---- ENDE ÜBERNAHME ----

% ===================================================
\section{Business Understanding}
% Umfang: ca. 3 Seiten -- CRISP-DM Phase 1
\subsection{Anwendungsfall und fachliche Zielsetzung}
\label{sec:anwendungsfall}
Die Leistung einer Feuerwehr bemisst sich in Minuten. Für schwedische Wohnungsbrände weist Jaldell nach, dass die Zahl der Todesopfer mit der Hilfsfrist steigt, und zwar nichtlinear: Der Zuwachs ist bei kurzen Anfahrten am größten und verliert sich mit wachsender Entfernung. Eine um eine Minute verkürzte mittlere Hilfsfrist entspräche knapp drei Prozent weniger Brandtoten in Wohngebäuden.\footcite[706]{Jaldell2017} Über die Hilfsfrist entscheidet jedoch nicht der einzelne Einsatz, sondern die Entfernung zwischen Wache und Einsatzstelle und die Frage, ob im Alarmfall überhaupt eine Einheit verfügbar ist. Beides ist Ergebnis von Entscheidungen über Wachstandorte, Fahrzeuge und Personal, die lange vor dem Alarm fallen. Sie stützen sich deshalb nicht auf die beobachtete, sondern auf eine erwartete Einsatzlast. Für die Standortplanung ist eine merkmalbasierte Prognose bereits erprobt\footcite[34]{Dey2023}.

Solche Vorhalteentscheidungen setzen eine Erwartung darüber voraus, wo künftig wie viel zu leisten sein wird, und sie werden nicht für den einzelnen Alarm getroffen, sondern für Gebiete und Zeiträume. Daraus ergibt sich die Entscheidungseinheit dieser Arbeit. Räumlich ist es der Stadtteil in der Abgrenzung der 41 Analysis Neighborhoods, die San Francisco aus ganzen Census Tracts gebildet hat, um seine Daten einheitlich zu berichten.\footcite{DataSFStadtteile} Feinere Zuschnitte wären verfügbar, entsprechen aber keiner Berichtseinheit der Stadt. Zeitlich ist es der Monat, der die Zufälligkeit des Tagesgeschehens glättet, ohne die jahreszeitliche Schwankung zu verwischen und ohne den Vorlauf zu unterschreiten, den eine Vorhalteentscheidung braucht.

Aus Strukturmerkmalen lassen sich Zeitpunkt und Ort eines konkreten Brandes nicht ableiten. Untersucht wird das strukturell erwartbare Niveau der Einsatzlast eines Stadtteils. Der Nutzen einer merkmalbasierten Prognose liegt somit weniger in San Francisco selbst, wo die Einsatzhistorie vorliegt, als in der Frage, ob sich das Niveau eines Stadtteils allein aus seinem Profil bestimmen lässt. Also für Gebiete, für die eine eigene Einsatzstatistik fehlt oder zu kurz ist.

Dass die Untersuchung an San Francisco ansetzt, folgt aus der Datenlage: Die Stadt veröffentlicht die Meldungen ihrer Feuerwehr ebenso offen wie die Merkmalsdaten, aus denen sich ein Stadtteilprofil bilden lässt, und bezieht beides auf dieselbe amtliche Gliederung, nach der sie auch plant. Analyseeinheit und Entscheidungseinheit fallen damit zusammen.


\subsection{Zielgrößen, Erfolgskriterien und Abgrenzung}
\label{sec:zielgroessen}
Die Einsatzlast wird in zwei Hinsichten erfasst: nach ihrer Menge und nach ihrer Zusammensetzung. Den Mengenstrang bildet die Einsatzhäufigkeit, die Zahl der Einsätze je Stadtteil und Monat, also eine Zählgröße, die als Regressionsaufgabe behandelt wird. Sie wird zusätzlich auf die Wohnbevölkerung normiert, weil sich erst dadurch trennen lässt, ob ein Stadtteil viele Einsätze aufweist, weil er groß ist, oder weil seine Struktur sie hervorbringt. Den Strukturstrang bildet die Einsatzart, die im jeweiligen Stadtteilmonat überwiegende Art des Einsatzgeschehens, also eine kategoriale Größe, die als Klassifikationsaufgabe behandelt wird. Damit stehen drei Zielgrößen nebeneinander, zwei des Mengen- und eine des Strukturstrangs. Berichtet wird der Mengenstrang auf der Anzahl, die Rate dient der Robustheitsprüfung.

Für diese Arbeit werden die Einsatzarten des NFIRS zu vier Gruppen zusammengefasst, die Einsätze fachlich trennen: Brandbekämpfung, Rettung und notfallmedizinische Hilfe, technische Hilfeleistung bei Gefahrenlagen ohne Feuer sowie Alarmierungen, die sich vor Ort als gegenstandslos erweisen.\footcite[42-48]{USFA2015} Da sie unterschiedliche Mittel und Dringlichkeiten erfordern, zählt für die Vorhaltung auch ihre Zusammensetzung.

Ob die betrachteten Merkmale tragen, entscheidet sich nicht an einem absoluten Gütewert, sondern am Abstand zu dem, was ohne sie erreichbar wäre. Als Vergleichsgröße dienen deshalb zwei Baselines. Die erste kommt ohne jede Modellierung aus und besteht aus dem Gesamtmittelwert beziehungsweise der Mehrheitsklasse. Die zweite ist ein einfaches Modell ohne einstellbaren Hyperparameter, das zur Form der jeweiligen Zielgröße passt: eine Poisson-Regression mit der Wohnbevölkerung als Exposition für die Menge, ein multinomiales Logit für die Art. Als erfolgreich gilt ein Verfahren, wenn es die zweite Baseline übertrifft.

Die Arbeit ist prognoseorientiert und erhebt keinen kausalen Anspruch: Geprüft wird, ob die betrachteten Merkmale einen prädiktiven Beitrag leisten, nicht, ob sie Einsätze verursachen. Ebenso wenig wird eine Empfehlung zur Verteilung von Einsatzmitteln abgeleitet. Diese Grenze ist nicht allein eine Frage des Umfangs: Sozioökonomische Merkmale können als Stellvertreter geschützter Eigenschaften wirken, und ihre Übersetzung in Ressourcenentscheidungen erforderte eine Prüfung, die hier nicht geleistet wird. Die Ergebnisse gelten schließlich auf der Ebene des Stadtteils und lassen keine Aussagen über einzelne Gebäude oder Haushalte zu.


\section{Data Understanding}

\subsection{Datengrundlage und Merkmale}
\label{sec:datengrundlage}

Sieben öffentlich zugängliche Quellen bilden die empirische Grundlage dieser Arbeit. Im Kern stehen die Einsatzmeldungen des SFFD, die mit Alarmzeitpunkt, Stadtteil und Einsatzarten-Code die Grundlage beider Zielstränge liefern\footcite{SFFD2026}. Die Quelle reicht bis in das Jahr 2003 zurück und umfasst insgesamt 720.258 Meldungen; auf den Analysezeitraum von Januar 2015 bis Dezember 2025 entfallen davon 371.316. Die Kriminalitätsbelastung stammt aus zwei Quellen des SFPD, einer monatlich verdichteten Fassung ab Januar 2018\footcite{SFPD2026} und einer Altquelle bis Dezember 2017\footcite{SFPD2018}. Die baulichen Merkmale entstammen einem Parzellenverzeichnis des Planning Department\footcite{SFPlanning2020}, hinzu kommen die Stadtteilgeometrie\footcite{SFPlanning2023}, eine Zuordnung von Census Tracts zu Stadtteilen\footcite{SFPlanning2022} und fünf Jahrgänge der Fünfjahresschätzungen der American Community Survey\footcite{CensusACS2023}. Räumliche Bezugsebene ist die in Abschnitt~\ref{sec:anwendungsfall} begründete Stadtteilgliederung mit 41 Analysis Neighborhoods.

Aus diesem Bestand gehen zwölf Merkmale, die drei in Abschnitt~\ref{sec:zielgroessen} bestimmten Zielgrößen und eine Expositionsgröße hervor. Fünf Merkmale beschreiben die sozioökonomische Lage eines Stadtteils, eines seine Kriminalitätsbelastung und drei seinen Baubestand: der Anteil der Gebäude von vor 1940, der Anteil der Wohnparzellen und der Flächenanteil brandlastreicher Gewerbenutzung. Dazu treten die Wohnbevölkerung als Größenkontrolle und zwei Terme, die den Kalendermonat abbilden. Skala, Einheit und Wertebereich aller Größen führt Tabelle~\ref{tab:codebook} in Abschnitt~\ref{sec:merkmale} auf, wo sie entstehen.

\subsection{Explorative Analyse und Datenqualität}
\label{sec:eda}

Die folgenden Befunde beziehen sich auf den aggregierten Analysedatensatz mit der Einheit Stadtteil~×~Monat: 4.752 Zeilen, 36 der 41 Analysis Neighborhoods über die 132 Monate von Januar 2015 bis Dezember 2025; seine Konstruktion beschreibt Kapitel~\ref{sec:data-preparation}. Wo eine Angabe darüber hinausgeht und den Rohbestand vor der Aggregation betrifft, ist das vermerkt. Gerechnet wird durchgehend auf allen 36 Stadtteilen: in diesem Abschnitt werden nur die Daten beschrieben. Die Eignungsprüfung in Abschnitt~\ref{sec:eignung} entscheidet über Verfahren und rechnet deshalb nur auf den Trainingsstadtteilen des ersten Folds.

Die sieben Quellen werden in verschiedenen Takten erhoben, und der zusammengeführte Bestand erbt diese Ungleichheit (Abbildung~\ref{abb:panelstruktur}). Der Kriminalitätsindex liegt monatlich vor und nimmt je Stadtteil im Mittel 128,1 verschiedene Werte an; die sozioökonomischen Merkmale nehmen höchstens vier an, einen je Jahrgang unter dem Publikationsversatz, und die drei baulichen sind über den gesamten Zeitraum konstant. Der Bestand zählt damit 4.752 Zeilen, aber die meisten seiner Spalten ändern sich innerhalb eines Stadtteils selten oder gar nicht. Die folgenden Befunde sind Ausprägungen dieser einen Eigenschaft.

Die Zielgröße des Mengenstrangs, die Zahl der Einsätze je Stadtteilmonat, ist deutlich rechtsschief; ihre Kennzahlen stehen denen der bevölkerungsnormierten Rate gegenüber (Tabelle~\ref{tab:kennzahlen}). Ihre Varianz übersteigt den Mittelwert um den Faktor 61,8, während beide bei einer Poisson-verteilten Zählgröße übereinstimmen würden. Nullwerte treten mit einem Anteil von 0,02\,\% praktisch nicht auf --- ein einziger Stadtteilmonat blieb ohne Einsatz.

\begin{table}[htbp]
\tabformat
\caption{Lage und Streuung der beiden Mengenzielgrößen}
\label{tab:kennzahlen}
\begin{tabularx}{\textwidth}{@{}Y r r@{}}
\toprule
\textbf{Kennzahl} & \textbf{Einsatzhäufigkeit} & \textbf{Rate je 1.000 Einwohner} \\
\midrule
Mittelwert & 75,2 & 3,53 \\
Median & 53 & 2,85 \\
Maximum & 451 & 16,99 \\
Schiefe & 1,93 & 1,93 \\
Wölbung & 3,63 & 4,32 \\
\addlinespace
kleinstes Stadtteilmittel & 6,4 & 1,04 \\
größtes Stadtteilmittel & 279,7 & 10,40 \\
Verhältnis der beiden & 44 & 10 \\
\bottomrule
\end{tabularx}
\tabquelle
\end{table}

Aufschlussreicher als die Form dieser Verteilung ist ihre Herkunft. Die Stadtteilmittel reichen von 6,4 Einsätzen je Monat in Seacliff bis 279,7 im Tenderloin, während ein einzelner Stadtteil über die 132 Monate vergleichsweise wenig schwankt. Die Schiefe des Gesamtbestands entsteht damit nicht aus schwankenden Monaten, sondern aus dem Abstand zwischen den Stadtteilen (Abbildung~\ref{abb:einsatzlast}).

\begin{figure}[H]
\centering
\includegraphics[width=\textwidth]{Abbildungen/a16_einsatzlast.pdf}
\caption{Einsatzhäufigkeit je Stadtteil}
\label{abb:einsatzlast}
\end{figure}

Die zweite Zielgröße des Mengenstrangs normiert die Einsatzzahl auf die Wohnbevölkerung und prüft damit, ob dieses Gefälle lediglich ein Größeneffekt ist. Der Abstand schrumpft deutlich, verschwindet aber nicht: Die Stadtteilmittel reichen von 1,04 bis 10,40 und damit um den Faktor 10 statt um den Faktor 44. Die Einwohnerzahl erklärt das Gefälle also zu einem erheblichen Teil, aber nicht vollständig. Die Form der Verteilung ändert die Normierung dagegen kaum: Die Schiefe bleibt mit 1,93 unverändert, die Wölbung steigt leicht von 3,63 auf 4,32.

Auch die Zielgröße des Strukturstrangs ist ungleich verteilt, allerdings über die Klassen statt über die Stadtteile. Sie bezeichnet die im jeweiligen Stadtteilmonat überwiegende Einsatzart und ist auf 4.751 Zeilen definiert; einem einzigen Stadtteilmonat ohne Einsatz lässt sich keine überwiegende Art zuordnen. Auf Fehlalarme und Good-Intent-Einsätze entfallen 79,6\,\% der Stadtteilmonate, auf technische Hilfeleistungen 15,9\,\% , auf Rettungsdiensteinsätze 3,1\,\% und auf Brände 1,5\,\%. Wer stets die häufigste Klasse nennt, die erste Baseline aus Abschnitt~\ref{sec:zielgroessen}, trifft damit rechnerisch in 79,6\,\% der Fälle, ohne eine einzige Regelmäßigkeit zu nutzen.

Wie sehr der Bestand von seiner räumlichen Dimension bestimmt wird, zeigt die Zerlegung seiner Varianz. Von der Gesamtvarianz der Einsatzhäufigkeit liegen 92,5\,\% zwischen den Stadtteilen und nur 7,5\,\% innerhalb eines Stadtteils über die Zeit; der Intraklassenkorrelationskoeffizient beträgt 0,926. Für die Rate verschiebt sich das Verhältnis auf 85,1\,\% zu 14,9\,\%, ohne sich umzukehren. Bei den Merkmalen selbst liegt der Zwischenanteil zwischen 52,5\,\% und 100\,\%; die einzige Ausnahme bilden die beiden Saisonterme (Abbildung~\ref{abb:panelstruktur}).

\begin{figure}[H]
    \centering
    \includegraphics[width=\textwidth]{Abbildungen/a17_panelstruktur.pdf}
    \caption{Varianzanteile und zeitliche Auflösung der zwölf Merkmale}
    \label{abb:panelstruktur}
\end{figure}

Für den Zusammenhang der Merkmale mit der Einsatzhäufigkeit sind zwei Maße angegeben: Der Korrelationskoeffizient nach Bravais-Pearson misst den linearen, der Rangkorrelationskoeffizient nach Spearman den monotonen Zusammenhang; weichen beide voneinander ab, verläuft der Zusammenhang monoton, aber gekrümmt\footcite[140, 144-145]{Fahrmeir2023}. Die Zusammenhänge sind durchweg schwach bis mittelstark und in ihrer Richtung erwartungskonform (Tabelle~\ref{tab:zusammenhaenge}). Am deutlichsten weichen die beiden Maße bei der Armutsquote, bei der logarithmierten Wohnbevölkerung und bei der Akademikerquote voneinander ab, mit Abständen von 0,147, 0,139 und 0,130. Bei der Akademikerquote und den beiden Saisontermen bleiben dagegen beide Maße so klein, dass ihre Differenz keine Aussage trägt.

Aufschlussreicher als die Höhe der einzelnen Werte ist ihr Verhältnis untereinander. Der größte Betrag der Korrelationsmatrix entfällt mit 0,921 auf Medianeinkommen und Medianmiete, gefolgt von Medianeinkommen und Armutsquote mit −0,778 sowie Kriminalitätsindex und Risikogewerbeanteil mit 0,777. Die Merkmale beschreiben einander damit stärker als die Zielgröße --- auch dies eine Eigenschaft von Stadtteilprofilen und nicht von Stadtteilmonaten.

\begin{table}[htbp]
    \tabformat
    \caption{Zusammenhänge mit der Einsatzhäufigkeit}
    \label{tab:zusammenhaenge}
    \begin{tabularx}{\textwidth}{@{}Y r r@{}}
        \toprule
        \textbf{Merkmal} & \textbf{Pearson} & \textbf{Spearman} \\
        \midrule
        Risikogewerbeanteil          & 0,680  & 0,576  \\
        Kriminalitätsindex, log.     & 0,632  & 0,508  \\
        Leerstandsquote              & 0,524  & 0,495  \\
        Wohnbevölkerung, log.        & 0,480  & 0,619  \\
        Armutsquote                  & 0,478  & 0,331  \\
        Wohngebäudeanteil            & $-0{,}384$ & $-0{,}274$ \\
        Medianeinkommen              & $-0{,}249$ & $-0{,}168$ \\
        Medianmiete                  & $-0{,}213$ & $-0{,}112$ \\
        Altbauanteil vor 1940        & 0,169  & 0,211  \\
        Akademikerquote              & $-0{,}130$ & 0,000  \\
        Saison, Kosinus              & 0,045  & 0,051  \\
        Saison, Sinus                & $-0{,}025$ & $-0{,}048$ \\
        \bottomrule
    \end{tabularx}
    \tabquelle
\end{table}

Die Einsatzdaten selbst sind praktisch fehlerfrei: Von 720.258 Meldungen tragen 269 eine bereits vergebene Einsatznummer, also 0,04\,\%; dahinter stehen Mehrfachmeldungen desselben Einsatzes.

Im Parzellenverzeichnis fehlt bei 16.962 der 155.395 Einträge das Baujahr, also bei 10,9\,\%. Bei den sozioökonomischen Merkmalen liegt es an der American Community Survey selbst: Eine Fünfjahresschätzung bündelt fünf Erhebungsjahre und erscheint erst nach deren Ablauf, der Jahrgang 2023 etwa deckt 2019 bis 2023 ab, sodass zum Zeitpunkt eines Einsatzes im Jahr $t$ der Jahrgang $t$ noch nicht veröffentlicht ist.\footcite[13]{CensusBureau2020} Der älteste Jahrgang 2009 enthält in keinem seiner 176 Census Tracts eine Angabe zum Bildungsabschluss. Und die Zahl der Tracts steigt über die Jahrgänge von 176 auf 197 und schließlich auf 244, weil die Zensuserhebungen die Grenzen periodisch neu ziehen; die verwendete Zuordnungstabelle beruht auf den Grenzen des Zensus 2020 und umfasst 242 Tracts. Weil die genutzten Jahrgänge 2014 und 2019 noch der Einteilung von 2010 folgen, ist die Zuordnung über die Zensusgrenzen hinweg zu führen --- wie das geschieht, beschreibt Abschnitt~\ref{sec:integration}. Keine dieser Lücken ist zufällig --- sie folgen daraus, wie und wann die jeweilige Quelle erhoben wird.

Fünf der 41 untersuchten Stadtteile fehlen im Analysedatensatz, und die Gründe dafür
liegen in den Daten selbst. Der Golden Gate Park, der Lincoln Park und der McLaren Park sind Parkanlagen
ohne nennenswerte Wohnbevölkerung: Der Golden Gate Park hat in keinem der
verwendeten ACS-Jahre mehr als 63 Einwohner, während der Medianwert für die übrigen Stadtteile niemals
unter 17.916 liegt. Jede auf der Einwohnerzahl basierende Kennzahl liefert dort unplausible
Werte. Für Lakeshore und Treasure Island enthält der Parzellenindex überhaupt kein
Baujahr, sodass der Anteil älterer Gebäude dort nicht
berechnet werden kann.

Der Datenbestand ist vollständig und in sich konsistent. Seine Besonderheiten liegen in seiner Struktur. Seine 4.752 Beobachtungen unterscheiden sich fast nur darin, welchen Stadtteil sie beschreiben, kaum darin, welchen Monat sie betreffen. Eine Vorhersage kann ihre Informationen deshalb nur aus dem Profil eines Stadtteils und nicht aus dem Verlauf eines Monats ziehen. Was das für die Prüfung der Verfahren bedeutet, klärt Kapitel~\ref{sec:data-preparation}; was es für die Reichweite der Ergebnisse bedeutet, Kapitel~\ref{sec:diskussion}.

% ===================================================
\section{Data Preparation}
\label{sec:data-preparation}

% #########################################################################
% SCHREIBANLEITUNG KAPITEL 5 -- UND ABGRENZUNG DER KAPITEL 4 BIS 8
% Stand 22.08.2026. Vor dem Schreiben lesen.
% #########################################################################
%
% Dieser Block ist die einzige Stelle, an der die Abgrenzung der Kapitel
% steht. Er lag frueher vor Kapitel 4; Kapitel 4 ist geschrieben, deshalb
% ist er hierher gewandert. Verweise aus 2.5 und aus Kapitel 6 zeigen auf
% "den Block vor Kapitel 5".
%
% WARUM ES DIESEN BLOCK GIBT
% Die haeufigste Schwaeche in CRISP-DM-Arbeiten ist, dass dieselbe Sache
% zweimal auftaucht: der Befund im Data Understanding und noch einmal als
% Begruendung in der Data Preparation. Schroeter hat "keine Dopplungen,
% Tiefe richtig" ausdruecklich angemerkt (27.07.2026), und das Gutachten
% zum Anwendungsprojekt hat "umfangreich, aber nicht fokussiert"
% abgestraft (R8).
%
% MERKSATZ: Ein Befund gehoert einmal genannt und danach nur noch
% zitiert. Wer ihn beim zweiten Mal erneut herleitet, schreibt das
% Kapitel doppelt.
%
% ------------------------------------------------------------------------
% 1. DIE TRENNUNG IN EINEM SATZ JE KAPITEL
% ------------------------------------------------------------------------
%   Kap. 4  Data Understanding  BEFUND    Was ist da, was faellt auf?
%                                         -- FERTIG, nicht mehr anfassen
%   Kap. 5  Data Preparation    EINGRIFF  Was wurde damit gemacht, warum?
%   Kap. 6  Modelling           VERFAHREN Was wurde angewandt, warum das?
%   Kap. 7  Evaluation          ERGEBNIS  Wie gut war es?
%   Kap. 8  Diskussion          DEUTUNG   Was heisst das, was gilt nicht?
%
% DREI TESTFRAGEN, wenn unklar ist, wohin ein Absatz gehoert:
%   1. Gaelte der Satz auch, wenn die Daten gar nicht verarbeitet
%      wuerden?                                        -> ja: Kapitel 4
%   2. Veraendert dieser Schritt die Datei auf der Platte?
%                                                      -> ja: Kapitel 5
%   3. Passiert das je Fold, innerhalb der sklearn-Pipeline?
%                                                      -> ja: Kapitel 6
%
% ------------------------------------------------------------------------
% 2. WAS KAPITEL 4 BEREITS VERBRAUCHT HAT
% ------------------------------------------------------------------------
% Diese Befunde stehen in 4.1 und 4.2. In Kapitel 5 duerfen sie NUR NOCH
% ZITIERT werden -- mit Verweis auf Kapitel 4, ohne Zahlenherleitung.
% Rechts steht, was in Kapitel 5 dazugehoert.
%
%   Befund in Kapitel 4                 Eingriff in Kapitel 5
%   ---------------------------------   --------------------------------
%   269 doppelte Einsatznummern         Entfernung (5.1), Zahl nicht
%   (0,04 %)                            neu herleiten
%   Antwortzeit durchgehend plausibel   Fenster als Plausibilitaets-
%                                       pruefung (5.1) -- OHNE
%                                       Entfallquote, es entfaellt
%                                       keine einzige Zeile
%   10,9 % der Parzellen ohne Baujahr   Nenner yrbuilt_count statt
%                                       parcel_count (5.3)
%   ACS ist Periodenschaetzung mit      Regel acs_jahr kleinergleich
%   Publikationsversatz                 Einsatzjahr minus 1 (5.2)
%   ACS 2009 ohne Bildungsangaben       Analysebeginn 2015-01 (5.1)
%   Tract-Grenzen aendern sich,         die drei Trefferquoten
%   Crosswalk beruht auf Zensus 2020    63,1 / 79,7 / 99,2 % (5.2)
%   Golden Gate Park, 45 Einwohner      Ausschluss der drei
%                                       Parkgebiete (5.1)
%   Klassenverteilung                   argmax-Bildung der vier
%   79,0 / 16,3 / 3,1 / 1,5 %           Klassen (5.3)
%   Dispersionsindex 62,8,              Poisson-GLM als Stufe-2-
%   rechtsschiefe Zielgroesse           Baseline (5.4)
%   92,5 % Zwischenvarianz              Stadtteil-Split (5.4), steht
%                                       als Spalte IN der Datei
%   zeitliche Aufloesung 128,3 / 4 / 1  NICHTS in Kapitel 5 -- der
%                                       Befund wird in 7.4 und 8.3
%                                       zitiert
%
% ------------------------------------------------------------------------
% 3. WAS NICHT IN KAPITEL 5 GEHOERT
% ------------------------------------------------------------------------
%   Linearitaetspruefung, Residuen, RESET-Test   -> 6.2, prueft ein
%                                                   VERFAHREN, nicht die
%                                                   Daten
%   VIF / Multikollinearitaet                    -> 6.2, betrifft Ridge
%   StandardScaler, log(1+y), Label-Encoder      -> 6, laeuft je Fold
%   class_weight "balanced"                      -> 6, Modellparameter
%   Hyperparameter-Suchraeume, Tuning-Budget     -> 6
%   Fold-Ergebnisse der Verfahren                -> 7
%   Trainings- und Inferenzzeiten                -> 7
%   Vergleich der Verfahren, Rangfolge oder      -> 8
%     deren Verweigerung bei ueberlappender Streuung
%   n_eff rund 38 (ICC 0,926), Extrapolation     -> 8, Limitation
%
% WAS SEHR WOHL HIERHER GEHOERT: die BASELINES samt ihren Werten
% (R2 0,542 / RMSE 33,98 usw.). Auflage Schroeter vom 27.07.2026 --
% sie gehoeren in die Data Preparation, nicht nach 6 oder 7.
%
% NOTIZ ZUR EXPOSE-ANKUENDIGUNG (stand bisher als Kommentar bei Kap. 4):
% Das Expose kuendigte an, "bei ausgepraegten Ungleichgewichten werden
% seltene Kategorien zusammengefasst oder die Klassifikation auf zwei
% Klassen vereinfacht, etwa Brand vs. Nicht-Brand". Das ist NICHT
% eingetreten: Die Zielgroesse hat vier Klassen, die Schieflage wird
% ueber class_weight "balanced" und Macro-F1 als Hauptmass behandelt,
% nicht ueber Zusammenlegung. In Kapitel 5 steht davon nur die
% argmax-Bildung; die Behandlung gehoert nach 6, die Abweichung vom
% Expose nach 8.
%
% ------------------------------------------------------------------------
% 4. ZWEI SONDERFAELLE
% ------------------------------------------------------------------------
%  (a) 4.2 EDA gegen 6.2 Eignungspruefung. Beide schauen Verteilungen an.
%      Die Regel: In 4.2 steht, WIE DIE DATEN BESCHAFFEN SIND ("die
%      Zielgroesse ist rechtsschief, Dispersionsindex 62,8"). In 6.2
%      steht, OB EIN VERFAHREN DAZU PASST ("weil sie rechtsschief ist,
%      wird Ridge auf log(1+y) geschaetzt; der RESET-Test verwirft die
%      lineare Spezifikation"). Jede Abbildung erscheint GENAU EINMAL.
%      ACHTUNG: A16 und A17 sind in 4.2 verbraucht -- 6.2 braucht eigene
%      Abbildungen oder verweist auf Kapitel 4.
%
%  (b) 2.3 Grundlagen der Verfahren gegen 6. Kapitel 2 erklaert, WIE ein
%      Verfahren funktioniert -- allgemein, lehrbuchhaft, ohne diesen
%      Datensatz. Kapitel 6 sagt, WARUM es hier eingesetzt wird und MIT
%      WELCHEN Einstellungen. Wer in 6 erklaeren muss, wie Boosting
%      arbeitet, hat es in 2.3 versaeumt.
%
% ------------------------------------------------------------------------
% 5. BLINDABSAETZE -- GILT FUER JEDE GLIEDERUNGSEBENE DER ARBEIT
% ------------------------------------------------------------------------
% Woertlich aus der Mitschrift vom 27.07.2026:
%   "Blindabsaetze: Zwischen 1 und 1.1 steht kein Inhalt -> redundanz
%    vermeiden"
% und aus der Mitschrift vom 10.08.2026:
%   "Keine zwischenfazits (keine reudundanz)! Blindabsaetze nutzen!"
%
% Ein BLINDABSATZ ist die BEWUSST LEERE Stelle zwischen einer
% Ueberschrift und ihrer ersten Unterueberschrift. "Blindabsaetze
% nutzen" heisst: diese Stelle LEER LASSEN. Also:
%   Zwischen \section und dem ersten \subsection steht KEIN Text.
%   Zwischen \subsection und dem ersten \subsubsection steht KEIN Text.
% BEGRUENDUNG (Schroeter): Ein Absatz an dieser Stelle kuendigt an, was
% gleich ohnehin kommt -- ein vorweggenommenes Zwischenfazit. Der Inhalt
% gehoert in den ersten Unterabschnitt, wo er ohne Doppelung steht.
%
% (In 01_VORGABEN.md stand diese Auflage bis zum 18.08.2026 falsch als
% "keine Blindabsaetze zwischen Gliederungsebenen" mit der Aufloesung
% "verbindende statt zusammenfassende Absaetze". Beides steht in keiner
% Mitschrift und ist gegen die Primaerquellen korrigiert worden. Einen
% Widerspruch zwischen dem 27.07. und dem 10.08. gab es nie.)
%
% OFFENER VERSTOSS -- GENAU HIER, direkt unter diesem Block:
% Die folgenden 13 Zeilen Fliesstext stehen zwischen
% \section{Data Preparation} und \subsection{Datenauswahl und
% Bereinigung}. Inhaltlich sind sie richtig und noetig -- sie tragen die
% zentrale Entwurfsentscheidung (Analyseeinheit) und den
% Verfuegbarkeitsgrundsatz. Zwei Loesungen:
%   a) an den Anfang von 5.1 ziehen -- billig, aber 5.1 heisst
%      "Datenauswahl und Bereinigung", und die Analyseeinheit ist keine
%      Bereinigung
%   b) einen eigenen Abschnitt 5.1 "Analyseeinheit und Grundsatz" davor
%      setzen, die uebrigen um eins verschieben -- sauberer
% EMPFEHLUNG: b). Alle uebrigen Kapitel sind regelkonform.
%
% ------------------------------------------------------------------------
% 6. ARBEITSWEISE FUER DIESES KAPITEL
% ------------------------------------------------------------------------
% Alle ZAHLEN kommen aus docs/03_STAND.md -- dort stehen sie, und nur
% dort. Nicht aus aelteren Fassungen abschreiben. Die Pipeline wurde am
% 28.07.2026 vollstaendig umgebaut; der fruehere Text dieses Kapitels ist
% inhaltlich ueberholt und wurde entfernt.
% Reihenfolge folgt der Pipeline: prep/s1_daten.py -> s2_datensaetze.py
% -> s3_baselines.py
% #########################################################################

\subsection{Datenauswahl und Bereinigung}
\label{sec:auswahl-bereinigung}

Der Rohbestand aus Abschnitt~\ref{sec:datengrundlage} umfasst 720.258
Einsatzmeldungen aus 41 Analysis Neighborhoods. Nicht alles davon darf und
nicht alles kann in die Analyse eingehen. Die erste Frage betrifft die
Zulässigkeit einer Größe und folgt dem Grundsatz aus
Abschnitt~\ref{sec:modellbewertung}, dass ein Prädiktor zum Prognosezeitpunkt
verfügbar gewesen sein muss; die zweite grenzt den Gegenstand zeitlich und
räumlich ab und stützt sich auf die Qualitätsbefunde aus
Abschnitt~\ref{sec:eda}. Beide Entscheidungen sind nicht über den
Verarbeitungscode verteilt, sondern als benannte Konstanten hinterlegt;
jede ist damit an einem Ort nachlesbar und an
einem Ort änderbar. Zwei von ihnen halten nicht eine Einstellung fest, sondern
das erwartete Ergebnis: \texttt{N\_STADTTEILE\_ERWARTET} und die
Plausibilitätsspanne der stadtweiten Wohnbevölkerung. Sie werden nach der
Aufbereitung geprüft, und eine Abweichung bricht den Lauf ab, statt eine still
verkleinerte Bezugsgröße in die Modelle zu tragen (Code~\ref{lst:zuschnitt}). Am Ende steht der Rahmen, in
dem gerechnet wird: 36 Stadtteile über 132 Monate.
 
% NOTIZ (07.09.2026) -- ersetzt den frueheren Ausschnitt mit der Konstantenliste.
% Gezeigt wird nicht mehr, DASS die Festlegungen benannt sind, sondern dass sie
% geprueft werden. Beide Konstanten aus dem Absatz darueber tauchen hier wieder
% auf: N_STADTTEILE_ERWARTET in Zeile 433 und BEV_PLAUSIBEL in Zeile 439.
% Erklaerbar im Text: (1) die zweite Pruefung ist die eigentliche -- ein Verbund,
% der Zeilen verwirft, verliert Stadtteile UND Bevoelkerung, und nur die
% Bevoelkerung faellt auf, wenn die Stadtteilzahl zufaellig stimmt; (2) die
% Fehlermeldung nennt die fehlenden Stadtteile beim Namen (Z. 435) und verweist
% auf den Crosswalk (Z. 443) -- sie sagt, wo zu suchen ist, nicht nur, dass
% etwas falsch ist; (3) die Pruefung laeuft NACH der Aufbereitung, nicht in ihr.
\begin{quellcode}[htbp]
\setstretch{1}
\caption{Prüfung von Zuschnitt und Wohnbevölkerung}
\label{lst:zuschnitt}\vspace{-3.5pt}
\begin{lstlisting}[style=pythonstil,firstnumber=417]
def pruefe_zuschnitt(r: pd.DataFrame) -> None:(*@\setcounter{lstnumber}{429}@*)
    ist = sorted(r["stadtteil"].unique())
    alle = sorted(pd.read_parquet(PFAD_EINSAETZE, columns=["stadtteil"])
                    ["stadtteil"].dropna().unique())
    assert len(ist) == N_STADTTEILE_ERWARTET, (
        f"{len(ist)} Analyseeinheiten statt {N_STADTTEILE_ERWARTET}. "
        f"Nicht enthalten: {sorted(set(alle) - set(ist))}")

    bev = (r.groupby(["jahr", "stadtteil"])[EXPOSURE_ROH].first()
             .groupby("jahr").sum())
    unten, oben = BEV_PLAUSIBEL
    assert unten <= bev.min() and bev.max() <= oben, (
        f"Stadtweite Wohnbevoelkerung {bev.min():,.0f} bis {bev.max():,.0f} "
        f"ausserhalb {unten:,}-{oben:,} - ACS-Zuordnung pruefen "
        f"(Tract-Grenzen der Jahrgaenge gegen den Crosswalk).")
\end{lstlisting}
\tabquelle[Eigene Darstellung, \texttt{prep/s2\_datensaetze.py}, Z.~417--443, gekürzt]
\end{quellcode}
 
In die Analysedatensätze gelangen ausschließlich Größen, die zum Zeitpunkt des
Alarms bereits feststanden --- das no-time-machine
requirement.\footcite[559]{Kaufman2011} Damit sind zehn Felder des
Feldverzeichnisses gesperrt: geschätzter Sachschaden, Löschfahrzeuge,
Löschkräfte, Rettungsdiensteinheiten, Alarmstufe, Antwortzeit, Ankunftszeitpunkt, zivile Tote, zivile Verletzte und die Angabe, ob die Flammenausbreitung eingedämmt wurde. Sie sind nicht nur unzulässig,
sondern gefährlich: Sie hängen eng mit der Zielgröße zusammen, und ein Modell
mit ihnen erreichte hohe Gütewerte, ohne eine Prognose zu leisten. Der
Ausschluss geschieht deshalb beim Laden und nicht erst beim Modellieren ---
was nicht in die Datei gelangt, kann später nicht versehentlich in ein Modell
geraten.\footcite[560]{Kaufman2011} Dass keines dieser Felder den Weg in die
erzeugten Dateien findet, wird nicht behauptet, sondern geprüft
(Code~\ref{lst:test-leakage}).
 
\begin{quellcode}[htbp]
\setstretch{1}
\caption{Prüfung auf Ergebnisvariablen in der erzeugten Datei}
\label{lst:test-leakage}\vspace{-3.5pt}
\begin{lstlisting}[style=pythonstil,firstnumber=279]
def test_keine_ergebnisvariablen():(*@\setcounter{lstnumber}{285}@*)
    gefunden = [c for c in ERGEBNISVARIABLEN if c in klassifikation().columns]
    assert not gefunden, f"Ergebnisvariable(n) im Datensatz: {gefunden}"
\end{lstlisting}
\tabquelle[Eigene Darstellung, \texttt{tests/test\_aufbereitung.py}, Z.~279--287, gekürzt]
\end{quellcode}
 
Bereinigt wird nur, was nachweislich Artefakt ist. 269 doppelte Einsatznummern
werden entfernt (Abschnitt~\ref{sec:eda}); Ausrückzeiten außerhalb von 0 bis 60
Minuten und Baujahre außerhalb von 1800 bis 2025 werden als fehlend markiert
statt gelöscht. Beide Grenzen sind gesetzte Plausibilitätsannahmen und folgen
nicht aus den Daten. An der Zielgröße selbst wird nichts gekappt oder getrimmt;
Bereinigung betrifft ausschließlich Platzhalterwerte der Quellen.
 
Fünf Stadtteile scheiden ganz aus: drei Parkgebiete ohne nennenswerte
Wohnbevölkerung sowie Lakeshore und Treasure Island, für die das
Parzellenverzeichnis kein einziges Baujahr führt, sodass der Altbauanteil dort
nicht bildbar ist. Entscheidend ist die Form
dieses Ausschlusses: Entfernt werden vollständige Analyseeinheiten, nicht
einzelne Zeilen. Ein zeilenweises Verwerfen erzeugte ein unbalanciertes Panel,
in dem ein Stadtteil mitten in der Zeitreihe hinzuträte und Summen im
Testfenster allein durch diesen Zutritt sprängen. Ein rechteckiges Panel ist
Voraussetzung dafür, dass die Folds aus Abschnitt~\ref{sec:rahmen-baselines}
miteinander vergleichbar sind.
 
Der Analysezeitraum ist auf Januar 2015 bis Dezember 2025 festgelegt, also auf
132 Monate. Der Beginn ist keine Setzung, sondern eine Folge der
Verfügbarkeitsregel: Für Einsätze vor 2015 wäre der jüngste zulässige
ACS-Jahrgang der von 2009, und der führt die Tabelle nicht, aus der die
Akademikerquote stammt (Abschnitt~\ref{sec:eda}). Das Ende ist das letzte
vollständige Kalenderjahr. Beide Grenzen sind fest hinterlegt und werden nicht
aus den Daten abgeleitet, weil die Analyse sonst mit jedem neuen Abzug wanderte
und ein angebrochener Randmonat unbemerkt ins Testfenster geriete.
 
% INHALT:
% - Strikte Spaltenauswahl: von 23 Einsatzspalten bleiben 6. Ausgeschlossen
%   werden Sachschaden, Loeschfahrzeuge, Loeschkraefte, Rettungsdiensteinheiten,
%   Alarmstufe, zivile Tote/Verletzte, Flammenausbreitung -- sie stehen erst
%   NACH dem Einsatz fest und waeren Leakage im engeren Sinn. Argument: Was
%   nicht in den Datensatz kommt, kann nicht versehentlich ins Modell geraten.
% - Bereinigung: 269 doppelte Einsatznummern (0,04 %) entfernt -> 719.989
%   Einsaetze; Ausrueckzeit auf 0-60 min begrenzt (~1,7 % entfallen);
%   Baujahre ausserhalb 1800-2025 als fehlend markiert.
% - Ausschluss ganzer Analyseeinheiten (41 -> 35 Stadtteile):
%     * Golden Gate Park, Lincoln Park, McLaren Park -- Parkgebiete mit 45 bis
%       507 Einwohnern; jede Pro-Kopf-Groesse wird dort beliebig gross
%       (Median der uebrigen Stadtteile: 14.435 Einwohner)
%     * Treasure Island, Lakeshore -- in KEINEM ACS-Jahrgang enthalten
%     * Mission Bay -- erst ab ACS 2021; ein Zutritt mitten in der Zeitreihe
%       erzeugte ein unbalanciertes Panel
%   Zu betonen: Unterschied zwischen zeilenweisem dropna und dem Ausschluss
%   ganzer Einheiten.
% - Analysezeitraum 2015-01 bis 2025-12, hart in config.py fixiert
% ZU FORMULIEREN -- der Start ist eine KONSEQUENZ, keine Setzung:
%   "Der Analysebeginn 2015-01 folgt nicht aus einer Setzung, sondern aus der
%    Publikationsrealitaet des ACS."
%   Herleitung: Die Regel lautet acs_jahr <= Einsatzjahr - 1. Verfuegbar sind
%   die Jahrgaenge 2009, 2014, 2019, 2021, 2023. Ein Einsatz aus 2014 bekaeme
%   damit ACS 2009 - und der enthaelt die Tabelle B15003 nicht, aus der die
%   akademikerquote_pct stammt. Erst ab 2015 greift ACS 2014 (Decision Log #11).
%   NICHT als Wahl darstellen: Es gab keine. und nicht
%   aus den Daten abgeleitet (sonst wandert die Analyse mit jedem Download,
%   und ein angebrochener Randmonat geraet unbemerkt ins Testfenster).
% KEIN Code-Listing.


\subsection{Zusammenführung der Datenquellen}
\label{sec:integration}

Die sieben Quellen aus Abschnitt~\ref{sec:datengrundlage} liegen in drei
Raumbezügen und zwei Zeitrastern vor. Sie werden auf der Ebene des einzelnen
Einsatzes zusammengeführt und nicht auf der späteren Analyseeinheit: Jeder
Einsatz trägt seinen eigenen Zeitpunkt, und nur auf dieser Ebene lässt sich ihm
derjenige Stand der zeitabhängigen Merkmale zuordnen, der damals verfügbar war.
Verdichtet wird erst danach (Abschnitt~\ref{sec:merkmale}); die Reihenfolge
folgt damit der Verarbeitung und nicht der Aufgabenliste des Vorgehensmodells.
Das Ergebnis ist eine Tabelle mit 719.989 Zeilen und 50 Spalten, eine je
Einsatz.
 
Die sozioökonomischen Merkmale werden in zwei Stufen angefügt. Räumlich ordnet
ein Crosswalk jeden Census Tract einem Stadtteil zu; dabei werden Zählgrößen
summiert, Medianwerte dagegen bevölkerungsgewichtet gemittelt, weil sich ein
Median nicht addieren lässt. Zeitlich gilt die Bedingung
$\text{ACS-Jahrgang} \le \text{Einsatzjahr} - 1$, sodass einem Einsatz nur ein
Jahrgang zugeordnet wird, der zu seinem Zeitpunkt bereits veröffentlicht war
(Code~\ref{lst:acs-versatz}). Der Crosswalk beruht auf den Tract-Grenzen
des Zensus 2020, die Jahrgänge 2014 und 2019 dagegen auf denen des Zensus 2010;
über die Kennung allein greift die Zuordnung dort nur für 79,7\,\% der Tracts.
Aufgefangen wird das über den vierstelligen Basiscode der Tract-Nummer: Eine
Teilung im Zensus 2020 behält ihn bei und ändert nur das zweistellige Suffix,
sodass sich ein Tract der alten Einteilung seinen Nachfolgern zuordnen lässt.
Die Rückfallebene greift nur, wenn sämtliche Nachfolger desselben Basiscodes in
denselben Stadtteil fallen, und lässt den Tract sonst offen --- sie rät nicht.
Damit steigt die Trefferquote auf 99,2 bis 99,5\,\% je Jahrgang; die
Aufbereitung bricht ab, sobald weniger als 95\,\% der Einwohner eines Jahrgangs
einem Stadtteil zugeordnet sind. Die stadtweite Wohnbevölkerung der enthaltenen
Stadtteile liegt danach zwischen 811.572 und 856.270 und damit im erwarteten
Bereich.

\begin{quellcode}[htbp]
\setstretch{1}
\caption{Zuordnung über Zensusgrenzen hinweg}
\label{lst:crosswalk}\vspace{-3.5pt}
\begin{lstlisting}[style=pythonstil,firstnumber=282]
def tract_zu_stadtteil(geoids: pd.Series, crosswalk: pd.DataFrame) -> pd.Series:(*@\setcounter{lstnumber}{302}@*)
    direkt = dict(zip(crosswalk["geoid"], crosswalk["neighborhood"]))
    basis = (crosswalk.assign(basis=crosswalk["geoid"].str[5:9])
                      .groupby("basis")["neighborhood"]
                      .agg(lambda s: s.iloc[0] if s.nunique() == 1 else None)
                      .dropna().to_dict())
    g = geoids.astype(str).str.zfill(11)
    return g.map(direkt).fillna(g.str[5:9].map(basis))(*@\setcounter{lstnumber}{330}@*)
    assert quote >= 0.95, (
        f"ACS-Zuordnung deckt nur {quote:.1%} der Wohnbevoelkerung - "
        f"Crosswalk gegen die Tract-Grenzen des Jahrgangs pruefen.")
\end{lstlisting}
\tabquelle[Eigene Darstellung, \texttt{prep/s1\_daten.py}, Z.~282--333, gekürzt]
\end{quellcode}

Die Eindeutigkeitsbedingung steht in der Aggregation: \texttt{nunique() == 1}
liefert nur dann einen Stadtteil, wenn alle Nachfolger desselben Basiscodes in
denselben fallen, und sonst \texttt{None}. Für die Jahrgänge 2014 und 2019
trifft das auf alle 39 betroffenen Tracts zu; offen bleiben danach allein die
unbewohnten Wasserflächen-Tracts. Die Schwelle von 95\,\% ist keine
Sorgfaltsroutine, sondern die Lehre aus einem Fehler: Ein reiner
Gleichheitsverbund verwarf 40 der 197 Tracts der Jahrgänge 2014 und 2019 --- und
mit ihnen gut ein Viertel der Wohnbevölkerung, konzentriert in den geteilten
Tracts der Innenstadt. Eine Zuordnung, die stillschweigend Zeilen verliert,
fällt in keiner Gütekennzahl auf; sie muss beim Einlesen auffallen.
 
\begin{quellcode}[htbp]
\setstretch{1}
\caption{Zeitbewusste Zuordnung des ACS-Jahrgangs}
\label{lst:acs-versatz}\vspace{-3.5pt}
\begin{lstlisting}[style=pythonstil,firstnumber=349]
def acs_snapshot(jahr: int, acs_years: list[int]) -> int:(*@\setcounter{lstnumber}{363}@*)
    return max([a for a in acs_years if a <= int(jahr) - ACS_PUBLIKATIONS_LAG],
               default=min(acs_years))(*@\setcounter{lstnumber}{379}@*)
sffd["acs_year"] = sffd["year"].apply(lambda y: acs_snapshot(y, jahrgaenge))
\end{lstlisting}
\tabquelle[Eigene Darstellung, \texttt{prep/s1\_daten.py}, Z.~349--380, gekürzt]
\end{quellcode}
 
Die Kriminalitätsbelastung stammt aus zwei Quellen, deren Schnitt im Januar
2018 liegt. Die Altquelle führt keine Stadtteilspalte, dafür Koordinaten; ihre
Meldungen werden über einen räumlichen Verbund derselben Geometrie zugeordnet,
die auch für die baulichen Merkmale gilt --- sonst bezögen sich Kriminalitäts-
und Baumerkmale desselben Stadtteils auf unterschiedliche Flächen.
Geschnitten wird am Beginn der
Nachfolgequelle, damit kein Monat aus zwei Erfassungssystemen gespeist wird.
 
Die baulichen Merkmale liegen je Parzelle vor und werden über den
Flächenmittelpunkt zugeordnet; das gelingt für 154.544 Parzellen und damit für
99,5\,\% des Verzeichnisses. Der Bestand ist ein Abzug des Jahres 2020 und der
einzige verfügbare Jahrgang. Die baulichen Merkmale sind deshalb je Stadtteil
über den gesamten Analysezeitraum konstant (Abschnitt~\ref{sec:eda}).
 
Alle drei Zusammenführungen sind nach demselben Muster abgesichert: Vor jedem
Verbund wird die Eindeutigkeit des Schlüssels geprüft, nach jedem die Zeilenzahl
gegen den Erwartungswert gehalten, und bei Abweichung bricht die Aufbereitung
ab --- ebenso, wenn die Zuordnungsquote des
räumlichen Verbunds unter 90\,\% fällt. Das ist keine Sorgfaltsroutine: Ist ein
Schlüssel mehrdeutig, entsteht ein kartesisches Produkt, und jeder betroffene
Einsatz geht mehrfach in die spätere Aggregation ein.

% INHALT (knapp halten, KEIN Code-Listing):
% - Join A, ACS: zwei Stufen. Census-Tract -> Stadtteil ueber den Crosswalk
%   (Mediane bevoelkerungsgewichtet, Zaehler summiert); dann ZEITBEWUSST mit
%   Publikationsversatz: acs_jahr <= Einsatzjahr - 1. Ohne den Versatz bekaeme
%   ein Einsatz aus 2023 den Jahrgang 2023, der erst Ende 2024 erschien.
%   Als Formel im Fliesstext, nicht als Quellcode.
%   Limitation: Trefferquoten je Jahrgang 2009: 63,1 %; 2014/2019: 79,7 %;
%   2021/2023: 99,2 % (Census-Tract-Grenzen aendern sich, Crosswalk von 2020).
% - Join B, Kriminalitaet: zwei SFPD-Quellen mit Systembruch 05/2018 (CABLE ->
%   Crime Data Warehouse). Die aeltere Quelle hat keine Stadtteilspalte, dafuer
%   Koordinaten -> Spatial Join gegen DIESELBE Neighborhood-Geometrie wie bei
%   Land Use. Sonst bezoegen sich Kriminalitaets- und Baumerkmale desselben
%   Stadtteils auf unterschiedliche Flaechen.
% - Join C, Land Use: 154.544 Parzellen, Zuordnung ueber den Parzellen-
%   Mittelpunkt, Match-Rate 99,5 %. Snapshot 2020, einziger Jahrgang ->
%   zeitkonstantes Strukturmerkmal.
% - Join-Hygiene: Schluesseleindeutigkeit vor jedem Merge, validate-Parameter,
%   Match-Quoten, Zeilenzahl vor/nach gegen den Erwartungswert. Schutz vor
%   kartesischen Produkten.
% ERGEBNIS: einsaetze.parquet, 719.989 x 50, ein Einsatz je Zeile.

% =========================================================================
\subsection{Konstruktion der Merkmale und Zielgrößen}
\label{sec:merkmale}
Die Einsatztabelle aus Abschnitt~\ref{sec:integration} führt 719.989 Zeilen,
eine je Ereignis. Auf dieser Ebene stehen Merkmale und Zielgröße nicht auf
derselben Beobachtung: Die Strukturmerkmale beschreiben einen Stadtteil, die
Zielgröße dessen Einsatzlast. Die Verdichtung auf Stadtteil und Monat stellt
beides in dieselbe Zeile und bildet dabei die Größen, mit denen gerechnet wird.
Eine Verdichtung auf Census Tract und Jahr scheidet aus, weil die Einsatzdaten
keine Tract-Zuordnung tragen und die Saisonalität dabei verloren ginge.
Tabelle~\ref{tab:codebook} führt die entstehenden Größen mit Skalenniveau,
Einheit und Wertebereich; die folgenden Absätze begründen ihre Bildung.
 
Verdichtet wird auf ein vollständiges Raster aus 36 Stadtteilen und 132
Monaten: Auch ein Monat ohne Einsatz erhält eine Zeile, und zwar mit dem Wert
null statt gar nicht. Fehlende Merkmalswerte werden ausschließlich vorwärts
fortgeschrieben; ein Rückwärtsfüllen schlösse eine Lücke mit einem später
erhobenen Wert und spielte damit genau die Information ein, die zum
Prognosezeitpunkt nicht vorlag (Code~\ref{lst:raster}). Am Ende stehen
4.752 Zeilen für die Menge der Einsatzlast und
4.751 für ihre Zusammensetzung; dem Strukturdatensatz fehlt der eine Monat ohne
Einsatz.
 
\begin{quellcode}[htbp]
\setstretch{1}
\caption{Vollständiges Raster und Vorwärtsfüllen}
\label{lst:raster}\vspace{-3.5pt}
\begin{lstlisting}[style=pythonstil,firstnumber=223]
    idx = pd.MultiIndex.from_product(
        [agg["stadtteil"].unique(),
         range(int(agg["jahr"].min()), int(agg["jahr"].max()) + 1),
         range(1, 13)],
        names=["stadtteil", "jahr", "monat"])
    raster = agg.set_index(["stadtteil", "jahr", "monat"]).reindex(idx).reset_index()
    raster[ZIELGROESSE] = raster[ZIELGROESSE].fillna(0).astype(int)(*@\setcounter{lstnumber}{233}@*)
    raster[_UEBERNOMMEN] = (raster.groupby("stadtteil")[_UEBERNOMMEN]
                                  .transform(lambda s: s.ffill()))
\end{lstlisting}
\tabquelle[Eigene Darstellung, \texttt{prep/s2\_datensaetze.py}, Z.~223--235, gekürzt]
\end{quellcode}
 
Sechs der zehn Strukturmerkmale sind Anteilswerte aus je zwei Zählgrößen
derselben Erhebung. Ein Nenner von null ergibt dabei einen fehlenden Wert und
nicht null, damit eine fehlende Bezugsgröße sichtbar bleibt
(Code~\ref{lst:quoten}). Welche Zählgröße zu welchem Nenner gehört, steht dabei als benannte Liste an einer einzigen Stelle und nicht verstreut in sechs Einzelrechnungen. Ein Nenner ist begründungspflichtig: Der
Altbauanteil bezieht sich auf die Parzellen mit bekanntem Baujahr, weil er
sonst dort am stärksten zu niedrig ausfiele, wo das Verzeichnis am
lückenhaftesten ist (Abschnitt~\ref{sec:eda}).
 
\begin{quellcode}[htbp]
\setstretch{1}
\caption{Anteilswerte aus Zählvariablen}
\label{lst:quoten}\vspace{-3.5pt}
\begin{lstlisting}[style=pythonstil,firstnumber=613]
def berechne_quoten(df: pd.DataFrame) -> pd.DataFrame:(*@\setcounter{lstnumber}{621}@*)
    for name, zaehler, nenner in QUOTEN:
        z = pd.to_numeric(df[zaehler], errors="coerce").astype(float)
        n = pd.to_numeric(df[nenner], errors="coerce").astype(float)
        df[name] = (z / n.where(n > 0, np.nan)).round(4)(*@\setcounter{lstnumber}{628}@*)
    return df
\end{lstlisting}
\tabquelle[Eigene Darstellung, \texttt{prep/s1\_daten.py}, Z.~613--629, gekürzt]
\end{quellcode}
 
Das kriminalitätsbezogene Merkmal ist kein Rohwert, sondern ein
Standortquotient: die Deliktrate eines Stadtteils, gemessen an der Deliktrate
der Gesamtstadt im selben Monat.
 
\begin{equation*}
    \mathrm{Rate}(i,t) = \frac{D(i,t)}{E(i,t)}
    \qquad
    \mathrm{Index}(i,t) = \frac{\mathrm{Rate}(i,t)}{\mathrm{Rate}(\mathrm{Stadt},t)}
\end{equation*}
 
Dabei ist $D(i,t)$ die Zahl der Delikte im Stadtteil $i$ über die zwölf Monate
bis einschließlich $t-1$ und $E(i,t)$ seine Einwohnerzahl im zugeordneten
Zensusjahrgang; ein Wert von eins bedeutet eine Belastung wie im
Stadtdurchschnitt desselben Monats. Die Form des Standortquotienten stammt aus
der Regionalwissenschaft und wird von Brantingham und Brantingham auf
Kriminalität übertragen.\footcite[268-269]{Brantingham1997} Deren Fassung
verwendet allerdings die Gesamtzahl der Delikte als Bezugsgröße und misst
damit, worauf ein Gebiet relativ spezialisiert ist; hier ist die Bezugsgröße
die Wohnbevölkerung, gemessen wird die relative Belastung je Einwohner.
 
Zwei Eigenschaften dieser Konstruktion sind Entscheidungen. Die relative Form
ist die Antwort auf den Quellenwechsel aus Abschnitt~\ref{sec:integration}: Ein
Wechsel des Erfassungssystems verschiebt das stadtweite Niveau annähernd
multiplikativ und kürzt sich deshalb weitgehend heraus. Was bleibt, ist eine
Verschiebung in der Zusammensetzung der erfassten Delikte, die einzelne
Stadtteile stärker trifft als andere. Die zweite Eigenschaft ist das Fenster:
Es endet strikt im Vormonat, weil eine im selben Monat gemessene Kriminalität
die Einsätze dieses Monats erklären statt vorhersagen würde
(Code~\ref{lst:krimindex}). Der Index geht logarithmiert ein, sodass der
Stadtdurchschnitt bei null liegt.
 
\begin{quellcode}[htbp]
\setstretch{1}
\caption{Relativer Kriminalitätsindex je Stadtteil und Monat}
\label{lst:krimindex}\vspace{-3.5pt}
\begin{lstlisting}[style=pythonstil,firstnumber=508]
    raster["delikte_fenster"] = (
        raster.groupby("neighborhood")["delikte"]
              .transform(lambda s: s.shift(1).rolling(CRIME_FENSTER_MONATE).sum()))
    raster = raster.dropna(subset=["delikte_fenster"])(*@\setcounter{lstnumber}{521}@*)
    stadt = (raster.groupby(["jahr", "monat"])
                   .agg(stadt_delikte=("delikte_fenster", "sum"),
                        stadt_bev=("total_population", "sum")).reset_index())
    raster = raster.merge(stadt, on=["jahr", "monat"], how="left")
    raster["crime_rate_raw"] = (raster["delikte_fenster"]
                                / raster["total_population"] * 1000).round(3)
    raster["crime_index"] = (raster["crime_rate_raw"]
                             / (raster["stadt_delikte"] / raster["stadt_bev"] * 1000)
                             ).round(4)
\end{lstlisting}
\tabquelle[Eigene Darstellung, \texttt{prep/s1\_daten.py}, Z.~508--530, gekürzt]
\end{quellcode}
 
Die Wohnbevölkerung geht als eigenes Merkmal ein, ebenfalls logarithmiert. Sie
wirkt auf die beiden Mengenzielgrößen gegenläufig: Ein größerer Stadtteil hat
mehr Einsätze, aber nicht zwangsläufig mehr Einsätze je Einwohner. Ohne diese
Größenkontrolle bildete ein Modell auf der absoluten Zielgröße vor allem die
Einwohnerzahl ab und nicht die Struktur, nach der gefragt ist. Im
Poisson-Modell des Abschnitts~\ref{sec:rahmen-baselines} tritt sie zusätzlich
als Offset auf.
 
Der Kalendermonat wird nicht als Zahl von eins bis zwölf geführt, sondern als
Sinus- und Kosinusterm (Code~\ref{lst:saison}). Als Zahl hätten Dezember und Januar den Abstand elf,
obwohl sie aufeinanderfolgen. Die beiden Terme sind die einzigen Merkmale, die
ausschließlich innerhalb der Stadtteile variieren (Abschnitt~\ref{sec:eda}).

% NOTIZ (07.09.2026) -- neu, der kuerzeste Ausschnitt der Arbeit und bewusst so.
% Er zeigt, dass eine Modellannahme -- Dezember und Januar sind benachbart --
% in zwei Zeilen kodiert ist und nicht in einer Vorverarbeitungsroutine steckt.
% Erklaerbar im Text (mit Zeilenbezug):
% (1) Zeile 274/275 bilden aus einem Monat zwei Merkmale; die Periode 12 im
%     Nenner ist die Annahme, nicht das Ergebnis einer Schaetzung;
% (2) beide Terme zusammen sind noetig -- ein Sinus allein waere im Jahr nicht
%     eindeutig, Monat 3 und Monat 9 haetten denselben Wert;
% (3) der Kommentar in Zeile 273 verweist auf die Konstante SAISON in config.py:
%     die Begruendung steht bei der Festlegung und nicht bei der Rechnung.
\begin{quellcode}[htbp]
\setstretch{1}
\caption{Der Kalendermonat als Sinus und Kosinus}
\label{lst:saison}\vspace{-3.5pt}
\begin{lstlisting}[style=pythonstil,firstnumber=273]
    # SAISON als sin/cos statt Monat 1-12 (Begruendung in config.py, SAISON).
    d["monat_sin"] = np.sin(2 * np.pi * d["monat"] / 12)
    d["monat_cos"] = np.cos(2 * np.pi * d["monat"] / 12)
\end{lstlisting}
\tabquelle[Eigene Darstellung, \texttt{prep/s2\_datensaetze.py}, Z.~273--275]
\end{quellcode}
 
Drei Verlaufsmerkmale --- Vormonat, Vorjahresmonat und Mittel der drei
Vormonate --- werden strikt rückwärtsgerichtet und je Stadtteil gebildet; damit
sie schon für den ersten Analysemonat definiert sind, beginnt die Aggregation
zwölf Monate früher, ohne dass diese Vorlaufmonate eigene Zeilen bilden. In die
Modelle gehen sie dennoch nicht ein: Unter der Aufteilung des
Abschnitts~\ref{sec:rahmen-baselines} wäre der Vormonatswert die eigene
Vergangenheit eines zurückgehaltenen Stadtteils. Sie verbleiben im Datensatz,
weil die Beschreibung in Abschnitt~\ref{sec:datengrundlage} sie benötigt.
 
Drei Zielgrößen stehen nebeneinander (Abschnitt~\ref{sec:zielgroessen}): die
Zahl der Einsätze je Stadtteilmonat, dieselbe Zahl auf je 1.000 Einwohner
normiert und die überwiegende Einsatzart. Die vier Anteile der NFIRS-Gruppen
sind dabei Zwischengröße und keine eigenen Zielgrößen; aus ihnen entsteht die
überwiegende Einsatzart als der größte der vier.Bei Gleichstand entscheidet die Reihenfolge der vier Anteilsspalten. Sie folgt der fachlichen Ordnung der NFIRS-Gruppen und stand vor dem ersten Modelllauf fest; der Gleichstand fällt damit stets zugunsten der selteneren Klasse. Betroffen sind 147 der 4.751 Zeilen, darunter 13 der 70 brandbestimmten und 41 der 145 rettungsbestimmten Monate. Lage und Streuung der beiden
Mengenzielgrößen stehen in Abschnitt~\ref{sec:eda}.
 
Dass die Einsatzart auf Stadtteilebene und nicht je Einsatz bestimmt wird, ist
die folgenreichste Festlegung dieses Abschnitts. Innerhalb eines
Stadtteilmonats tragen alle Einsätze dieselben Strukturmerkmale; die 357.553
Einsätze des Analysezeitraums verteilen sich deshalb auf nur 4.751 verschiedene
Merkmalsprofile. Ein Modell, das jedem Profil die dort häufigste Einsatzart
zuwiese, träfe in 50,0\,\% der Fälle zu, die häufigste Klasse ohne jedes
Merkmal zu nennen in 48,3\,\% --- der gesamte Spielraum beträgt 1,7
Prozentpunkte. Auf Stadtteilebene ist die Frage dagegen beantwortbar. Die
Zielgröße bleibt dabei eine echte Klasse und kein an einem Schwellwert
geteilter stetiger Anteil; damit entfällt die Begründungslast einer solchen
Teilung, deren Wirkung Altman und Royston mit dem Verlust eines Drittels der
Beobachtungen beziffern.\footcite{Altman2006}
 
% LAYOUT -- nicht ohne Nachrechnen aendern. Die Fassung mit fuenf Spalten
% (Einheit als eigene Spalte) war 119 pt zu hoch fuer eine Seite. Deshalb:
% Einheit steht in der Spalte Wertebereich, \singlespacing haelt die Zeilen
% zusammen (\onehalfspacing wirkt sonst auch in der Tabelle), und
% \tabcolsep 5 pt schafft der Skalenspalte Platz, damit "Verhaeltnis" nicht
% umbricht. Neu am 25.08.: aboverulesep/belowrulesep auf 0 und die
% Kriminalitaetszeile ohne "0 ist der Stadtdurchschnitt" -- zusammen rund
% 16 pt, damit meldet LaTeX den Float nicht mehr als zu hoch.
% Quelle der Angaben: results/codebook/merkmale.md, erzeugt von
% tools/codebook.py. Erfuellt die Auflage Schroeter vom 10.08.2026 und loest
% P1 aus docs/01_VORGABEN.md.
\begin{table}[htbp]
    \tabformat
    \caption{Merkmale und Zielgrößen des Analysedatensatzes}
    \label{tab:codebook}
    \begin{tabularx}{\textwidth}{@{}L{3.4cm} L{2.0cm} L{3.6cm} Y@{}}
        \toprule
        \textbf{Größe} & \textbf{Skala} & \textbf{Wertebereich und Einheit} & \textbf{Quelle und Bildung} \\
        \midrule
        \multicolumn{4}{@{}l}{\textbf{Merkmale --- sozioökonomisch}} \\
        Medianeinkommen & Verhältnis & 21.110 bis 246.635 USD je Jahr & ACS B19013, bevölkerungsgewichtet \\
        Armutsquote & Verhältnis & 0,007 bis 0,321 & ACS B17001, Quote \\
        Akademikerquote & Verhältnis & 0,139 bis 0,553 & ACS B15003, Quote \\
        Medianmiete & Verhältnis & 589 bis 3.501 USD je Monat & ACS B25064, bevölkerungsgewichtet \\
        Leerstandsquote & Verhältnis & 0,009 bis 0,237 & ACS B25002, Quote \\
        \addlinespace
        \multicolumn{4}{@{}l}{\textbf{Merkmale --- kriminalitätsbezogen}} \\
        Kriminalitätsindex, logarithmiert & Intervall & $-2{,}68$ bis 1,72 & SFPD, Standortquotient der Deliktrate je Einwohner über zwölf Vormonate \\
        \addlinespace
        \multicolumn{4}{@{}l}{\textbf{Merkmale --- baulich}} \\
        Altbauanteil vor 1940 & Verhältnis & 0,128 bis 0,908 & Land Use 2020, Anteil an den Parzellen mit bekanntem Baujahr \\
        Wohngebäudeanteil & Verhältnis & 0,105 bis 0,971 & Land Use 2020, Anteil an allen Parzellen \\
        Risikogewerbeanteil & Verhältnis & 0,000 bis 0,222 & Land Use 2020, Flächenanteil brandlastreicher Gewerbenutzung \\
        \addlinespace
        \multicolumn{4}{@{}l}{\textbf{Merkmale --- Größenkontrolle und Saison}} \\
        Wohnbevölkerung, logarithmiert & Intervall & 7,77 bis 11,32 & ACS B01003, Summe über die Tracts, logarithmiert \\
        Saison, Sinus & Intervall & $-1{,}000$ bis $1{,}000$ & abgeleitet, $\sin(2\pi \cdot \mathrm{Monat}/12)$ \\
        Saison, Kosinus & Intervall & $-1{,}000$ bis $1{,}000$ & abgeleitet, $\cos(2\pi \cdot \mathrm{Monat}/12)$ \\
        \addlinespace
        \multicolumn{4}{@{}l}{\textbf{Zielgrößen und Expositionsgröße}} \\
        Einsatzhäufigkeit & Verhältnis & 0 bis 451 Einsätze je Monat & SFFD, Zählung je Stadtteil und Monat \\
        Einsätze je 1.000 Einwohner & Verhältnis & 0,00 bis 16,99 & abgeleitet, auf die Wohnbevölkerung normiert \\
        überwiegende Einsatzart & nominal & Brand, Rettung, technische Hilfe, Fehlalarm & SFFD, größter der vier Anteile desselben Monats \\
        Wohnbevölkerung & Verhältnis & 2.357 bis 82.682 Personen & ACS B01003, Summe über die Tracts; Exposition \\
        \bottomrule
    \end{tabularx}
    \tabquelle
\end{table}

\subsection{Analyserahmen und Baselines}
\label{sec:rahmen-baselines}

Die Einsatzhäufigkeit unterscheidet sich weit stärker zwischen den Stadtteilen
als innerhalb eines Stadtteils über die Zeit; auf die räumliche Dimension
entfallen 92,5\,\% der Varianz (Abschnitt~\ref{sec:eda}). Ein Modell, das einen
Stadtteil bereits aus dem Training kennt, muss dessen Niveau deshalb nicht mehr
erklären. Die Aufteilung in Trainings- und Testmengen wird daher nicht entlang
der Zeit, sondern entlang der Stadtteile gezogen; Ziel der Validierung ist die
Generalisierung auf unbekannte Stadtteile. Sie entsteht in der Aufbereitung und
wird als Spalten \texttt{fold} und \texttt{ist\_holdout} in beide
Analysedateien geschrieben, ist also eine Eigenschaft des Datensatzes und nicht
der Algorithmen --- das Vorgehensmodell führt sie folgerichtig unter der
Auswahl der Daten.\footcite[23-26]{Chapman2000}
 
Geprüft wird, indem ganze Stadtteile zurückgehalten werden. Die 30
Entwicklungsstadtteile verteilen sich auf fünf Folds mit je sechs
Stadtteilen; jeder ist genau einmal Testfall, und zwar mit allen
132 Monaten (Abbildung~\ref{abb:foldstruktur}, Code~\ref{lst:test-folds}).
Ein Zeitschnitt prüfte die Forschungsfrage nicht: Dort steht jeder Stadtteil in
Training und Test, und das Modell kennt sein Niveau bereits --- der
Stadtteilmittelwert allein erreicht ein Bestimmtheitsmaß von 0,925
(Abschnitt~\ref{sec:eda}). Für genau diesen Fall empfehlen Roberts et al.\ eine
räumlich geblockte Aufteilung, wenn auf neue Orte vorhergesagt werden
soll.\footcite[925]{Roberts2017} Was der Rahmen dadurch nicht mehr misst, ist
die Fortschreibung eines bereits bekannten Stadtteils in die Zukunft.
 
\begin{figure}[H]
    \centering
    \includegraphics[width=\textwidth]{Abbildungen/a18_foldstruktur.pdf}
    \caption{Struktur der Aufteilung: Foldgröße, Wohnbevölkerung und brand-dominierte Monate je Gruppe}
    \label{abb:foldstruktur}
\end{figure}
 
\begin{quellcode}[htbp]
\setstretch{1}
\caption{Prüfung der Stadtteiltrennung an der erzeugten Datei}
\label{lst:test-folds}\vspace{-3.5pt}
\begin{lstlisting}[style=pythonstil,firstnumber=136]
def test_folds_ordnung_und_holdout():(*@\setcounter{lstnumber}{143}@*)
    d = regression()
    holdout = set(d.loc[d["ist_holdout"] == 1, "stadtteil"])
    gesehen = set()
    for k in range(1, N_FOLDS + 1):
        train, test = fold_masken(d, k)
        st_train = set(d.loc[train, "stadtteil"])
        st_test = set(d.loc[test, "stadtteil"])
        assert st_train.isdisjoint(st_test), \
            f"Fold {k}: Stadtteil in Training UND Test: {st_train & st_test}"
        assert st_test.isdisjoint(holdout), f"Fold {k} greift ins Hold-out"(*@\setcounter{lstnumber}{154}@*)
        assert gesehen.isdisjoint(st_test), "Stadtteil in zwei Folds im Test"
        gesehen |= st_test
    assert gesehen | holdout == set(d["stadtteil"]), \
        "Nicht jeder Stadtteil ist genau einmal Testfall oder im Hold-out"
\end{lstlisting}
\tabquelle[Eigene Darstellung, \texttt{tests/test\_aufbereitung.py}, Z.~136--158, gekürzt]
\end{quellcode}
 
Für die zweite Unterfrage war eine zeitreihengerechte Kreuzvalidierung und ein
Testfenster der letzten zwölf Monate vorgesehen. Der Wechsel  ist
die Folge eines Befundes über die Daten und keine Anpassung an ein Ergebnis.
Die Abwägung beider Rahmen steht in Kapitel~\ref{sec:diskussion}.
 
Welcher Stadtteil in welchen Fold fällt, ist nicht zufällig. Sortiert wird
zuerst nach der Zahl brand-dominierter Monate, bei Gleichstand nach der
Wohnbevölkerung; die so geordneten Stadtteile werden reihum auf sechs Gruppen
verteilt, von denen die erste das Hold-out bildet (Code~\ref{lst:folds}).
Das erste Kriterium ist nötig, weil die Hälfte der 70 brand-dominierten Monate
auf einen einzigen Stadtteil entfällt: Ohne Stratifizierung blieb ein Fold
regelmäßig ganz ohne Brand-Testfall, und das Macro-F1 mittelte über eine
Klasse, die dort nicht vorkam. Das zweite verhindert, dass die Streuung
zwischen den Folds ein Größeneffekt ist. Ein Leckageproblem entsteht nicht:
Festgelegt wird allein, welche Stadtteile gemeinsam getestet werden.
 
\begin{quellcode}[htbp]
\setstretch{1}
\caption{Fold-Zuteilung mit doppelter Stratifizierung}
\label{lst:folds}\vspace{-3.5pt}
\begin{lstlisting}[style=pythonstil,firstnumber=63]
def ergaenze_aufteilung(daten: pd.DataFrame, versatz: int = 0,
                        selten: pd.Series | None = None) -> pd.DataFrame:(*@\setcounter{lstnumber}{81}@*)
    bev = daten.groupby("stadtteil")[EXPOSURE_ROH].mean()
    if selten is None:
        selten = pd.Series(0, index=bev.index)
    ordnung = (pd.DataFrame({"selten": selten.reindex(bev.index).fillna(0),
                             "bev": bev})
                 .sort_values(["selten", "bev"], ascending=False).index)(*@\setcounter{lstnumber}{89}@*)
    gruppe = {st: (i + versatz) % (N_FOLDS + 1) for i, st in enumerate(ordnung)}
 
    d = daten.copy()
    d["fold"] = d["stadtteil"].map(gruppe).astype(int)
    d["ist_holdout"] = (d["fold"] == 0).astype(int)
    return d
\end{lstlisting}
\tabquelle[Eigene Darstellung, \texttt{prep/s2\_datensaetze.py}, Z.~63--95, gekürzt]
\end{quellcode}
 
Sechs Stadtteile bleiben während der gesamten Entwicklung unberührt und werden
erst für die Schlussbewertung geöffnet. Sie sind die erste Gruppe derselben
Ordnung und deshalb kein verkleinertes Abbild der Stadt: Auf sie entfallen 37
der 70 brand-dominierten Monate, ihr Brandanteil liegt bei 4,67\,\% gegenüber
0,83\,\% in den Entwicklungsstadtteilen. Ergebnisse auf dem Hold-out sind daher
nicht als Durchschnittsleistung über die Stadt zu lesen, sondern als Probe
unter erschwerten Bedingungen.
 
Die Zahl der Zeilen überzeichnet, wie viel unabhängige Information vorliegt.
Weil das Raster rechteckig ist, gilt die Näherung für gleich große Gruppen
genau: Aus einer Intraklassenkorrelation von 0,926 und 132 Monaten je
Stadtteil folgt ein Designeffekt von 122 und damit eine effektive Stichprobe
von rund 39 statt 4.752 Beobachtungen.\footcite[206]{Killip2004} Die
Aufteilung nach Stadtteilen trägt dem Rechnung, indem sie den Stadtteil und
nicht die Zeile als Einheit behandelt.
 
Dass alle Verfahren dieselben Merkmale und dieselbe Aufteilung sehen, ist keine
Zusicherung, sondern eine Eigenschaft des Aufbaus, und beides ist am Code
ablesbar. Der Merkmalssatz steht an einer einzigen Stelle: Die Aufbereitung führt die zehn Strukturmerkmale und die beiden Saisonterme als die Konstanten PRAEDIKTOREN und SAISON (Tabelle~\ref{tab:codebook}). Beide Modellskripte importieren genau diese und führen keine eigene Liste --- ein Merkmal kann so nicht in einem Strang stehen und im anderen fehlen. Die Fold-Zuteilung wiederum steht als
Spalte in beiden Analysedateien, und eine automatisierte Prüfung vergleicht sie
Stadtteil für Stadtteil zwischen Mengen- und Strukturdatensatz
(Code~\ref{lst:test-fairness}). Eine zweite Prüfung baut die Aufteilung
aus den Dateien nach und hält sie gegen die gespeicherten Spalten; ohne sie
liefen bei einer veralteten Aufteilung alle Verfahren auf falschen, aber
untereinander identischen Mengen: Der Vergleich bliebe fair und das Ergebnis
wäre trotzdem falsch.

\begin{quellcode}[htbp]
\setstretch{1}
\caption{Beide Datensätze tragen dieselbe Fold-Zuordnung}
\label{lst:test-fairness}\vspace{-3.5pt}
\begin{lstlisting}[style=pythonstil,firstnumber=290]
def test_struktur_gleiche_abgrenzung_wie_regression():(*@\setcounter{lstnumber}{296}@*)
    k, d = klassifikation(), regression()
    assert set(k["stadtteil"]) == set(d["stadtteil"])
    assert k["jahr_monat"].min() == d["jahr_monat"].min() == START
    assert k["jahr_monat"].max() == d["jahr_monat"].max() == ENDE(*@\setcounter{lstnumber}{302}@*)
    zuordnung = d.groupby("stadtteil")["fold"].first()
    assert (k.groupby("stadtteil")["fold"].first() == zuordnung).all(), \
        "Struktur- und Mengendatensatz nutzen verschiedene Folds"
\end{lstlisting}
\tabquelle[Eigene Darstellung, \texttt{tests/test\_aufbereitung.py}, Z.~290--305, gekürzt]
\end{quellcode}

% NOTIZ (07.09.2026) -- neu; ersetzt den frueheren Ausschnitt zur
% Datentypumwandlung. Nach drei einzelnen Pruefungen (Code auf Ergebnisvariablen,
% Stadtteiltrennung, gleiche Fold-Zuordnung) steht hier, wie sie ausgefuehrt
% werden.
% Erklaerbar im Text: (1) die Sammlung ueber die Namenskonvention in Zeile 399/400
% -- eine neue Pruefung wird gefunden, sobald sie test_ heisst, sie muss nirgends
% registriert werden; (2) Zeile 407 faengt nur AssertionError, kein bare except:
% ein Tippfehler im Test bricht den Lauf ab, statt als Fehlschlag durchzugehen;
% (3) der Exitcode in Zeile 415 macht den Lauf verkettbar -- build.py haengt ihn
% als optionalen letzten Schritt an. Daher die Angabe "20 von 20 bestanden".
\begin{quellcode}[htbp]
\setstretch{1}
\caption{Sammeln und Ausführen aller Prüfungen}
\label{lst:testlauf}\vspace{-3.5pt}
\begin{lstlisting}[style=pythonstil,firstnumber=398]
def main() -> int:
    tests = [(n, f) for n, f in sorted(globals().items())
             if n.startswith("test_") and callable(f)]
    print(f"Prüfungen der Datenaufbereitung ({len(tests)} Tests)\n")
    fehler = 0
    for name, funktion in tests:
        try:
            funktion()
            print(f"  [OK]     {name}")
        except AssertionError as e:
            fehler += 1
            print(f"  [FEHLER] {name}: {e}")
    print(f"\n{len(tests) - fehler}/{len(tests)} bestanden.")(*@\setcounter{lstnumber}{414}@*)
    return 1 if fehler else 0
\end{lstlisting}
\tabquelle[Eigene Darstellung, \texttt{tests/test\_aufbereitung.py}, Z.~398--415, gekürzt]
\end{quellcode}
 
Zuletzt werden beide Dateien auf modelltaugliche Typen gebracht: Merkmale als
\texttt{float64}, Schlüssel und Zählgrößen als \texttt{int64}, keine fehlenden
Werte, weil das Raster vollständig ist. Das ist
kein Formalismus --- eine einzige nullbare Ganzzahlspalte macht die
Merkmalsmatrix zu einem \texttt{object}-Feld, das XGBoost ablehnt und
scikit-learn stillschweigend umwandelt. Skaliert wird im Datensatz nicht; die
Standardisierung gehört zur Modellpipeline und wird je Fold nur auf den
Trainingsdaten angepasst.
 
Die Messlatte ist zweistufig angelegt (Abschnitt~\ref{sec:zielgroessen}).
Stufe~1 kommt ohne ein einziges Merkmal aus: der Gesamtmittelwert der
Trainingsstadtteile für die Mengenzielgrößen, die häufigste Klasse für die
Einsatzart. Sie beantwortet, ob in den Merkmalen überhaupt Information steckt.
Stufe~2 ist die einfachste Form, die zur Datenform passt: ein verallgemeinertes
lineares Modell mit kanonischem Link, unpenalisiert angepasst. Beide Stufen
laufen über dieselben 50 Läufe, dieselbe Aufteilung und dieselben zwölf
Merkmale wie die späteren Vergleichsverfahren; die Baseline ist damit
Mitbewerber unter identischem Protokoll und nicht bloß ein Erwartungswert.
Maßgeblich ist Stufe~2 --- ein Verfahren, das nur Stufe~1 schlägt, hat nichts
gezeigt (Tabelle~\ref{tab:baselines}). Gemittelt wird dabei zweistufig: erst über die fünf Folds einer Wiederholung, dann über die zehn Wiederholungen; die Streuung, die in dieser und in allen späteren Ergebnistabellen hinter dem Mittelwert steht, ist die zweite (Code~\ref{lst:mittelung}). Warum nicht über alle fünfzig Fold-Ergebnisse gestreut wird, begründet Abschnitt~\ref{sec:belastbarkeit}.
 
\begin{table}[htbp]
    \tabformat
    \caption{Referenzwerte beider Stufen, Mittel über 50 Läufe; Streuung über die zehn Wiederholungsmittel}
    \label{tab:baselines}
    \begin{tabularx}{\textwidth}{@{}L{3.9cm} c Y r r@{}}
        \toprule
        \textbf{Zielgröße} & \textbf{Stufe} & \textbf{Referenz} & \textbf{R\textsuperscript{2}} & \textbf{RMSE} \\
        \midrule
        Einsatzhäufigkeit & 1 & Gesamtmittelwert & $-0{,}865 \pm 0{,}778$ & $68{,}08 \pm 0{,}86$ \\
        Einsatzhäufigkeit & 2 & Poisson-Modell mit Offset & $0{,}607 \pm 0{,}070$ & $32{,}99 \pm 3{,}25$ \\
        Einsätze je 1.000 Einwohner & 1 & Gesamtmittelwert & $-0{,}653 \pm 0{,}332$ & $2{,}55 \pm 0{,}06$ \\
        Einsätze je 1.000 Einwohner & 2 & Poisson-Modell mit Offset & $0{,}306 \pm 0{,}110$ & $1{,}74 \pm 0{,}15$ \\
        \midrule
        & & & \textbf{Macro-F1} & \textbf{Trefferquote} \\
        \cmidrule(l){4-5}
        Überwiegende Einsatzart & 1 & Häufigste Klasse & 0,221 & 0,795 \\
        Überwiegende Einsatzart & 2 & Multinomiales Logit & $0{,}301 \pm 0{,}015$ & 0,590 \\
        \bottomrule
    \end{tabularx}
    \tabquelle
\end{table}

% NOTIZ (07.09.2026) -- neu. Die Beschriftung der Tabelle darueber sagt "Mittel
% ueber 50 Laeufe; Streuung ueber die zehn Wiederholungsmittel"; woher der
% Unterschied kommt, stand bisher nirgends als Konstruktion, sondern erst in
% Abschnitt 7.3 als Argument.
% Erklaerbar im Text (mit Zeilenbezug):
% (1) Zeile 531 bildet die Streuung ueber alle 50 Fold-Ergebnisse, Zeile 532/533
%     erst das Mittel je Wiederholung und dann die Streuung darueber -- zwei
%     Zahlen aus denselben Daten, die sich um ein Vielfaches unterscheiden;
% (2) berichtet wird die zweite, weil die 50 Fold-Ergebnisse dieselben 30
%     Stadtteile in zehn Gruppierungen sind und ein Intervall aus
%     std_folds/sqrt(50) zu eng waere;
% (3) beide Spalten wandern bewusst mit in die Ergebnisdatei -- der Unterschied
%     ist selbst ein Befund und nicht nur ein Zwischenschritt.
\begin{quellcode}[htbp]
\setstretch{1}
\caption{Zweistufiges Mitteln: erst je Wiederholung, dann darüber}
\label{lst:mittelung}\vspace{-3.5pt}
\begin{lstlisting}[style=pythonstil,firstnumber=517]
def aggregiere(folds: pd.DataFrame) -> pd.DataFrame:(*@\setcounter{lstnumber}{527}@*)
    schluessel = ["zielgroesse", "verfahren"]
    g = folds.groupby(schluessel, sort=False)
    z = g[MASSE].mean().add_suffix("_mean")
    z = z.join(g[MASSE].std().add_suffix("_std_folds"))
    je_wdh = folds.groupby(schluessel + ["wiederholung"], sort=False)[MASSE].mean()
    z = z.join(je_wdh.groupby(schluessel, sort=False).std()
                     .add_suffix("_std_wiederholungen"))
\end{lstlisting}
\tabquelle[Eigene Darstellung, \texttt{modelle/m02\_menge.py}, Z.~517--534, gekürzt]
\end{quellcode}
 
Für die Mengenzielgrößen ist Stufe~2 ein Poisson-Modell mit Offset: Die
logarithmierte Wohnbevölkerung geht mit fest auf eins gesetztem Koeffizienten
ein, geschätzt wird damit die Rate (Code~\ref{lst:poisson}). Über den
Logarithmus als Link wirkt das Modell auf der Skala der Zielgröße
multiplikativ; es erfasst Krümmung, aber keine Wechselwirkungen --- anders als
die beiden Baumverfahren (Abschnitt~\ref{sec:referenzmodelle}). Einen frei wählbaren
Parameter hat es nicht. Die zweite Mengenzielgröße entsteht aus derselben
Anpassung, geteilt durch die Wohnbevölkerung; ein eigenes Ratenmodell wäre eine
zweite Spezifikation und damit nicht mehr dasselbe Protokoll. Für die
Einsatzart ist Stufe~2 das Gegenstück: ein multinomiales Logit, ebenfalls ohne
Strafterm und mit ausgeglichenen Klassengewichten (Code~\ref{lst:logit}).
 
\begin{quellcode}[htbp]
\setstretch{1}
\caption{Poisson-Modell mit Offset, Stufe 2 der Menge}
\label{lst:poisson}\vspace{-3.5pt}
\begin{lstlisting}[style=pythonstil,firstnumber=74]
def poisson_glm(train: pd.DataFrame, test: pd.DataFrame,
                merkmale: list[str] | None = None) -> np.ndarray:(*@\setcounter{lstnumber}{92}@*)
    import statsmodels.api as sm
 
    spalten = MERKMALE if merkmale is None else list(merkmale)
    X_tr = sm.add_constant(train[spalten].astype(float), has_constant="add")
    X_te = sm.add_constant(test[spalten].astype(float),  has_constant="add")
    y_tr = train[ZIELGROESSE].astype(float)(*@\setcounter{lstnumber}{101}@*)
    off_tr = np.log(train[EXPOSURE_ROH].astype(float))
    off_te = np.log(test[EXPOSURE_ROH].astype(float))(*@\setcounter{lstnumber}{105}@*)
    modell = sm.GLM(y_tr, X_tr, family=sm.families.Poisson(),
                    offset=off_tr).fit()
    return np.asarray(modell.predict(X_te, offset=off_te))
\end{lstlisting}
\tabquelle[Eigene Darstellung, \texttt{vorpruefung/v1\_baselines.py}, Z.~74--108, gekürzt]
\end{quellcode}
 
\begin{quellcode}[htbp]
\setstretch{1}
\caption{Multinomiales Logit, Stufe 2 der Struktur}
\label{lst:logit}\vspace{-3.5pt}
\begin{lstlisting}[style=pythonstil,firstnumber=111]
def logit_glm(train: pd.DataFrame, merkmale: list[str] | None = None):(*@\setcounter{lstnumber}{126}@*)
    from sklearn.linear_model import LogisticRegression
    from sklearn.pipeline import make_pipeline
    from sklearn.preprocessing import StandardScaler
 
    spalten = MERKMALE if merkmale is None else list(merkmale)(*@\setcounter{lstnumber}{134}@*)
    return make_pipeline(
        StandardScaler(),
        LogisticRegression(max_iter=2000, C=np.inf, class_weight="balanced")
    ).fit(train[spalten].astype(float), train[ZIELKLASSE])
\end{lstlisting}
\tabquelle[Eigene Darstellung, \texttt{vorpruefung/v1\_baselines.py}, Z.~111--138, gekürzt]
\end{quellcode}

% NOTIZ (07.09.2026) -- neu. Die beiden Ausschnitte darueber zeigen, WAS die
% Stufe-2-Baselines sind; dieser zeigt, wie sie laufen -- und dass sie unter
% demselben Protokoll laufen wie die Vergleichsverfahren.
% Erklaerbar im Text: (1) die beiden Schleifen in Zeile 154/156 sind dieselben
% zehn Wiederholungen und fuenf Folds wie spaeter in Kapitel 6, also 50 Laeufe;
% (2) Zeile 161 ruft das Poisson-Modell EINMAL, und Zeile 163/164 leiten die Rate
% aus derselben Vorhersage ab (geteilt durch die Bevoelkerung) -- ein eigenes
% Ratenmodell waere eine zweite Spezifikation und die beiden Zielgroessen waeren
% nicht mehr vergleichbar; (3) Zeile 166-168 stellt Stufe 1 und Stufe 2
% nebeneinander in dieselbe Ergebniszeile, weshalb der Abstand zwischen ihnen
% ohne weitere Rechnung ablesbar ist.
\begin{quellcode}[htbp]
\setstretch{1}
\caption{Beide Baseline-Stufen über dieselben 50 Läufe}
\label{lst:baselinelauf}\vspace{-3.5pt}
\begin{lstlisting}[style=pythonstil,firstnumber=141]
def regression(panel: pd.DataFrame, selten: pd.Series) -> pd.DataFrame:(*@\setcounter{lstnumber}{152}@*)
    zeilen = []
    for w in range(WIEDERHOLUNGEN):
        d = wiederholte_aufteilung(panel, wiederholung=w, selten=selten)
        for k in range(1, N_FOLDS + 1):
            tr, te = fold_masken(d, k)
            train, test = d[tr], d[te]
            bev = test[EXPOSURE_ROH].to_numpy()

            anzahl = poisson_glm(train, test)

            for ziel, referenz in ((ZIELGROESSE, anzahl),
                                   (RATE, anzahl / bev * 1000)):
                y = test[ziel].to_numpy()
                for stufe, name, y_hat in (
                        (1, NULLMARKE, np.full(len(test), train[ziel].mean())),
                        (2, POISSON, referenz)):
                    zeilen.append({"wiederholung": w, "fold": k, "stufe": stufe,
                                   "zielgroesse": ziel, "modell": name,
                                   **bewerte_regression(y, y_hat)})
\end{lstlisting}
\tabquelle[Eigene Darstellung, \texttt{vorpruefung/v1\_baselines.py}, Z.~141--171, gekürzt]
\end{quellcode}
 
Die Zähldaten sind deutlich überdispers (Dispersionsindex 61,8;
Abschnitt~\ref{sec:eda}). Die naheliegende Erweiterung, eine negative
Binomialverteilung, war ursprünglich vorgesehen. Sie löst jedoch ein
Inferenzproblem, das diese Baseline nicht hat
(Abschnitt~\ref{sec:referenzmodelle}). Die Reduktion auf die einfacheren
Varianten ist damit methodisch sauber, vermeidet willkürliche Parameter und
liefert stärkere Vergleichswerte. Tragend ist das zweite Argument: Für einen
Strafterm gäbe es zwei Wege, beide begründungsbedürftig. Der Vorgabewert einer
Softwarebibliothek ist keine für diese Daten getroffene Wahl, Vorgabewerte
entstehen häufig ad hoc oder heuristisch\footcite[4]{Probst2019}, und ein
optimierter Wert verlangte ein Auswahlverfahren je Trainingsfold. Die
unpenalisierte Anpassung hat dagegen keinen zu wählenden Parameter. Das dritte
Argument gilt nicht durchgehend: Im Mengenstrang ist die unpenalisierte Form
die härtere Messlatte, im Strukturstrang die weichere. Dass dieselbe Festlegung
einen Strang härter, den anderen weicher macht, ist zugleich das Argument
dagegen, sie nach dem Ergebnis auszuwählen.

Drei Eigenschaften dieser Referenzwerte werden leicht als Fehler gelesen. Ein
negatives Bestimmtheitsmaß ist keiner, sondern Folge des Rahmens: Gemessen wird
gegen den Mittelwert der jeweiligen Testdaten; wer für einen unbekannten
Stadtteil den Gesamtdurchschnitt vorhersagt, liegt schlechter als dessen
Mittelwert, genau diese Lücke sollen die Strukturmerkmale schließen. Das Bestimmtheitsmaß ist auf der Rate kein Hauptmaß, weil deren
Niveau zwischen den Stadtteilen um ein Vielfaches auseinanderliegt; maßgeblich
sind dort RMSE und mittlerer absoluter Fehler. In der Klassifikation liegt die
Trefferquote der Stufe~2 unter der von Stufe~1, während das Macro-F1 steigt;
maßgeblich ist Macro-F1 (Abschnitt~\ref{sec:modellbewertung}). Eine
Fortschreibung des Vormonatswerts fehlt bewusst: Sie verwendete die eigene
Vergangenheit des Teststadtteils, die im Training fehlt.

 
% \begin{table}[htbp]
%     \centering
%     \caption{Steckbrief der beiden Analysedatensätze}
%     \label{tab:steckbrief}
%     \footnotesize
%     \singlespacing
%     \setlength{\tabcolsep}{5pt}
%     \begin{tabularx}{\textwidth}{@{}>{\raggedright\arraybackslash}p{4.6cm} >{\raggedright\arraybackslash}X@{}}
%         \toprule
%         \textbf{Eigenschaft} & \textbf{Ausprägung} \\
%         \midrule
%         Analyseeinheit & Stadtteil und Monat, in beiden Dateien \\
%         Zeitraum & Januar 2015 bis Dezember 2025, 132 Monate \\
%         Beobachtungen Menge & 4.752, das sind 36 mal 132 \\
%         Beobachtungen Struktur & 4.751, ein Monat ohne Einsatz entfällt \\
%         Merkmale & 10 Strukturmerkmale und 2 Saisonterme \\
%         Aufteilung & 5 Folds mit je 6 Stadtteilen; 30 Entwicklungs- und 6 Hold-out-Stadtteile \\
%         Ausschlüsse & 3 Parkgebiete ohne Wohnbevölkerung, 2 Stadtteile ohne Baujahresangabe \\
%         Absicherung & 19 automatisierte Prüfungen an den erzeugten Dateien \\
%         \bottomrule
%     \end{tabularx}
% \end{table}
% ==========================================================================
% AUFLAGEN AUS SCHROETERS FREIGABE VOM 04.08.2026 (#35, 01_VORGABEN.md §0)
% Alle drei sind am 25.08.2026 noch OFFEN. Dieser Abschnitt ist ihr Ort.
% Quelle: E-Mail 04.08.2026, zitiert in 02_ENTSCHEIDUNGEN.md #35
% ==========================================================================
%
% AUFLAGE A - Die Abweichung transparent erlaeutern
%   Unterfrage 2 lautete im angemeldeten Expose "mittels zeitreihengerechter
%   Kreuzvalidierung"; im Text steht jetzt "stadtteilweiser". Diese Aenderung
%   ist NIRGENDS erlaeutert - das Wort "Abweichung" kommt in main.tex nicht vor.
%   Schroeter: das begruendete Verwerfen in Kapitel 8 ist VERLANGT, nicht
%   optional. Hier ein Verweis, die Ausfuehrung in 8.3.
%   Vorgeschichte, die dazugehoert: #14 (26.07.2026) sah ein End-Hold-out der
%   letzten 12 Monate vor - der Zeitschnitt war der urspruengliche Plan und
%   wurde am 28.07. durch #29 ersetzt. Das ist ein verworfener Weg mit Datum,
%   kein nachtraeglicher Schwenk.
%
% AUFLAGE B - Woertlich "Generalisierung auf unbekannte Stadtteile"
%   Schroeter hat die Formulierung selbst verlangt; sie gehoert in die
%   Zielsetzung (1.2) UND hierher. ACHTUNG, R-17: In der E-Mail vom
%   08.08.2026 wurde ihm gemeldet, sie sei bereits formuliert. Sie ist es
%   nicht - sie steht bislang nur in Kommentaren (Z. 437 und Z. 1629).
%   Damit ist aus der Auflage eine ZUSAGE geworden. Er wird beim Lesen
%   darauf achten.
%
% AUFLAGE C - Identische Merkmale und Splits im Text BELEGEN
%   Nicht behaupten, sondern den Mechanismus nennen:
%     - Die Fold-Zuteilung entsteht in EINEM Aufruf und wird als Spalten
%       `fold` und `ist_holdout` in BEIDE Parquet-Dateien geschrieben (#22).
%       Mengen- und Strukturstrang koennen konstruktiv keine
%       unterschiedlichen Aufteilungen sehen.
%     - Die Merkmalsliste steht genau einmal, in prep/config.py. Es gibt
%       keine zweite Stelle, an der eine Merkmalsliste definiert ist.
%     - Kein Modellskript kann davon abweichen.
%   Das ist die Antwort auf "Wie stellen Sie sicher, dass der Vergleich fair
%   ist?" - und derzeit das staerkste Argument des Designs, das in der
%   Arbeit nicht vorkommt.
%
% WEITERE PFLICHTINHALTE DIESES ABSCHNITTS
%   - Auflage D woertlich (#47): "methodisch sauber, vermeidet willkuerliche
%     Parameter, liefert zudem staerkere Vergleichswerte" - mit dem
%     Vorbehalt, dass das dritte Argument NUR fuer den Mengenstrang gilt:
%     Regressionslatte 37,27 -> 33,98 RMSE, Klassifikationslatte SINKT
%     0,314 -> 0,297.
%   - Zusammensetzung des Hold-outs (results/panelprofil/): sechs
%     Stadtteile, 37 der 70 brand-dominierten Monate, brand-Anteil 4,67 %
%     gegen 0,86 % in der Entwicklung. Rang 1 der Zuteilungsordnung faellt
%     nach der Regel aus #30 zwangslaeufig ins Hold-out.
%   - Effektive Stichprobe: ICC 0,926, Designeffekt 122, n_eff etwa 38.
%     Kommt derzeit im GESAMTEN Text nicht vor.
% ==========================================================================

% INHALT (Schwerpunkt, Code-Listing 3):
% - Warum der Split hierher gehoert und nicht nach Kapitel 6: Die Aufteilung
%   steht als Spalten fold und ist_holdout IN DEN DATEIEN. Sie ist eine
%   Eigenschaft des Datensatzes, nicht der Algorithmen -- damit ist die
%   Fairness-Regel konstruktiv abgesichert statt behauptet.
% - STADTTEIL-Split statt Zeitschnitt: 29 Entwicklungsstadtteile auf 5 Folds
%   (6 | 6 | 6 | 6 | 5), dazu 6 Hold-out-Stadtteile; jeder Stadtteil genau
%   einmal Testfall, mit allen 132 Monaten. Trainingszeilen je Fold 3.036 bis
%   3.168, Testzeilen 660 bis 792.
%   Begruendung (zentrale methodische Aussage der Arbeit): Ein Zeitschnitt
%   prueft die Forschungsfrage NICHT, weil dort jeder Stadtteil in Training
%   und Test steht und das Modell sein Niveau bereits kennt.
%   Beleg: 92,5 % der Varianz liegen zwischen den Stadtteilen; der
%   Stadtteil-Mittelwert allein erreicht dort R2 0,888.
% - Doppelte Stratifizierung: sortiert nach brand-dominierten Monaten, bei
%   Gleichstand nach Bevoelkerung. Von 70 brand-dominierten Monaten liegen 35
%   allein in Bayview Hunters Point; ohne Stratifizierung hatte in drei von
%   vier Aufteilungen ein Fold NULL Brand-Testfaelle. Jetzt 13, 9, 6, 3, 2.
%   Kein Leakage -- vgl. StratifiedGroupKFold.  -> LISTING 3: Fold-Zuteilung
% - Modelltauglichkeit: alle Merkmale float64, keine fehlenden Werte. Die
%   Int64-Falle erwaehnen (eine nullable Spalte macht X.to_numpy() zum
%   object-Array; sklearn faengt das still ab, XGBoost lehnt es ab).
% - Baselines (Auflage Schroeter 27.07.2026: gehoeren in die Data Preparation).
%   ZWEI STUFEN, #32/#33, neu gefasst mit #45. Stand 03_STAND.md Abschnitt 4,
%   ueber 50 Laeufe; Streuung ist std_wiederholungen ueber die 10
%   Wiederholungsmittel, NICHT std_folds:
%
%     Zielgroesse            Stufe  Baseline               R2               RMSE
%     anzahl_einsaetze         2    Poisson-GLM/Offset     0,542 +/- 0,082  33,98 +/- 3,11
%     anzahl_einsaetze         1    Gesamtmittelwert      -0,744 +/- 0,325  69,93 +/- 1,92
%     einsaetze_je_1000_ew     2    Poisson-GLM/Offset     0,367 +/- 0,261   4,08 +/- 0,62
%     einsaetze_je_1000_ew     1    Gesamtmittelwert      -1,054 +/- 0,875   7,54 +/- 0,13
%
%     dominante_einsatzart     1    Mehrheitsklasse    Macro-F1 0,223 | Acc. 0,806
%     dominante_einsatzart     2    Multinom. Logit    Macro-F1 0,297 +/- 0,014
%                                                      | AUROC 0,705 | Acc. 0,584
%
%   DIE RATE ENTSTEHT AUS DERSELBEN ANPASSUNG, geteilt durch die Bevoelkerung.
%   Ein zweites Modell waere eine zweite Spezifikation und damit unfair
%   gegenueber den Vergleichsverfahren (#43).
%
%   AUF DER RATE IST R2 KEIN HAUPTMASS. In Fold 4 faellt es auf -0,920, obwohl
%   das Poisson-GLM die Nullmarke bei RMSE in JEDEM Fold schlaegt. Ursache: R2
%   misst gegen den Mittelwert der TESTdaten, und die Rate streut zwischen den
%   Stadtteilen um Faktor 32 (Excelsior 1,04 gegen Financial District 33,80).
%   Konsequenz fuer Kapitel 7: bei der Rate RMSE/MAE berichten, R2 nur
%   nachrichtlich -- mit genau dieser Begruendung.
%
%   DER VERGLEICH DER BEIDEN KLASSIFIKATIONSZEILEN IST SELBST EIN ARGUMENT:
%   Das Logit hat die deutlich SCHLECHTERE Trefferquote (0,584 gegen 0,806) und
%   zugleich das deutlich BESSERE Macro-F1. Kein Widerspruch, sondern die
%   Wirkung von class_weight="balanced" -- das Modell gibt Treffer bei der
%   dominanten Klasse auf, um die drei seltenen ueberhaupt zu finden. Wer
%   Accuracy als Hauptmass naehme, kaeme zum Schluss, das Modell sei
%   schlechter geworden. Genau deshalb ist Macro-F1 massgeblich.
%
%   Zwei Punkte muessen in den Text:
%     (a) Negative R2 sind korrekt und aussagekraeftig -- wer fuer einen
%         unbekannten Stadtteil den Gesamtdurchschnitt vorhersagt, liegt
%         schlechter als dessen eigener Mittelwert. Genau diese Luecke sollen
%         die Strukturmerkmale schliessen.
%     (b) Die naive Vormonats-Baseline entfaellt: Sie wuerde die eigene
%         Vergangenheit des Teststadtteils nutzen.
%   Zur Auflage "nichtlineare Baseline" (Schroeter 27.07.2026) -- HIER IST
%   EHRLICHKEIT NOETIG: Die frueher gegebene Begruendung (Vormonatswert und
%   saisonaler Durchschnitt seien nichtparametrisch) traegt nicht mehr. Der
%   Vormonatswert ist mit dem Stadtteil-Split entfallen, siehe (b) oben.
%   Gesamtmittelwert und saisonaler Durchschnitt benutzen KEIN EINZIGES
%   MERKMAL und fielen nachgerechnet praktisch zusammen -- faktisch eine
%   triviale Referenz, nicht zwei.
%
%   GEWAEHLT WURDE STATTDESSEN ein verallgemeinertes lineares Modell mit
%   kanonischem Link als Stufe-2-Referenz: fuer die Menge ein POISSON-GLM MIT
%   OFFSET, fuer die Struktur ein MULTINOMIALES LOGIT OHNE STRAFTERM
%   (#32/#33, neu gefasst mit #45, freigegeben 08.08.2026). Fuenf Argumente,
%   ausformuliert in 01_VORGABEN.md:
%     1. Kein Strohmann -- dieselben zwoelf Merkmale, dieselben Zeilen,
%        dieselben Folds wie die Vergleichsverfahren.
%     2. Auf der Originalskala nichtlinear -- ueber den Log-Link multiplikativ,
%        bildet also ueberproportionale Effekte ab. Damit ist die Auflage
%        "keine lineare Baseline, wenn es keine Linearitaet gibt" beantwortet.
%     3. Verteilungsgerecht fuer Zaehldaten mit Exposition -- die Bevoelkerung
%        geht als Offset ein. Die Ueberdispersion (Dispersionsindex 62,8)
%        beschaedigt beim Poisson-Schaetzer nur die Standardfehler, nicht die
%        Konsistenz des bedingten Mittelwerts (Gourieroux, Monfort & Trognon
%        1984) -- und Standardfehler verwendet eine Baseline mit reinen
%        Punktvorhersagen nicht.
%     4. Sie zieht genau die relevante Trennlinie: Kruemmung ja, WECHSEL-
%        WIRKUNGEN nein -- und genau die finden RF und XGBoost konstruktions-
%        bedingt. Schlagen sie die Latte, ist der Mehraufwand belegt; schlagen
%        sie sie nicht, reicht die einfachere Struktur. Beides ist ein Ergebnis.
%     5. Kein frei waehlbarer Parameter -- unpenalisierte Maximum-Likelihood.
%
%   WARUM POISSON UND NICHT NEGATIVE BINOMIAL (#45): Die Negative Binomial
%   waere die Erweiterung fuer korrekte INFERENZ. Sie loest ein Problem, das
%   diese Baseline nicht hat, und bringt mit dem Dispersionsparameter eine
%   zusaetzliche Groesse mit -- sie ist damit nicht mehr "die einfachste Form,
%   die zur Datenform passt". Gemessen ist das Poisson-GLM zudem die HAERTERE
%   Latte: 33,98 gegen 37,27 RMSE.
%   ZU BENENNEN: Schroeter hatte die Negative Binomial namentlich freigegeben;
%   der Wechsel wurde ihm am 08.08.2026 mitgeteilt (R-14).
%
%   AUFLAGE D (Schroeter, 08.08.2026) -- WOERTLICH ZU DOKUMENTIEREN. Er hat
%   die Begruendung selbst formuliert und angewiesen, sie "genau so" zu
%   dokumentieren. Seine drei Argumente, in seiner Reihenfolge:
%     1. Die Reduktion auf die einfacheren Varianten ist methodisch sauber.
%     2. Sie vermeidet willkuerliche Parameter.
%     3. Sie liefert zudem staerkere Vergleichswerte.
%   Punkt 2 ist der tragende: Ein Strafterm auf dem Vorgabewert der Software
%   waere ein willkuerlich gesetzter Parameter, ein getunter eine zu
%   begruendende Wahl. Die unpenalisierte Anpassung hat nichts zu waehlen.
%   DER VORBEHALT GEHOERT DAZU -- Punkt 3 trifft auf den Mengenstrang zu (die
%   Latte STEIGT von 37,27 auf 33,98 RMSE), auf den Strukturstrang jedoch
%   NICHT (dort SINKT sie von 0,314 auf 0,297 Macro-F1). Beide Zahlen standen
%   in der Anfrage vom 08.08., die Freigabe erfolgte also in Kenntnis. Punkt 3
%   nicht pauschal uebernehmen, sondern den Unterschied benennen (R-2).
%   Das ist zugleich das Argument gegen Rosinenpicken: dieselbe Aenderung
%   hilft im einen Strang und schadet im anderen.
%
%   DIE KLASSIFIKATION HAT SEIT #33 EIN PENDANT. Fruehere Fassungen dieses
%   Abschnitts hielten fest, dort bleibe nur die Mehrheitsklasse -- das gilt
%   NICHT mehr. Das multinomiale Logit ist das Gegenstueck zum Poisson-GLM:
%   dieselbe Modellklasse, derselbe kanonische Link, derselbe Verzicht auf
%   einen Strafterm. Zu benennen bleibt, dass die Latte dort trotzdem
%   niedriger liegt als in der Regression (R-2) -- und dass ein flacher
%   Entscheidungsbaum (Macro-F1 0,270) das Logit NICHT schlaegt, der
%   Mehraufwand von RF und XGBoost im zweiten Strang vorab also nicht
%   begruendet war.
%
%   AUFLAGE SCHROETER 04.08.2026 (R-17): Die Zielsetzung der Validierung ist
%   hier WOERTLICH als "Generalisierung auf unbekannte Stadtteile" zu
%   formulieren. Steht noch nicht im Text.
% - Steckbrief des finalen Datensatzes als Tabelle:
%     Analyseeinheit           Stadtteil x Monat, beide Datensaetze
%     Zeitraum                 2015-01 bis 2025-12 (132 Monate)
%     Beobachtungen Menge      4.620 (35 x 132)
%     Beobachtungen Struktur   4.619 (ein Monat ohne Einsatz)
%     Merkmale                 10 Struktur + 2 Saison = 12
%     Aufteilung               5 Folds (6 | 6 | 6 | 6 | 5 Stadtteile),
%                              29 Entwicklungsstadtteile + 6 Hold-out
%     Ausschluesse             3 Parkgebiete, 3 ohne durchgaengige ACS-Abdeckung
%     Absicherung              19 automatisierte Pruefungen an den Dateien

% ===================================================
\section{Modelling}
% Umfang: ca. 8 Seiten -- CRISP-DM Phase 4
% ABGRENZUNG (Block vor Kapitel 5): Hier steht, WELCHE Verfahren auf den
% fertigen Datensatz angewandt werden und WARUM -- nicht, wie sie
% funktionieren (das steht in 2.3) und nicht, wie gut sie waren (Kapitel 7).
% Alles, was je Fold in der Pipeline laeuft, gehoert hierher: Skalierung,
% log(1+y), Klassengewichte, Label-Encoding, Tuning. Die Eignungspruefung in
% 6.2 prueft VERFAHREN gegen Daten, nicht die Daten selbst -- der Unterschied
% zu 4.2 ist im Abgrenzungsblock unter (a) beschrieben.
\subsection{Versuchsaufbau und technische Umsetzung}
% Der VALIDIERUNGSRAHMEN steht jetzt in Abschnitt 5.4 -- die Aufteilung ist
% eine Eigenschaft des Datensatzes (Spalten fold, ist_holdout) und wird hier
% nur noch referenziert, nicht erneut hergeleitet.
% Hier: identischer Datensatz fuer alle Verfahren, Seeds, wiederholte Splits
% (10 Wiederholungen), Python, Bibliotheken, Versionen. Code in den Anhang.

% ACHTUNG: Der folgende Absatz stammt aus dem Expose und nennt
% Time-Series-Cross-Validation. Das gilt seit dem Umbau NICHT mehr -- der
% Hauptvergleich laeuft ueber einen STADTTEIL-Split (siehe 5.4). Der Absatz
% ist beim Ausformulieren entsprechend zu ersetzen; die Abweichung vom Expose
% gehoert in die Diskussion (8.3) und in die Sprechstunde.

% ---- ÜBERNAHME AUS EXPOSÉ (Kap. 3, Abs. 2 letzter Satz + Abs. 4) -- 1:1, Tempus angepasst ----

% ERSETZT AM 26.08.2026 durch den ausformulierten Abschnitt darunter. Die
% beiden Saetze stammen 1:1 aus dem Expose (Kap. 3, Abs. 2 und 4); sie sind
% inhaltlich in Absatz 1 (eine Spezifikation fuer alle Verfahren) und in
% Tabelle tab:techstack (Python) aufgegangen. Auskommentiert statt geloescht,
% damit der Abgleich mit dem Expose moeglich bleibt:
% Alle Modelle erhalten denselben aufbereiteten Datensatz, um Performance-Unterschiede ausschließlich
% auf die Verfahren zurückführen zu können. Die technische Umsetzung erfolgt in Python.

% ENTFERNT AM 18.08.2026: Der Satz "Die Evaluation erfolgt zeitbewusst mittels
% Time-Series-Cross-Validation, die bei zeitlich strukturierten Daten
% Datenleckagen aus der Zukunft vermeidet.\footcite[192-193]{Bergmeir2012}"
% stand hier aus dem Expose und ist seit #29 SACHLICH FALSCH -- der
% Hauptvergleich laeuft ueber einen Stadtteil-Split. Bergmeir & Benitez bleibt
% als Beleg brauchbar, aber an anderer Stelle: in 5.0 fuer den Grundsatz, dass
% jeder Praediktor zum Prognosezeitpunkt verfuegbar sein muss (dort bereits
% zitiert), und in 2.6.3 fuer die allgemeine Darstellung zeitbewusster
% Validierung. Die ABWEICHUNG vom Expose gehoert nach 8.3.
%
% ERGAENZEN, konkret (Quellen: 03_STAND.md Abschnitte 3 und 5.0/5.4):
%   - Validierungsrahmen wird aus 5.4 nur REFERENZIERT, nicht neu hergeleitet:
%     Stadtteil-Split, 10 Wiederholungen x 5 Folds, Tuning einmal auf
%     Wiederholung 0, Hold-out einmalig nach Abschluss von Modellwahl und Tuning.
%   - Die eine Spezifikation fuer alle Modelle (#43): alle Verfahren
%     modellieren die RATE, fuer die absolute Zahl wird mit der Einwohnerzahl
%     zurueckmultipliziert. Verlustfunktionen: Ridge log(1+y) mit
%     Ruecktransformation ueber expm1, Random Forest criterion="poisson",
%     XGBoost reg:tweedie mit getuntem Varianzexponenten (#42).
%   - Nebeneffekt, der hierher gehoert: KEINE negativen Vorhersagen in keinem
%     der 300 Laeufe -- Tweedie und Poisson haben eine Log-Verknuepfung, die
%     Zielgroesse wird strukturell respektiert statt nachtraeglich geprueft.
%   - REPRODUZIERBARKEIT: RANDOM_STATE = 42 UND die KERNZAHL nennen. XGBoost
%     ist nicht threaddeterministisch (B-24); alle berichteten Guetemasse
%     stammen aus dem EINKERNIGEN Fit. Ohne die Kernzahl ist die Angabe
%     unvollstaendig.
%   - Messumgebung und Bibliotheksversionen: Intel Core i5-7300U @ 2,60 GHz,
%     2 physische / 4 logische Kerne, 7,8 GB RAM, Windows 10 Pro, Python 3.14.0,
%     Pakete in requirements_lauf.txt. Gehoert hierher, weil erst damit die
%     Laufzeiten in 7.3 lesbar sind.
%   - Begruendung fuer die einkernige Messung (#39/#40) -- der zwischen den
%     Verfahren vergleichbare Wert; der Parallelisierungsgewinn wird getrennt
%     ausgewiesen.
% ---- ENDE ÜBERNAHME ----

Der Vergleich läuft in zwei Strängen: auf der Einsatzhäufigkeit mit drei
Verfahren, auf der Einsatzart mit zweien. Innerhalb jedes Strangs sehen alle
Verfahren dieselben zwölf Merkmale, dieselben Zeilen und dieselbe Aufteilung in
Folds --- alle drei stammen aus der Aufbereitung und stehen als Spalten in den
Analysedateien (Abschnitt~\ref{sec:rahmen-baselines}). Zwischen zwei Läufen
unterscheidet sich damit genau eine Sache, nämlich das Verfahren; ohne diese
Gleichheit wäre jede Differenz auch als Folge einer anderen Spezifikation
lesbar.

Im Mengenstrang kommt eine vierte Gleichheit hinzu, und sie wird hier
festgelegt: Alle drei Verfahren schätzen dieselbe Größe wie die Referenz der
Stufe~2, nämlich die Einsätze je 1.000 Einwohner, und rechnen erst in der
Vorhersage mit der Wohnbevölkerung des jeweiligen Stadtteil-Monats auf die
absolute Zahl zurück. Angepasst wird damit auf der Rate und berichtet die
Anzahl --- ein Modell mit zwei Berichtsskalen und nicht zwei Modelle (Code~\ref{lst:exposition}). Was dadurch nicht gilt: Die Ergebnisse der beiden Mengenzielgrößen sind keine zwei unabhängigen Befunde, denn sie stammen aus derselben Anpassung. Was diese Gleichheit nicht herstellt, ist die Gewichtung der Anpassung: Beim Poisson-Modell geht die Exposition über den Offset zusätzlich in die Likelihood ein, sodass dort große Stadtteile die Schätzung stärker bestimmen als bei den drei Vergleichsverfahren. Beide Gewichtungen liegen jedoch als Auswertung vor, weil die Zählskala rechnerisch die mit dem Quadrat der Exposition gewichtete Ratenskala ist --- und die Rangfolge ist zwischen ihnen invariant (Abschnitt~\ref{sec:ergebnisse-menge}).

\begin{quellcode}[htbp]
\setstretch{1}
\caption{Geschätzt wird die Rate, berichtet wird die Anzahl}
\label{lst:exposition}\vspace{-3.5pt}
\begin{lstlisting}[style=pythonstil,firstnumber=276]
    auf_rate = ziel == ZIELGROESSE
    y_tr = train[RATE if auf_rate else ziel].astype(float)
    zurueck = (test[EXPOSURE_ROH].astype(float).to_numpy() / 1000.0
               if auf_rate else 1.0)

    modell = verfahren(name).set_params(**parameter)(*@\setcounter{lstnumber}{283}@*)
    modell.fit(X_tr, y_tr)(*@\setcounter{lstnumber}{287}@*)
    y_hat = modell.predict(X_te) * zurueck
\end{lstlisting}
\tabquelle[Eigene Darstellung, \texttt{modelle/m02\_menge.py}, Z.~276--288, gekürzt]
\end{quellcode}


Verschieden sein muss dagegen, was ein Verfahren minimiert, denn das hängt an
der Form der Zielgröße. Auf den Trainingsstadtteilen des ersten Folds
übersteigt die Varianz der Einsatzzahlen ihren Mittelwert um das 80-Fache. Ein
quadratischer Fehler gewichtet große Abweichungen quadratisch und wird dann von
wenigen großen Stadtteilen bestimmt --- eine Eigenschaft, die bei
rechtsschiefen Zielgrößen bekannt ist.\footcite[349]{Hastie2009} Random Forest
minimiert deshalb die Poisson-Devianz und XGBoost eine Tweedie-Devianz, deren
Varianzexponent mitgesucht wird.

Ridge hat kein zähldatengerechtes Gegenstück und wird stattdessen auf
$\log(1+y)$ geschätzt; die Gütemaße entstehen nach der Rücktransformation auf
der Originalskala. Eine Korrektur der Rücktransformationsverzerrung wird nicht gerechnet: expm1 liefert näherungsweise den bedingten Median statt des Erwartungswerts. Sie wirkt zulasten von Ridge und nicht zu seinen Gunsten — alle drei Verfahren unterschätzen im Mittel, Ridge um 7,1, XGBoost um 5,8 und der Random Forest um 2,8 Einsätze, und das Entfernen dieser Verschiebung senkt den Fehler bei allen dreien um eine vergleichbare Größe. Der Grund ist messbar: Die Anpassung im Training steigt
durch die Transformation von einem Bestimmtheitsmaß von 0,754 auf 0,859.
Dieselbe Sonderstellung hat die Standardisierung der Merkmale --- sie wird je
Fold allein auf den Trainingszeilen angepasst und betrifft nur Ridge, weil die
beiden Baumverfahren ausschließlich die Ordnung der Werte nutzen und gegenüber
deren Skala unempfindlich sind. Über die Prognosegüte sagt die Zahl nichts: Sie
ist an denselben Zeilen gemessen, an denen das Modell angepasst wurde.

Dass diese Wahl den Wertebereich der Zielgröße respektiert, ist geprüft: Poisson- und Tweedie-Devianz arbeiten über eine logarithmische Verknüpfung und können keine negative Vorhersage erzeugen, bei Ridge lässt die Rücktransformation sie zu. Jeder Lauf schreibt deshalb die Zahl negativer Vorhersagen mit, ohne sie abzuschneiden; über alle 300 Läufe des Mengenstrangs bleibt sie bei null.

\begin{quellcode}[htbp]
\setstretch{1}
\caption{Verlustfunktion und Zieltransformation je Verfahren}
\label{lst:verfahren}\vspace{-3.5pt}
\begin{lstlisting}[style=pythonstil,firstnumber=130]
def verfahren(name: str, n_jobs: int = N_JOBS_MODELL):(*@\setcounter{lstnumber}{158}@*)
    if name == "ridge":
        return make_pipeline(
            StandardScaler(),
            TransformedTargetRegressor(regressor=Ridge(),
                                       func=np.log1p, inverse_func=np.expm1))
    if name == "random_forest":
        return RandomForestRegressor(random_state=RANDOM_STATE, n_jobs=n_jobs,
                                     criterion="poisson")
    if name == "xgboost":
        from xgboost import XGBRegressor
        return XGBRegressor(random_state=RANDOM_STATE, n_jobs=n_jobs,
                            objective="reg:tweedie")
\end{lstlisting}
\tabquelle[Eigene Darstellung, \texttt{modelle/m02\_menge.py}, Z.~130--170, gekürzt]
\end{quellcode}

Alle drei Verfahren entstehen in derselben Funktion \texttt{verfahren} (Code~\ref{lst:verfahren}) und unterscheiden sich allein in ihrem Rückgabezweig: Die Poisson-Devianz steht als \texttt{criterion} in Zeile~166, die Tweedie-Verlustfunktion als \texttt{objective} in Zeile~170. Dass die Sonderbehandlung von Ridge in den Zeilen~160 bis~163 innerhalb der Pipeline steht und nicht davor, ist der Grund, warum Standardisierung und Rücktransformation je Fold ausschließlich die Trainingszeilen sehen; vorgelagert angewandt, hätten beide die Testzeilen mitgelesen. Beide Baumverfahren beziehen ihren \texttt{random\_state} aus derselben Konstante, und \texttt{n\_jobs} ist über \texttt{N\_JOBS\_MODELL} auf einen Kern voreingestellt, damit die berichteten Laufzeiten vergleichbar bleiben.


Im Strukturstrang betrifft die Sonderbehandlung nicht die Skala, sondern die
Klassen. Ihre vier Ausprägungen werden einmal global in Ganzzahlen übersetzt
und nicht je Fold, weil sonst in einem Fold ohne brand-dominierten Monat die
Wahrscheinlichkeitsspalten auf die falschen Klassen zeigten, ohne dass eine
Fehlermeldung entstünde. Die Schieflage zwischen den Klassen wird über
ausgeglichene Gewichte behandelt und nicht über Resampling: Keine Zeile wird
vervielfacht, keine entfernt --- die Zeilenmenge bleibt damit dieselbe wie im
Mengenstrang.

\begin{quellcode}[htbp]
\setstretch{1}
\caption{Globale Klassenkodierung und ausgeglichene Gewichte}
\label{lst:kodiere}\vspace{-3.5pt}
\begin{lstlisting}[style=pythonstil,firstnumber=162]
def kodiere(y: pd.Series) -> np.ndarray:(*@\setcounter{lstnumber}{172}@*)
    index = {k: i for i, k in enumerate(KLASSEN)}
    unbekannt = set(y.unique()) - set(index)
    assert not unbekannt, f"Unbekannte Klassen im Datensatz: {unbekannt}"
    return y.map(index).to_numpy()(*@\setcounter{lstnumber}{178}@*)
def _gewichte(y_int: np.ndarray) -> np.ndarray:(*@\setcounter{lstnumber}{185}@*)
    return compute_sample_weight("balanced", y_int)
\end{lstlisting}
\tabquelle[Eigene Darstellung, \texttt{modelle/m03\_struktur.py}, Z.~162--186, gekürzt]
\end{quellcode}

Drei Festlegungen betreffen nicht die Modelle, sondern die Maschine, auf der
sie rechnen. Die erste ist die Umgebung selbst: Tabelle~\ref{tab:techstack}
nennt sie, und die Verfahren stammen aus scikit-learn, das sie über eine
einheitliche Schnittstelle bereitstellt und selbst nur numpy und scipy
voraussetzt.\footcite[2826]{Pedregosa2011} Ohne Prozessor, Kernzahl und
Bibliotheksversionen wären die in Abschnitt~\ref{sec:aufwand} berichteten
Sekunden keine Größe, die sich andernorts wiederholen ließe.

\begin{table}[htbp]
\tabformat
\caption{Rechenumgebung und verwendete Bibliotheken}
\label{tab:techstack}
\begin{tabularx}{\textwidth}{@{}L{4.2cm} Y@{}}
\toprule
\textbf{Bestandteil} & \textbf{Ausprägung} \\
\midrule
Prozessor & Intel Core i5-7300U, 2,60 GHz, 2 physische/4 logische Kerne \\
Umgebung & 7,8 GB Arbeitsspeicher, Windows 10 Pro, Python 3.14.0 \\
\addlinespace
Modellierung & scikit-learn 1.8.0, xgboost 3.4.0 \\
Referenz und Tests & statsmodels 0.14.6, scipy 1.17.1 \\
Datenhaltung und Raumbezug & pandas 3.0.2, pyarrow 24.0.0, polars 1.40.1, geopandas 1.1.3, shapely 2.1.2, pyproj 3.7.2 \\
Erklärbarkeit und Abbildungen & shap 0.52.0, matplotlib 3.10.9, seaborn 0.13.2 \\
\addlinespace
Vollständige Liste & \texttt{requirements\_lauf.txt}, 41 Pakete \\
\bottomrule
\end{tabularx}
\tabquelle
\end{table}

Die zweite betrifft die Parallelisierung. Anpassen ließe sich auf zwei Ebenen
zugleich, genutzt wird nur eine: Die Zufallssuche verteilt ihre Ziehungen über
alle logischen Kerne, der Schätzer darin bleibt einkernig. Gemessen an einem
Fold kostet die geschachtelte Variante auf zwei physischen Kernen das
Vierzehnfache bei Ridge und das Dreifache beim Random Forest; hochgerechnet
fällt der volle Lauf damit von rund sieben auf gut zwei Stunden.

\begin{quellcode}[htbp]
\setstretch{1}
\caption{Innere Suche, nach Stadtteilen gruppiert}
\label{lst:tune}\vspace{-3.5pt}
\begin{lstlisting}[style=pythonstil,firstnumber=210]
def tune(name: str, train: pd.DataFrame, ziel: str) -> dict:(*@\setcounter{lstnumber}{227}@*)
    suche = RandomizedSearchCV(
        estimator=verfahren(name, n_jobs=N_JOBS_MODELL),
        param_distributions=suchraum(name),
        n_iter=TUNING_BUDGET,
        cv=GroupKFold(n_splits=4),
        scoring="neg_root_mean_squared_error",
        random_state=RANDOM_STATE,
        n_jobs=N_JOBS_SUCHE,
    )
    suche.fit(train[MERKMALE].astype(float), train[ziel].astype(float),
              groups=train["stadtteil"])
    return suche.best_params_
\end{lstlisting}
\tabquelle[Eigene Darstellung, \texttt{modelle/m02\_menge.py}, Z.~210--239, gekürzt]
\end{quellcode}

Die dritte ist die Zahl der Kerne im Bewertungslauf. Ob sie allein die Dauer
betrifft oder auch das Ergebnis, prüft der
Lauf selbst: Jeder Fit wird ein zweites Mal über alle Kerne gerechnet und beide
Vorhersagen werden verglichen. Für Ridge und Random Forest stimmen sie überein,
für XGBoost nicht. Dort summiert die parallele Zusammenfassung der Histogramme
in wechselnder Reihenfolge, was knapp benachbarte Split-Kandidaten kippen lässt
und sich über hunderte Bäume aufschaukelt, gemessen bis zu 50,88 Einsätze bei
einem Mittelwert von 75. Sämtliche berichteten Gütemaße stammen deshalb aus dem
einkernigen Fit; ohne die Kernzahl wäre die Angabe zur Reproduzierbarkeit
unvollständig.

\begin{quellcode}[htbp]
\setstretch{1}
\caption{Derselbe Fit einkernig und über alle Kerne}
\label{lst:parallel}\vspace{-3.5pt}
\begin{lstlisting}[style=pythonstil,firstnumber=293]
if auch_parallel:
    parallel = verfahren(name, n_jobs=-1).set_params(**parameter)
    t = time.perf_counter()
    parallel.fit(X_tr, y_tr)
    train_par = time.perf_counter() - t(*@\setcounter{lstnumber}{298}@*)
    y_par = parallel.predict(X_te) * zurueck(*@\setcounter{lstnumber}{309}@*)
    abweichung = float(np.max(np.abs(y_hat - y_par)))
\end{lstlisting}
\tabquelle[Eigene Darstellung, \texttt{modelle/m02\_menge.py}, Z.~293--310, gekürzt]
\end{quellcode}

% NOTIZ (07.09.2026) -- neu. Schliesst 6.1 ab: Nach Schaetzerbau, Zielgroesse,
% Klassenkodierung, innerer Suche und Parallelmessung steht hier der Lauf, der
% alles zusammenfuehrt. Gegenstueck zum Baseline-Lauf in Abschnitt 5.4 -- gleiche
% Schleifen, gleiche Aufteilung.
% Erklaerbar im Text: (1) die vier geschachtelten Schleifen (Z. 482, 484, 487,
% 488) ergeben 10 x 5 x 2 x 3 Laeufe, und der Bewertungslauf sieht dieselben
% Wiederholungen und Folds wie die Baselines; (2) Zeile 490 setzt die Parameter
% aus Phase 1 in ein FRISCHES Modell -- best_estimator_ aus der Suche waere nur
% auf drei Vierteln der Trainingsstadtteile gefittet und der Vergleich waere
% nicht mehr gleich; (3) auch_parallel=True in Zeile 491 ist der Grund, warum die
% Kernzahl-Diagnose des vorigen Absatzes in JEDEM Lauf anfaellt und nicht in
% einer Stichprobe.
\begin{quellcode}[htbp]
\setstretch{1}
\caption{Der Bewertungslauf über alle Wiederholungen und Folds}
\label{lst:modelllauf}\vspace{-3.5pt}
\begin{lstlisting}[style=pythonstil,firstnumber=468]
def phase_bewertung(panel: pd.DataFrame, parameter: pd.DataFrame,
                    selten: pd.Series) -> pd.DataFrame:(*@\setcounter{lstnumber}{479}@*)
    param = _parameter_je_fold(parameter)
    zeilen, vorhersagen = [], []
    for w in range(WIEDERHOLUNGEN):
        d = wiederholte_aufteilung(panel, wiederholung=w, selten=selten)
        for k in range(1, N_FOLDS + 1):
            tr, te = fold_masken(d, k)
            train, test = d[tr], d[te]
            for ziel in ZIELE:
                for name in VERFAHREN:
                    # Parallelmessung in jedem Lauf - keine Ausnahmen.
                    z = ein_lauf(name, param[(ziel, name, k)], train, test,
                                 ziel, auch_parallel=True,
                                 mit_vorhersagen=True)
                    vorhersagen.append(z.pop("_vorhersagen")
                                        .assign(wiederholung=w, fold=k))
                    zeilen.append({"wiederholung": w, "fold": k, **z})
\end{lstlisting}
\tabquelle[Eigene Darstellung, \texttt{modelle/m02\_menge.py}, Z.~468--495, gekürzt]
\end{quellcode}



Damit steht der Rahmen. Er entscheidet nichts darüber, welche Verfahren in
ihm verglichen werden --- er sorgt allein dafür, dass ein Unterschied
zwischen ihnen auch einer ist.

\subsection{Eignungsprüfung und Begründung der Verfahrenswahl}
\label{sec:eignung}
% Die BASELINES sind hierher NICHT mehr zu schreiben -- sie stehen seit dem
% Umbau in Abschnitt 5.4 (Auflage Schroeter, 27.07.2026). Der frueher hier
% stehende Expose-Absatz ist dorthin gewandert.
%
% Dieser Abschnitt beantwortet stattdessen die Frage, warum ueber ein
% einfaches Regressionsmodell hinausgegangen wird. Die Begruendung entsteht
% aus der Regression selbst, in vier Schritten:
%
% 1. Das einfache Modell funktioniert. Das unpenalisierte Poisson-GLM erreicht
%    in der Kreuzvalidierung R2 0,542 bei RMSE 33,98, Ridge auf log(1+y)
%    R2 0,514 bei 36,51. Ein lineares Modell ist also ein ernstzunehmender
%    Kandidat, kein Strohmann -- und wie sich in Kapitel 7 zeigt, schlaegt
%    KEIN Vergleichsverfahren das GLM.
%    (Fruehere Fassung nannte hier R2 0,472 gegen die Negative Binomial und
%    0,588 fuer ein L2-regularisiertes Poisson-GLM. Beide Zahlen stammen aus
%    einer ueberholten Konfiguration und sind seit #45/#52 nicht mehr zu
%    verwenden -- es wird genau ein Lauf berichtet.)
% 2. Der formale Spezifikationstest verwirft es trotzdem. RESET-Test nach
%    Ramsey (1969), H0 "die lineare Spezifikation ist adaequat":
%      Potenzen bis 2: F = 215,23  p = 4e-47   -> H0 verworfen
%      Potenzen bis 3: F = 125,59  p = 4e-53   -> H0 verworfen
% 3. Die Ursache ist lokalisierbar, und sie ist ZWEIFACH (Lauf 03.08.2026):
%    (a) KRUEMMUNG. Von den sechs Merkmalen mit substanzieller Korrelation
%        (|Pearson| > 0,2) zeigt log_bevoelkerung den groessten Abstand
%        zwischen Pearson (+0,416) und Spearman (+0,559) -- ein monotoner,
%        aber gekruemmter Zusammenhang. Wichtig fuer den Text: Bei den
%        uebrigen vier Merkmalen liegt die Korrelation nahe null, dort ist
%        der Pearson-Spearman-Abstand Rauschen und KEIN Befund.
%    (b) WECHSELWIRKUNGEN. Ergaenzt man die 45 Interaktionsterme, steigt das
%        adjustierte R2 von 0,805 auf 0,919.
%    ACHTUNG: Eine fruehere Fassung dieses Kommentars behauptete, die
%    Einzeleffekte seien praktisch linear und nur Wechselwirkungen fehlten.
%    Das ist durch den Lauf vom 03.08.2026 widerlegt und NICHT mehr zu
%    schreiben.
%    UNBEDINGT DAZU -- SONST WIDERSPRICHT 6.2 DEM KAPITEL 7 (B-41):
%    Die hier diagnostizierte Nichtlinearitaet ist IN-SAMPLE nachgewiesen. Der
%    RESET-Test laeuft auf 3.828 Zeilen, die als unabhaengig behandelt werden --
%    tatsaechlich liegen 29 unabhaengige Stadtteile vor. Die Gegenprobe
%    (vorpruefung/v3_spezifikation.py, derselbe Split, dieselben 50 Laeufe)
%    zeigt, dass sich die Struktur NICHT auf unbekannte Stadtteile uebertraegt:
%      linear (12 Terme)        RMSE  33,98   R2   0,542
%      + Quadrate (22)                101,11      -5,971
%      + Interaktionen (57)           121,63      -6,247
%      + beides (67)                  180,86     -16,454
%    Alle 50 Laeufe konvergiert, also kein numerisches Artefakt.
%    WIE ES ZU SCHREIBEN IST: Die Eignungspruefung begruendet, warum ueber ein
%    einfaches lineares Modell HINAUSGEGANGEN wurde -- das ist eine Entscheidung
%    VOR dem Ergebnis und bleibt richtig. Dass sie sich nicht ausgezahlt hat,
%    ist der Befund von Kapitel 7 und gehoert dorthin, nicht hierher.
%    Ein Vorwaertsverweis genuegt. Was NICHT passieren darf: dass 6.2 die
%    Baumverfahren als zwingend darstellt und Kapitel 7 dann das Gegenteil
%    berichtet, ohne dass die Bruecke gebaut ist.
%
% 4. Daraus folgt die Verfahrenswahl: Ein lineares Modell bildet Kruemmung
%    nur ueber vorab gewaehlte Transformationen ab und Interaktionen nur,
%    wenn man sie von Hand spezifiziert -- bei zehn Merkmalen 45 Terme,
%    willkuerlich ausgewaehlt und bei 23 Trainingsstadtteilen (Fold 1)
%    ueberangepasst. Baumverfahren fangen beides konstruktionsbedingt ab:
%    Ein Split kann an beliebiger Stelle schneiden, und jeder Split bedingt
%    auf die vorherigen.
%
% Ebenfalls hierher: die Linearitaetspruefung als Auflage Schroeter (R7,
% Scatterplots und Residuenanalyse VOR dem Einsatz von Ridge), der VIF
% (hoechster Wert 12,29 bei median_haushaltseinkommen -- daher die BLOCKWEISE
% Auswertung in 7.4, einzelne Merkmalsbeitraege waeren Scheinpraezision),
% sowie die Extrapolationspruefung.
%
% ZUR EXTRAPOLATION -- ZWEI KORREKTUREN GEGEN DIE FRUEHERE FASSUNG:
%   (a) Die Zahl ist nicht 31 %. Sie lautet 33,7 % (Mittel der Wiederholung 0),
%       34,6 % (ueber alle 50 Laeufe) und 7,6 % (Hold-out). Wo eine davon
%       steht, ist die BEZUGSMENGE zu nennen. Je Fold: 40,9 / 33,3 / 57,4 /
%       33,3 / 3,6 %.
%   (b) Die Behauptung, sie treffe Baumverfahren und Ridge unterschiedlich,
%       ist GEMESSEN UND WIDERLEGT (B-31). Der Zusammenhang zwischen
%       Extrapolationsanteil und RMSE ist bei RIDGE am staerksten
%       (Spearman +0,302) und beim Random Forest am schwaechsten (+0,108,
%       p 0,456) -- waere Extrapolation der Hebel gegen die Baumverfahren,
%       muesste es umgekehrt sein. Diese Erklaerung darf NICHT mehr auftauchen.
%       Die Extrapolation bleibt eine Eigenschaft des Validierungsrahmens und
%       begrenzt die Generalisierbarkeit (Limitation, 8.3) -- sie erklaert
%       aber keinen Verfahrensunterschied.
%   (c) Sie ist eine Eigenschaft von STADTTEILEN, nicht von Zeilen: 9 von 29
%       Stadtteilen liegen mit 100 % ihrer Zeilen ausserhalb, 16 von 29 mit
%       0 %. Weil die Strukturmerkmale innerhalb eines Stadtteils nahezu
%       konstant sind, bricht ein Stadtteil ganz aus oder gar nicht.
%
% AUFLAGE SCHROETER 10.08.2026 -- ANFORDERUNGEN JE VERFAHREN. Hierher gehoert
% die Tabelle aus 03_STAND.md Abschnitt 7: 15 gepruefte Anforderungen, davon 7
% verletzt, 3 fuer das jeweilige Verfahren gar nicht bestehend. Mit
% Teststatistik und p-Wert, wie verlangt:
%     Cameron & Trivedi (1990), Hilfsregression   t  =  17,2   p < 0,001
%     Breusch-Pagan auf log(1+y)                  LM = 722,7   p < 0,001
%     Jarque-Bera auf log(1+y)                    JB =  98,2   p < 0,001
%   Schiefe -0,21, Woelbung 3,78. Abbildung A10 (QQ-Plot der Residuen): Kern
%   auf der Geraden, nur die Raender biegen ab -- schwere Raender, kein
%   verbogener Kern.
%   DIE ZEILEN "nicht erforderlich" SIND KEIN FUELLMATERIAL: Eine Tabelle, die
%   nur verletzte Annahmen zeigt, liesse die Baumverfahren voraussetzungslos
%   aussehen. Deren eigentliche Anforderung ist der Interpolationsbereich --
%   und die ist verletzt.
%
% AUFLAGE SCHROETER (E-Mail 03.08.2026): Hierher gehoert ausserdem die
% Begruendung, WARUM Regression und Klassifikation unterschiedlich viele
% Verfahren einsetzen -- drei gegen zwei. Schroeter woertlich: "Sie muessen
% nicht fuer Regression und Klassifikation dieselbe Anzahl oder dieselben
% Algorithmen verwenden. Entscheidend ist, dass die Verfahren zur jeweiligen
% Zielgroesse passen." Die Argumentation, ausformuliert in Decision Log #31:
%
% a) Die Klassifikationsmenge ist eine ECHTE TEILMENGE der Regressionsmenge.
%    Es kommt kein Verfahren hinzu, genau eines faellt weg. Damit entstehen
%    gerade KEINE heterogenen Spezifikationen (Gutachten R1) -- beide
%    Zielgroessen laufen durch denselben Preprocessing- und
%    Validierungsrahmen und werden von denselben Modellfamilien bearbeitet.
% b) RF und XGBoost uebertragen sich, weil bei ihnen nur die Verlustfunktion
%    wechselt (Gini/Entropie statt Varianzreduktion beim Split;
%    multi:softprob statt reg:squarederror). Der Ensemble-Mechanismus ist
%    identisch, beide sind nativ mehrklassenfaehig.
% c) Ridge uebertraegt sich nicht: Es minimiert den quadratischen Fehler auf
%    einer metrischen Zielgroesse. dominante_einsatzart ist nominal mit vier
%    ungeordneten Klassen -- ohne Ordnung und Abstand ist der quadratische
%    Fehler nicht definiert. RidgeClassifier wurde geprueft und verworfen
%    (+-1-Kodierung, One-vs-Rest, keine kalibrierten Wahrscheinlichkeiten).
%    Das Pruefen und Verwerfen gehoert in den Text -- es zeigt, dass die
%    Option gesehen und nicht uebersehen wurde.
% d) Die logistische Regression waere das korrekte lineare Gegenstueck und
%    schlaegt im Vortest die Baseline (Macro-F1 0,282 gegen 0,220 der
%    Mehrheitsklasse, RF 0,301). Sie entfaellt aus FOKUSGRUENDEN, nicht
%    mangels Eignung: ein vierter Verfahrensstrang bei zwei Zielgroessen ist
%    genau die Verbreiterung, die das Gutachten unter R4 als groesstes
%    Notenrisiko benennt. DAS IST SO ZU SCHREIBEN -- eine funktionierende
%    Option als untauglich darzustellen waere unbelegt (R9).

Das einfachste Modell, das zu diesen Daten passt, funktioniert bereits: Das
Poisson-Modell der Stufe~2 erreicht die in
Abschnitt~\ref{sec:rahmen-baselines} ausgewiesene Güte, ohne einen einzigen frei
wählbaren Parameter. Aufwendigere Verfahren
einzusetzen ist deshalb begründungspflichtig. Die Frage ist nicht, ob das
lineare Modell rechnet, sondern ob seine Bauform zur Struktur der Daten passt
--- und es legt zwei Dinge fest: dass sich jedes Merkmal über den gesamten
Wertebereich gleich auswirkt, und dass diese Wirkung nicht davon abhängt,
welche Werte die übrigen Merkmale annehmen.

Beide Festlegungen halten einer Prüfung nicht stand. Trifft ein Modell die Form
des Zusammenhangs, dann steckt in seinen eigenen Vorhersagen keine Information
mehr, die es nicht schon genutzt hätte. Der RESET-Test nimmt genau diese
Vorhersagewerte, fügt ihre Potenzen als zusätzliche Merkmale hinzu und misst,
ob die Anpassung dadurch besser wird.\footcite[]{Ramsey1969} Hier wird sie deutlich
besser, $F = 360{,}4$ bei $p < 0{,}001$, gerechnet auf den 24
Trainingsstadtteilen des ersten Folds.

Die erste Ursache ist Krümmung. Alle zehn Strukturmerkmale erreichen mit der
Einsatzhäufigkeit eine Korrelation über 0,20 in mindestens einem der beiden
Maße; am stärksten weicht der Rangkorrelationskoeffizient bei der Armutsquote
vom Korrelationskoeffizienten nach Bravais-Pearson ab, $+0{,}342$ gegen
$+0{,}617$. Der Abstand beider Maße zeigt einen monotonen, aber gekrümmten
Zusammenhang (Abschnitt~\ref{sec:eda}); dass Pearson hier über Spearman liegt,
weist auf einzelne Hebelpunkte hin. Für die beiden Saisonterme gilt dieser
Schluss nicht: Dort liegen beide Koeffizienten nahe null, und ihr Abstand ist
Rauschen.

Die zweite Ursache sind Wechselwirkungen. Ergänzt man die lineare Spezifikation
um alle 45 paarweisen Produkte der zehn Merkmale, steigt das adjustierte
Bestimmtheitsmaß von 0,859 auf 0,932. Beide Befunde beziehen sich auf dieselben
Zeilen, an denen das Modell angepasst wurde: Sie belegen, dass die lineare Form
die vorliegenden Stadtteile schlechter beschreibt als eine flexiblere Form ---
nicht, dass sie auf unbekannten Stadtteilen schlechter vorhersagt.

Diese zweite Frage beantwortet eine eigene Gegenprobe: Sie ergänzt systematisch
die Terme, deren Fehlen der Test anzeigt --- quadratische für alle zehn
Merkmale, Produkte für alle Paare ---, und wertet sie im selben Rahmen aus wie
die Vergleichsverfahren. Ihr Ergebnis steht in
Abschnitt~\ref{sec:ergebnisse-menge}.

Aus beiden Ursachen folgt die Verfahrenswahl. Ein lineares Modell bildet
Krümmung nur über vorab festzulegende Transformationen ab und Wechselwirkungen
nur über von Hand spezifizierte Produktterme, für deren Auswahl es keine Regel
gibt und die bei 24 Trainingsstadtteilen je Fold überanpassen. Baumverfahren
brauchen beides nicht (Abschnitt~\ref{sec:referenzmodelle}).

Die drei Vergleichsverfahren decken drei verschiedene Antworten auf dasselbe
Problem ab (Abschnitt~\ref{sec:untersuchte-verfahren}). Ridge dämpft die Folgen
korrelierter Merkmale --- der höchste Varianzinflationsfaktor liegt hier bei
10,64 für das mittlere Haushaltseinkommen ---, der Random Forest senkt die
Varianz über das Mittel vieler Bäume, XGBoost setzt an der Verzerrung an. Der
Vergleich fragt damit, welche der beiden Fehlerquellen bei diesen Daten
überwiegt.

Ebenso begründungspflichtig ist, was nicht verglichen wird. Lasso und Elastic
Net setzen Koeffizienten auf exakt null und betreiben damit Merkmalsselektion;
die zehn Merkmale sind aber aus drei Faktorgruppen abgeleitet, deren
Erklärungsbeitrag die erste Unterfrage gerade prüfen soll --- ein Verfahren,
das eine Gruppe entfernt, umginge die Frage. LightGBM und CatBoost wären eine
weitere Variante derselben Familie ohne neue Frage. Neuronale Netze scheiden
aus, weil 36 Analyseeinheiten für sie zu wenig sind und Baumverfahren auf
tabellarischen Daten dieser Größenordnung ohnehin vorn
liegen.\footcite[1]{Grinsztajn2022}

Auf der Einsatzart wird nur mit zwei der drei Verfahren verglichen, und die
Auswahl folgt der Skala der Zielgröße. Random Forest und XGBoost übertragen
sich ohne Bruch: Es wechselt allein die Verlustfunktion, der
Ensemble-Mechanismus bleibt derselbe, und beide sind ohne Umweg
mehrklassenfähig. Ridge überträgt sich nicht, weil es den quadratischen Fehler
minimiert und dieser auf einer nominalen Zielgröße mit vier ungeordneten
Klassen nicht definiert ist. Die klassifizierende Variante wurde geprüft und
verworfen: Sie liefert keine kalibrierten Wahrscheinlichkeiten, ohne die sich
die Macro-AUROC nicht bilden lässt.

Das lineare Gegenstück zu den beiden Baumverfahren ist damit die multinomiale
logistische Regression --- und sie ist im Rahmen, allerdings als Referenz der
Stufe~2 und nicht als viertes Vergleichsverfahren
(Abschnitt~\ref{sec:rahmen-baselines}). Als eigener Verfahrensstrang entfällt
sie aus Gründen des Zuschnitts, nicht mangels Eignung: Ein vierter Strang
verbreiterte den Vergleich, ohne eine neue Frage zu beantworten. Im Strukturstrang ist das lineare Modell dadurch die Messlatte und nicht der
Mitbewerber.

Welche Anforderungen die vier Modelle aus der Herleitung ihres Schätzers mitbringen und welche davon diese Daten verletzen, zeigt Tabelle~\ref{tab:anforderungen}: Von fünfzehn geprüften Anforderungen sind sieben verletzt und drei für das jeweilige Verfahren gar nicht zu stellen; unabhängige Beobachtungen setzen alle vier voraus, und die Panelstruktur verletzt diese eine Annahme für alle gleichermaßen (Abschnitt~\ref{sec:rahmen-baselines}). Die verletzten Annahmen betreffen durchweg die Inferenz und nicht die Punktprognose: Bei ungleicher Fehlervarianz verliert die Kleinste-Quadrate-Schätzung an Effizienz, vor allem aber werden die geschätzten Standardfehler verzerrt\footcite[1287]{BreuschPagan1979} --- und Standardfehler werden hier nirgends berichtet.

% NOTIZ (07.09.2026) -- neu in 6.2. Von den drei formalen Tests der
% Anforderungstabelle ist dies der einzige, der nicht aus einer Bibliothek
% kommt: statsmodels kennt ihn nicht, er ist aus Cameron & Trivedi (1990)
% gebaut. Breusch-Pagan und Jarque-Bera sind je ein Aufruf und stehen deshalb
% bewusst NICHT als Quellcode in der Arbeit.
% Erklaerbar im Text (mit Zeilenbezug):
% (1) Zeile 458 ist die Hilfsgroesse z = ((y-mu)^2 - y)/mu -- ihr Erwartungswert
%     ist unter Equidispersion null, unter Ueberdispersion proportional zu mu;
% (2) Zeile 459 regressiert z OHNE Konstante auf mu, weil der Koeffizient genau
%     das alpha der NB2-Form ist und H0 alpha = 0 lautet;
% (3) die Halbierung des p-Werts in Zeile 460 macht den Test einseitig --
%     Unterdispersion waere bei Einsatzzahlen keine sinnvolle Gegenhypothese;
% (4) der Offset in Zeile 456 ist derselbe wie in der Stufe-2-Baseline
%     (Code~\ref{lst:poisson}): geprueft wird das Modell, das spaeter die
%     Messlatte ist, und nicht ein anderes.
\begin{quellcode}[htbp]
\setstretch{1}
\caption{Überdispersionstest als Hilfsregression}
\label{lst:dispersionstest}\vspace{-3.5pt}
\begin{lstlisting}[style=pythonstil,firstnumber=448]
    X = sm.add_constant(train[MERKMALE].astype(float), has_constant="add")
    y = train[ZIELGROESSE].astype(float).to_numpy()(*@\setcounter{lstnumber}{454}@*)
    poisson = sm.GLM(y, X, family=sm.families.Poisson(),
                     offset=np.log(train[EXPOSURE_ROH].astype(float))).fit()
    mu = np.asarray(poisson.fittedvalues, float)
    z = ((y - mu) ** 2 - y) / mu
    hilfs = sm.OLS(z, mu).fit()
    ct_t, ct_p = float(hilfs.tvalues[0]), float(hilfs.pvalues[0]) / 2
\end{lstlisting}
\tabquelle[Eigene Darstellung, \texttt{vorpruefung/v2\_eignung.py}, Z.~448--460, gekürzt]
\end{quellcode}

\begin{table}[htbp]
\tabformat
\caption{Anforderungen je Verfahren, geprüft auf den Trainingsstadtteilen des ersten Folds; alle ausgewiesenen Teststatistiken sind auf dem Niveau $p < 0{,}001$ signifikant}
\label{tab:anforderungen}
\begin{tabularx}{\textwidth}{@{}L{2.1cm} L{3.3cm} L{3.4cm} Y@{}}
\toprule
\textbf{Verfahren} & \textbf{Anforderung} & \textbf{Befund} & \textbf{Konsequenz} \\
\midrule
alle & unabhängige Beobachtungen & verletzt, Panelstruktur & Stadtteil-Split, Streuung über die Wiederholungsmittel \\
alle & identische Merkmale und Folds & erfüllt, Spalte in der Datei & konstruktiv abgesichert \\
\addlinespace
Poisson-Modell & Equidispersion & verletzt, $t = 18{,}0$ & folgenlos, nur Punktvorhersagen \\
Poisson-Modell & Linearität im Log-Link & verletzt, $F = 360{,}4$ & bewusst in Kauf genommen \\
\addlinespace
Ridge & Linearität & verletzt, $F = 360{,}4$ & Schätzung auf $\log(1+y)$ \\
Ridge & Homoskedastizität & verletzt, $LM = 525{,}1$ & betrifft Standardfehler, nicht die Prognose \\
Ridge & normalverteilte Residuen & nicht erforderlich, $JB = 150{,}0$ & Voraussetzung der Inferenz, nicht der Prognose \\
\addlinespace
Random Forest, XGBoost & Verteilungsannahme & nicht erforderlich & verteilungsfrei \\
Random Forest, XGBoost & Testpunkte im gelernten Bereich & verletzt, 38,2\% außerhalb (Wiederholung~0) & Limitation des Stadtteil-Splits \\
Random Forest, XGBoost & Verlust zur Datenform & erfüllt, Index 79,5 & Poisson- bzw. Tweedie-Devianz, Varianzexponent getunt \\
\addlinespace
Logit & Linearität in den Log-Odds & angenommen, nicht formal geprüft & Trennlinie zu den Baumverfahren \\
Logit & jede Klasse im Testfold besetzt & erfüllt, Selbsttest & doppelte Stratifizierung \\
\bottomrule
\end{tabularx}
\tabquelle
\end{table}

Die Prüfung auf ungleiche Fehlervarianz ist in der studentisierten Fassung gerechnet, weil ihre ursprüngliche Form normalverteilte Störgrößen voraussetzt\footcite[1288]{BreuschPagan1979} und genau diese Annahme die Zeile darunter mit $JB = 150{,}0$: verwirft; die zweite Tabellenzeile ist damit der Grund für die Form der ersten. Woraus die Abweichung besteht, zeigt Abbildung~\ref{abb:qq}: aus schweren Rändern und nicht aus einem verbogenen Kern.

\begin{figure}[htbp]
  \centering
  \includegraphics[width=\textwidth]{Abbildungen/a10_qq_residuen.pdf}
  \caption{Verteilung der Residuen des linearen Modells}
  \label{abb:qq}
\end{figure}

Verletzt ist auch die Anforderung der Baumverfahren: Ein Baum sagt außerhalb des
im Training beobachteten Wertebereichs den Wert des nächstgelegenen bekannten
Blattes vorher, und der davon betroffene Anteil der Testzeilen
(Tabelle~\ref{tab:anforderungen}) streut über alle fünfzig Läufe zwischen 0,0 und 83,3 Prozent. Als Erklärung für Unterschiede zwischen den Verfahren taugt der Befund
gleichwohl nicht: Der Zusammenhang zwischen dem Extrapolationsanteil eines Folds
und dem dort erreichten Fehler ist bei Ridge am stärksten und beim Random Forest
am schwächsten, mit Rangkorrelationen von $+0{,}335$ gegen $+0{,}249$ --- wäre
die Extrapolation der Hebel gegen die Baumverfahren, müsste es umgekehrt sein.
Was bleibt, ist eine Eigenschaft des Validierungsrahmens, deren Folgen für die
Übertragbarkeit in Abschnitt~\ref{sec:limitationen} stehen.


\subsection{Modellkonfiguration und Hyperparameter-Suchräume}
% Tabelle: Verfahren | Hyperparameter | Suchraum | Iterationen | bester Wert.

% ---- ÜBERNAHME AUS EXPOSÉ (Kap. 3, Abs. 3) -- 1:1, Tempus angepasst ----

% ERSETZT AM 26.08.2026 durch den ausformulierten Abschnitt darunter. Der
% Absatz stammt 1:1 aus dem Expose (Kap. 3, Abs. 3); sein Inhalt steht jetzt
% in Absatz 1 (Randomized Search), Absatz 3 (Probst2019, warum trotzdem
% getunt wird) und in Tabelle tab:suchraeume. Auskommentiert statt
% geloescht, damit der Abgleich mit dem Expose moeglich bleibt:
% In der Modeling-Phase wurden Ridge Regression, Random Forest und XGBoost auf beiden Zielgrößen,
% Einsatzhäufigkeit (Regression) und Einsatzart (Klassifikation), trainiert. Für jedes Modell erfolgt
% ein Hyperparameter-Tuning mittels Randomized Search. Random Forest liefert bereits mit den
% Standardeinstellungen brauchbare Ergebnisse und reagiert auf Hyperparameter deutlich weniger
% empfindlich als etwa Support Vector Machines.\footcite[Abschn. 3.1]{Probst2019} Getunt wurde
% dennoch, um alle drei Verfahren unter denselben Bedingungen zu vergleichen.

% ERGÄNZEN: die Konfigurationstabelle selbst -- das Exposé liefert nur den Rahmen.
% ---- ENDE ÜBERNAHME ----

Aus einer Verfahrensfamilie wird ein Modell erst durch die Wahl seiner
Hyperparameter, und diese Wahl darf die Testdaten nicht kennen. Gesucht wird
deshalb in einer inneren Kreuzvalidierung über vier Folds, die wie die äußere
nach Stadtteilen gruppiert: Ein Stadtteil steht mit allen 132 Monaten entweder
im inneren Training oder in der inneren Validierung, nie in beidem. Je Fold
entsteht ein eigener Parametersatz. Suchraum und Ziehungszahl sind allein an
diesen inneren Folds festgelegt --- kein Testfold geht in die Festlegung ein.

Gezogen wird zufällig statt gerastert, und zwar hundertmal je Suchlauf. Für
die Wahrscheinlichkeit, mit $T$ Zufallsziehungen mindestens einmal in einen
Zielbereich vom relativen Volumen $v/V$ zu treffen, gilt
$P = 1-(1-v/V)^{T}$.\footcite[296]{Bergstra2012} Angewandt auf ein Zielvolumen
von fünf Prozent ergibt die Formel 99,4 Prozent bei hundert Ziehungen
gegenüber 92,3 Prozent bei fünfzig; beide Werte sind eigene Rechnung und keine
Angabe der Quelle. Der Ausdruck enthält die Dimension des Suchraums nicht:
Ridge mit einem und XGBoost mit sieben Hyperparametern erhalten dasselbe
Budget, ohne dass dem höherdimensionalen Verfahren daraus ein Nachteil
entsteht.

Getunt wird für alle drei Verfahren, auch wo wenig zu gewinnen ist. Random
Forest liefert bereits mit den Vorgabewerten brauchbare Ergebnisse und
reagiert auf Hyperparameter deutlich weniger empfindlich als etwa Support
Vector Machines.\footcite[13]{Probst2019} Ihn ungetunt zu lassen wäre
trotzdem eine ungleiche Behandlung: Der Vergleich soll die Verfahren trennen
und nicht den Aufwand, der in sie gesteckt wurde. Die Räume sind für beide
Zielgrößen dieselben, weil dort die Verlustfunktion wechselt und nicht der
Mechanismus; allein der Varianzexponent der Tweedie-Verlustfunktion entfällt
in der Klassifikation, wo er ohne Wirkung bliebe.

Tabelle~\ref{tab:suchraeume} nennt die Räume und die Werte, die die Suche
darin gefunden hat. Vier Parameter enden an natürlichen Grenzen: alle Merkmale
je Split, eine Beobachtung je Blatt, der ganze Datensatz --- dahinter
existiert nichts, was sich noch öffnen ließe. Dass der Random Forest im
Mengenstrang dort anschlägt und in vier von fünf Folds sämtliche Merkmale je
Split zulässt, ist deshalb keine Frage der Sucheinstellung, sondern eine
Aussage über das Verfahren: Es verlangt maximale Flexibilität.

\begin{table}[htbp]
\tabformat
\caption{Suchräume und die je Fold gewählten Werte, Spanne über die fünf Folds}
\label{tab:suchraeume}
\begin{tabularx}{\textwidth}{@{}L{2.0cm} L{2.9cm} Y L{2.9cm} L{2.5cm}@{}}
\toprule
\textbf{Verfahren} & \textbf{Hyperparameter} & \textbf{Suchraum} & \textbf{gewählt, Menge} & \textbf{gewählt, Struktur} \\
\midrule
Ridge & Strafterm & $10^{-5}$ bis $10^{5}$, logarithmisch & 42,2 bis 1.428 & entfällt \\
\addlinespace
Random Forest & Zahl der Bäume & 200 bis 1.000 & 240 bis 898 & 236 bis 995 \\
 & Baumtiefe & 8, 12, 16, 24, 32, 48, unbegrenzt & 12 bis 32, einmal unbegrenzt & 8 bis 48 \\
 & Beobachtungen je Blatt & 1 bis 20 & 9 bis 19 & 5 bis 18 \\
 & Merkmale je Split & Wurzel, Logarithmus, 0,3, 0,5, 1,0 & Wurzel bis 1,0 & Wurzel bis 0,5 \\
\addlinespace
XGBoost & Zahl der Bäume & 200 bis 1.000 & 406 bis 845 & 201 bis 997 \\
 & Lernrate & 0,01 bis 0,3, logarithmisch & 0,014 bis 0,222 & 0,012 bis 0,124 \\
 & Baumtiefe & 1 bis 14 & 2 bis 14 & 1 bis 9 \\
 & Zeilenanteil je Baum & 0,6 bis 1,0 & 0,69 bis 0,99 & 0,68 bis 1,00 \\
 & Merkmalsanteil je Baum & 0,6 bis 1,0 & 0,63 bis 0,95 & 0,72 bis 0,99 \\
 & Strafterm & $10^{-4}$ bis $10^{4}$, logarithmisch & 0,0006 bis 3.966 & 1,7 bis 9.333 \\
 & Varianzexponent & 1,01 bis 1,90 & 1,07 bis 1,73 & entfällt \\
\bottomrule
\end{tabularx}
\tabquelle
\end{table}

Bei den gesetzten Grenzen fällt das Bild uneinheitlich aus, und das gehört so
berichtet. Der Strafterm von Ridge bleibt mit 42 bis 1.428 in einer schmalen
Zone weit innerhalb des Raums; der von XGBoost streut im Mengenstrang dagegen
über sieben Zehnerpotenzen, von 0,0006 bis 3.966. An die gesetzten Ränder stößt
die Suche an vier Stellen: Die Baumtiefe von XGBoost erreicht im Mengenstrang
in einem Fold die Obergrenze von vierzehn, der Random Forest im Strukturstrang
in zwei Folds die Untergrenze von acht, und der Strafterm von XGBoost liegt dort
in zwei Folds dicht unter $10^{4}$. Der Raum ist damit an den meisten Stellen
weit genug, um die Wahl nicht zu erzwingen. Eng genug, um sie zu stabilisieren,
ist er nirgends.

\subsection{Modellspezifische Tuning-Befunde}
\label{sec:tuning-befunde}
% ===== BLOCK 2 — vor Zeile 1854, Kapitel 6.4 ==============================
% PARAMETERSENSITIVITAET (vorpruefung/v6_parametersensitivitaet.py)
% Beleg: results/parametersensitivitaet/ · R-4, B-14 · Entwurf der Gegenprobe
% seit 11.08.2026 in 06_RISIKEN.md (Commit 12dff7c)
% ORT IST ENTSCHEIDEND: Das gehoert hierher, zur Suche - nicht nach Kapitel 8.
%
% Zu schreiben: Die Suche liefert je Fold einen eigenen Parametersatz; fuer
% die Schlussbewertung gilt einer davon. Ob diese Wahl das Ergebnis bewegt,
% prueft eine Kreuzprobe - jeder Testfold wird mit jedem der fuenf Saetze
% ausgewertet. KEINE Zeitangabe, kein "vorher", kein "nachtraeglich".
%
%   Spanne der fuenf Saetze IM SELBEN Fold:
%     Ridge 7,48 · Random Forest 10,06 · XGBoost 17,18 RMSE
%     Struktur: Random Forest 0,022 · XGBoost 0,050 Macro-F1
%   Zum Vergleich: der Abstand zwischen den drei Verfahren betraegt 2,2 RMSE.
%
%   Der eigene Satz ist NICHT systematisch besser:
%     Ridge Menge      +0,89 RMSE   ( 5 von 20 fremd besser)
%     XGBoost Menge    +0,20 RMSE   (11 von 20)
%     Random Forest M. -2,24 RMSE   (15 von 20)  <- eigener Satz schlechter
%     Random Forest S. +0,0077 F1   ( 5 von 20)
%     XGBoost Struktur -0,0031 F1   (12 von 20)
%
%   FORMULIERUNGSAUFLAGE: n = 5 Folds. Deskriptiv, KEIN Signifikanztest -
%   einzelne Folds dominieren (Ridge Fold 5 dreht dort das Vorzeichen).
%
%   DEUTUNG, die hierher gehoert: Bei rund 38 effektiven Einheiten ist die
%   innere Kreuzvalidierung zu verrauscht, um Modellkomplexitaet zuverlaessig
%   zu waehlen. Das ist die DRITTE unabhaengige Bestaetigung desselben
%   Musters - neben v3_spezifikation (Flexibilitaet im GLM zerstoert die
%   Prognose) und #49 (groesserer Suchraum verbessert den inneren Wert um
%   0,19 RMSE und verschlechtert XGBoost out-of-sample um 1,89).
%   Vorwaertsverweis auf UF4 in Kapitel 8: dort wird daraus die uebertragbare
%   Regel.
%
%   NICHT VERSCHWEIGEN: Die Baselines haben keine Hyperparameter und sind von
%   dieser Streuung nicht betroffen. Das ist ein eigenstaendiges Argument
%   fuer die Referenz, kein Nebensatz.
% ==========================================================================
% Auffälligkeiten pro Verfahren -- KEINE Wiederholung der Funktionsweise (steht in 2.3).
%
% ===================================================================
% STAND 16.08.2026. Quellen: Decision Log #49, #50, #52;
% docs/07_BEFUNDE.md B-45; results/suchdiagnose/; Abbildung A8.
% ===================================================================
%
% BUDGET -- hergeleitet, nicht gewaehlt. Bergstra und Bengio (2012, S. 296)
% geben die geschlossene Form P = 1 - (1 - v/V)^T fuer die Wahrscheinlichkeit
% an, mit T Zufallsziehungen mindestens einmal in einen Zielbereich vom
% relativen Volumen v/V zu treffen. Fuer v/V = 0,05: T = 50 -> 92,3 %,
% T = 100 -> 99,4 %. Verwendet wird T = 100.
%   Der Ausdruck enthaelt die DIMENSION des Suchraums nicht -- deshalb erhalten
%   Ridge mit einem und XGBoost mit sieben Hyperparametern dasselbe Budget,
%   ohne Nachteil fuer das hoeherdimensionale Verfahren. Das ist die Antwort
%   auf den naheliegenden Einwand.
%   ACHTUNG BEIM ZITIEREN: Die beiden Prozentzahlen sind die EIGENE Anwendung
%   ihrer Formel, nicht ihre Aussage. Die verbreitete Angabe "60 Ziehungen fuer
%   das beste Fuenf-Prozent-Fenster" steht NICHT in dem Papier (im Volltext
%   geprueft am 13.08.2026); dessen eigene Simulation rechnet mit 1 %.
%   Formel als eigene Rechnung kennzeichnen, sonst ist es ein Zitatfehler.
%
% SUCHRAEUME -- vier davon erweitert, weil das Optimum an der Grenze lag:
%   alpha (Ridge)        1e-3..1e3  ->  1e-5..1e5
%   max_depth (RF)       [None,8,12,16,24] -> [8,12,16,24,32,48,None]
%   max_depth (XGBoost)  3..10      ->  1..14
%   reg_lambda (XGBoost) 1e-2..1e2  ->  1e-4..1e4
%   NICHT erweitert: max_features, min_samples_leaf, subsample,
%   colsample_bytree -- deren Grenzen sind NATUERLICH (alle Merkmale, eine
%   Beobachtung je Blatt, der ganze Datensatz), dahinter existiert nichts.
%   Dass der Random Forest dort ans Limit geht, ist ein BEFUND ueber das
%   Verfahren: Es will maximale Flexibilitaet.
%
% DIE BEGRUENDUNG, in einem Satz: In sechs von sieben geprueften Parametern lag
% mindestens ein Fold-Sieger an oder jenseits der alten Grenze. Am deutlichsten
% max_depth von XGBoost im Strukturstrang -- dort waehlten VIER VON FUENF Folds
% den Wert 3, also die Untergrenze; mit geoeffneter Grenze waehlen sie 2, 2, 1.
% Es ist eine Aussage ueber die SUCHE, nicht ueber das Ergebnis.
% Abbildung A8 zeigt das: Lage jedes gewaehlten Werts im eigenen Suchraum.
%
% ES WIRD GENAU EIN LAUF BERICHTET -- Decision Log #52.
% Kein Vorher-Nachher, keine zweite Ergebnisreihe, keine Zahl aus einer
% frueheren Konfiguration. Die Begruendung der Suchraeume steht auf den
% GEWAEHLTEN HYPERPARAMETERN einer Vorabpruefung (Werte an der Grenze), nicht
% auf deren Guetemassen -- deshalb entsteht hier nirgends der Eindruck zweier
% Analysen. Wer beim Schreiben versucht ist, eine Verbesserung oder
% Verschlechterung gegenueber einer aelteren Einstellung zu erwaehnen: nicht
% tun. Die internen Vergleichsbefunde stehen in docs/07_BEFUNDE.md und sind
% ausdruecklich NICHT fuer die Kapitel bestimmt.
%
% ========================================================================
% DARSTELLUNGSREGEL TUNING -- FESTGELEGT 20.08.2026, GILT FUER DIE GANZE ARBEIT
% Die Endkonfiguration (100 Ziehungen, erweiterte Suchraeume) wird als DIE
% Methode beschrieben, nicht als Ergebnis einer Korrektur. Keine Chronologie
% des eigenen Vorgehens ("zuerst 50, dann erweitert"), sondern:
%   1. Methode = 100 Ziehungen ueber vier innere Folds.
%   2. Das halbierte Budget von 50 erscheint als SENSITIVITAETSANALYSE.
%   3. Die Randpruefung erscheint als Eigenschaft des Designs: die gewaehlten
%      Werte liegen im Inneren der Suchraeume, XGBoost brauchte dafuer einen
%      weiteren Bereich als die uebrigen Verfahren.
%   4. IMMER mitschreiben: Suchraeume und Ziehungszahl sind ausschliesslich
%      an den INNEREN Folds festgelegt, die Testfolds gehen nicht ein.
%      Das ist der Satz, der die Frage "ist das nicht nachtraeglich
%      angepasst?" entkraeftet. Von Lukas am 20.08.2026 bestaetigt.
% In 7.1 ist das bereits so umgesetzt. Kapitel 6.3/6.4 muessen folgen.
% ========================================================================
% ZUR EHRLICHKEIT GEHOERT DAGEGEN: Die Verdopplung des Budgets ist fuer VIER
% VON FUENF Verfahren wirkungslos -- gemessen in der Vorabpruefung vom
% 13.08.2026 (results/suchdiagnose/). Random Forest im Mengenstrang gewinnt
% auf vier Nachkommastellen exakt null, Ridge +0,0065, beide Strukturverfahren
% +0,0001; nur XGBoost in der Regression gewinnt spuerbar. So gehoert es in den
% Text: fuer ein Verfahren belegt, fuer vier ohne Wirkung. Das ist eine
% Aussage ueber die SUCHE und beruehrt keine Ergebniszahl.

Die Suche liefert je Fold einen eigenen Parametersatz; für die Schlussbewertung
gilt einer davon. Ob diese Wahl das Ergebnis bewegt, prüft eine Kreuzprobe, die
jeden Testfold mit jedem der fünf Sätze auswertet. Innerhalb desselben Folds
trennen die fünf Sätze bei Ridge 2,54, beim Random Forest 4,91 und bei XGBoost
9,14 RMSE; im Strukturstrang sind es 0,056 und 0,102 Macro-F1. Zwischen dem
besten und dem schwächsten der drei Verfahren liegen dagegen 6,5 RMSE
(Abschnitt~\ref{sec:ergebnisse-menge}).

Der auf dem inneren Fenster gewählte Satz ist dabei nicht systematisch der
bessere. Bei Ridge bringt er im Mittel 0,46 RMSE gegenüber einem fremden
Satz; bei XGBoost kostet er 1,05 und beim Random Forest 1,85, und bei den
Baumverfahren liegt in 14 beziehungsweise 17 von 20 Paarungen ein auf anderen
Daten gefundener Satz vorn. Der auf einem Fold gesuchte Parametersatz ist dort
also nicht nur nicht besser, sondern im Mittel schlechter als ein fremder.
Getestet wird darauf nicht: Die Kreuzprobe hat fünf Folds, und einzelne von
ihnen bestimmen das Vorzeichen --- bei Ridge dreht es ein einziger.

Der Grund liegt nicht in der Suche, sondern in der Menge unabhängiger
Information. Bei rund 39 effektiven Einheiten
(Abschnitt~\ref{sec:rahmen-baselines}) schwankt der innere Gütewert stärker,
als die Unterschiede zwischen den Parametersätzen groß sind; die Suche wählt
dann zu einem Teil Rauschen. Dasselbe Muster zeigt sich unabhängig davon ein
zweites Mal, wenn nicht die Parameter, sondern die Spezifikation des
Referenzmodells flexibler wird (Abschnitt~\ref{sec:ergebnisse-menge}).

Ein halbiertes Budget ändert daran nichts. Eine Suchdiagnose, die jede der hundert Ziehungen mit ihrem inneren Gütewert mitschreibt, zeigt: Die zweite Hälfte der Suche verbessert den besten bis dahin gefundenen Wert bei keinem der fünf Verfahren im Mittel über die Folds um mehr als 0,020, gemessen an der Streuung zwischen den Folds, die für dieselben Verfahren zwischen 0,040 und 0,342 liegt, in jedem Fall ein Bruchteil. Das Budget ist damit nicht empirisch gerechtfertigt, sondern hergeleitet, und die Suche hat sich totgelaufen, lange bevor es ausgeschöpft war. Als Erklärung für einen Verfahrensunterschied scheidet ihre Größe aus.

Wo die Suche landet, ist selbst ein Befund
(Abbildung~\ref{abb:hyperparameter}). Im Strukturstrang wählt XGBoost Baumtiefen
zwischen eins und neun bei Straftermen, die von 1,7 bis 9.333 reichen und in
zwei von fünf Folds dicht an der Obergrenze des Suchraums liegen. Der Random
Forest geht im selben Strang in die Gegenrichtung und lässt seine Bäume in zwei
Folds bis achtundvierzig Ebenen tief wachsen, während er in zwei weiteren die
flachste zugelassene Tiefe von acht wählt. Zwei Verfahren derselben Familie, zwei entgegengesetzte
Antworten auf dieselben 24 Trainingsstadtteile.

\begin{figure}[H]
  \centering
  \includegraphics[width=\textwidth]{Abbildungen/a8_hyperparameter.pdf}
  \caption{Lage der gewählten Werte im jeweiligen Suchraum}
  \label{abb:hyperparameter}
\end{figure}

Von dieser Streuung unberührt bleiben die Referenzmodelle. Sie haben keinen
frei wählbaren Parameter, ihr Wert hängt also nicht davon ab, welcher Fold die
Suche gesehen hat. Das ist kein Nebensatz, sondern ein eigenständiges Argument
für sie: Eine Messlatte, die nicht selbst schwankt, macht den Abstand zu ihr
überhaupt erst lesbar.

% ===================================================
\section{Evaluation}
% Umfang: ca. 11 Seiten -- CRISP-DM Phase 5
%
% =========================================================================
% ERGEBNISINVENTAR -- was am Ende vorliegt und was davon in die Arbeit kommt
% =========================================================================
% Gerechnet werden 400 Einzellaeufe. Sie kommen NICHT in die Arbeit, sondern
% ins Abgabe-Zip; berichtet werden Aggregate.
%
%   Rohdaten in results/          Regression   Klassifikation
%     Einzellaeufe (*_folds.csv)         300              100
%       = Zielgroessen x Verfahren x 10 Wiederholungen x 5 Folds
%     Aggregate (*_mittel.csv)             6                2
%     Tuning-Ergebnisse                   30               10
%     Hold-out                             6                2
%     Baselines je Fold (liegen vor)      20               10
%
%   Berichtete Aussagen          Regression   Klassifikation
%     PRIMAER: gegen Stufe-2-Baseline      6                2
%     sekundaer: paarweise Verfahren       6                1
%     Laufzeitpaare (Training/Inferenz)    6                2
%     Hold-out-Messungen                   6                2
%
% ACHT PRIMAERAUSSAGEN tragen das Kapitel. Jede lautet: "Verfahren X schlaegt
% auf Zielgroesse Y die Stufe-2-Baseline um Z, bei Streuung S ueber zehn
% Wiederholungen." Das ist die Antwort auf Unterfrage 2 (Decision Log #34).
% Jede steht fuer sich - sie haengt nicht davon ab, ob sich die Verfahren
% UNTEREINANDER trennen lassen.
%
% Die sieben paarweisen Vergleiche sind die Rangfolge. Sie kommen nur ins
% Kapitel, wenn der gepaarte Wilcoxon-Test sie hergibt.
%
% EHRLICHKEIT BEI DER TESTZAHL: Insgesamt sind es 15 Signifikanztests - acht
% primaere und sieben paarweise. Die Gesamtzahl gehoert in den Text, sonst
% entsteht der Eindruck, sie sei verschwiegen worden (docs/06_RISIKEN.md, R-10).
%
% MASSGEBLICHE STREUUNG ist std_wiederholungen (ueber die 10
% Wiederholungsmittel), NICHT std_folds ueber alle 50 Laeufe - die 50 sind
% nicht unabhaengig, es sind dieselben 29 Stadtteile in anderer Gruppierung
% (R-5). Beide Spalten stehen in *_mittel.csv.
%
% VIER TABELLEN, DREI ABBILDUNGEN decken alles ab:
%   T1  Aggregate Regression (6 Zeilen)
%   T2  Aggregate Klassifikation (2 Zeilen)
%   T3  Vergleichstabelle mit Wilcoxon-Ergebnissen
%   T4  Hold-out (eine einzige Messung an 6 Stadtteilen, KEIN Mittelwert)
%   A1  Boxplot je Verfahren ueber die 50 Laeufe, je Strang
%   A2  Balken: Verfahren gegen Stufe-2-Baseline -- die Primaeraussage
%   A3  Streudiagramm Laufzeit gegen Prognoseguete -- beantwortet UF3 und UF4
%
% Alle Abbildungen erzeugt modelle/m05_abbildungen.py aus den CSV-Dateien,
% nicht von Hand. Sie sind damit reproduzierbar und bei neuen Laeufen in einem
% Befehl aktualisiert.
% =========================================================================
%
% =========================================================================
% BEFUNDLAGE -- Stand 17.08.2026, Durchsicht aller Dateien in results/
% Vollstaendige Tabellen: docs/Guetemasse_Uebersicht_2026-08-17.md
% =========================================================================
%
% B1  IM MENGENSTRANG GEWINNT DIE STUFE-2-BASELINE -- IN BEIDEN VALIDIERUNGEN
%     UND BEI BEIDEN MENGEN-ZIELGROESSEN.
%     ACHTUNG, HAEUFIGE FEHLLESUNG: Diese Aussage gilt AUSSCHLIESSLICH fuer den
%     Mengenstrang. Es gibt zwei verschiedene Stufe-2-Baselines --
%     Poisson-GLM fuer die Regression (beide Zielgroessen, die Rate entsteht
%     aus derselben Vorhersage geteilt durch die Bevoelkerung, #43) und
%     multinomiale logistische Regression fuer die Klassifikation. Im
%     STRUKTURSTRANG faellt das Ergebnis anders aus, siehe B3.
%     Ein Satz wie "die Baseline schlaegt alle Verfahren" ist FALSCH.
%     Kreuzvalidierung anzahl_einsaetze: Poisson-GLM RMSE 33,978 / R2 0,542
%     gegen ridge 36,509 / 0,514, random_forest 35,606 / 0,412, xgboost
%     37,772 / 0,510. Hold-out: GLM RMSE 22,350 / R2 0,801 gegen ridge
%     23,387 / 0,782, xgboost 26,179 / 0,727, random_forest 30,501 / 0,629.
%     Die Rangfolge auf dem Hold-out ist ueber beide Zielgroessen identisch:
%     Poisson-GLM > ridge ~ xgboost > random_forest.
%     Das ist die zentrale Aussage von 7.1 und darf nicht relativiert werden.
%     Quelle: results/regression/baselines_mittel.csv, menge_mittel.csv,
%     holdout.csv.
%
% B2  DIE SPEZIFIKATIONSPRUEFUNG IST DER UNABHAENGIGE BELEG FUER B1 --
%     SIE GEHOERT IN DEN HAUPTTEXT, NICHT IN DEN ANHANG.
%     Erweitert man das GLM um nichtlineare Terme, bricht die Prognose
%     zusammen: linear (12 Terme) R2 0,542 -> quadrate (22) -5,971 ->
%     interaktionen (57) -6,247 -> beides (67) -16,454. Alle 50 Laeufe
%     konvergiert, also kein numerisches Artefakt.
%     ARGUMENTATIVE FUNKTION: Der Befund zeigt OHNE Rueckgriff auf die drei
%     ML-Verfahren, dass Flexibilitaet in diesen Daten aktiv schadet. Damit
%     ist erklaert, WARUM Random Forest und XGBoost das GLM nicht schlagen --
%     es liegt an der Datenstruktur, nicht an Tuning oder Implementierung.
%     Das nimmt der Pruefungsfrage "haben Sie nicht einfach schlecht
%     getunt?" die Spitze. Vgl. auch die Suchdiagnose mit Budget 100
%     (results/suchdiagnose/), die zeigt, dass ein groesseres Budget die
%     Guete nicht bewegt.
%     Quelle: results/spezifikation/spezifikation_mittel.csv.
%
% B3  DER STRUKTURSTRANG IST SCHWACH -- UND DAS IST ZU BERICHTEN.
%     GEGENSTUECK ZU B1: Hier ist die Stufe-2-Baseline die multinomiale
%     logistische Regression, NICHT das Poisson-GLM. Und hier schlagen die
%     Baumverfahren die Baseline in der Kreuzvalidierung sehr wohl --
%     random_forest +0,031 und xgboost +0,035 Macro-F1, beide p = 0,0020.
%     Die Aussage "die Baseline gewinnt immer" gilt also NICHT projektweit.
%     Sie gilt im Mengenstrang; im Strukturstrang gewinnt die Baseline erst
%     auf dem Hold-out, und genau dieser Umschlag ist der Befund.
%     Kreuzvalidierung Macro-F1: Mehrheitsklasse 0,223, Logit 0,297,
%     random_forest 0,328, xgboost 0,332. Hold-out: Logit 0,327,
%     xgboost 0,260, random_forest 0,257 -- die Rangfolge DREHT SICH.
%     Random-Forest-Macro-AUROC faellt von 0,739 (CV) auf 0,521 (Hold-out),
%     also praktisch Zufall.
%     KEIN Verfahren erreicht die Accuracy der Mehrheitsklasse (0,806 CV /
%     0,711 Hold-out). Das ist Folge der Metrikwahl, nicht ein Fehler --
%     muss aber offen dastehen, sonst wirkt es verschwiegen.
%     Quelle: results/klassifikation/baselines_klasse_mittel.csv,
%     struktur_mittel.csv, holdout.csv.
%
% B4  IM STRUKTURSTRANG TRAEGT KEINE MERKMALSGRUPPE.
%     Groesster Ablationseffekt ueberhaupt: 0,024 Macro-F1 (Logit,
%     groessenkontrolle). Bei random_forest liegt der Kriminalitaetseffekt
%     bei 0,000 mit 5 von 10 Wiederholungen -- exakt Muenzwurf.
%     Zum Vergleich Mengenstrang: kriminalitaetsbezogen weglassen kostet
%     +24,274 RMSE bei 10 von 10 Wiederholungen.
%     Das ist ein eigenstaendiges Ergebnis zu UF1: Die Einsatzart laesst
%     sich aus diesen Merkmalen im Kern nicht vorhersagen.
%     Quelle: results/shap/ablation_faktorgruppen_mittel.csv.
%
% B10 BERICHTSUMFANG: DER MENGENSTRANG WIRD AUF EINE ZIELGROESSE VERDICHTET.
%     Entscheidung vom 17.08.2026. GERECHNET bleiben beide Zielgroessen -
%     results/ und der Code sind unveraendert, es wird KEIN Lauf wiederholt.
%     Verdichtet wird allein die BERICHTERSTATTUNG.
%
%     WAS BLEIBT UND WAS GEHT
%       Fliesstext, Tabellen im Hauptteil, Signifikanztests   nur anzahl_einsaetze
%       einsaetze_je_1000_ew                                  ein Absatz + Anhangstabelle
%       ABBILDUNGEN                                           BEIDE Zielgroessen
%
%     WARUM anzahl_einsaetze UND NICHT DIE RATE
%     1  Kapitel 1.2 nennt woertlich "die Einsatzhaeufigkeit als Anzahl
%        Einsaetze pro Stadtteil und Monat". Die Rate kommt in der
%        Forschungsfrage NICHT vor. Streicht man sie, bleibt die Fragestellung
%        unangetastet; streicht man die Anzahl, muss die FF umgeschrieben werden.
%     2  Die Ratenmodellierung geht NICHT verloren. Seit #43 modellieren alle
%        Verfahren die Rate und multiplizieren zurueck; das Poisson-GLM traegt
%        log(Bevoelkerung) als Offset. Der methodische Beitrag - die
%        Expositionsbehandlung - bleibt vollstaendig und ist ueber
%        results/shap/ablation_exposition.csv weiter belegt (22 bis 29 RMSE).
%     3  Bei der Rate sind die Werte nicht rangierbar: R2-Streuung ueber die
%        Folds bis 1,247 bei einem Mittelwert von 0,241 (random_forest).
%        Bei anzahl_einsaetze maximal 0,600 bei 0,412.
%     4  Nennerartefakt (R-19): Bevoelkerung steht im Nenner der Rate UND des
%        Kriminalitaetsindex. Die Within-Korrelation faellt nach Kontrolle von
%        +0,644 auf +0,230. Mit der Rate verschwindet diese offene Flanke.
%     5  R2 war bei der Rate ohnehin nur nachrichtlich - sie streut zwischen
%        den Stadtteilen um Faktor 32 (06_RISIKEN.md, Abschnitt 6, Punkt 2).
%
%     WAS DER SCHNITT KOSTET - GEPRUEFT, NICHTS
%     Alle sechs Rate-Tests stimmen in Richtung UND Signifikanzmuster mit den
%     Anzahl-Tests ueberein. Kein sekundaerer Vergleich wird bei der Rate
%     signifikant (nach Holm p >= 0,117). Es gibt keine Aussage, die nur die
%     Rate traegt. Einziger Unterschied: die staerkste Signifikanz liegt bei
%     der Rate (ridge gegen GLM p = 0,0039 statt 0,0098) - gleiche Richtung,
%     eine Nachkommastelle.
%
%     ZAHLEN NACH DEM SCHNITT
%       Signifikanztests gesamt (beide Straenge)   15  ->   9
%       davon primaer im Mengenstrang               6  ->   3
%       Zeilen menge_mittel / baselines_mittel      6/4 ->  3/2
%       Zeilen holdout.csv (Regression)            10  ->   5
%       berichtete Einzellaeufe                   300  -> 150
%
%     ZWEI STELLEN, DIE DABEI RICHTIG BLEIBEN MUESSEN
%     a) LAUFZEIT (UF3). Die Zeiten sind JE ZIELGROESSE gemessen, jede
%        Zielgroesse bekommt einen eigenen Fit. Berichtet wird die Zeile
%        anzahl_einsaetze: ridge 0,0092 s, xgboost 0,9795 s, random_forest
%        7,9227 s Training. KEIN Mittelwert ueber beide Zielgroessen und keine
%        Gesamtlaufzeit des Skripts, die beide enthaelt.
%        CHANCE: Die Rate liefert dieselben Zeiten mit unter 6 % Abweichung
%        (ridge 0,0092 / xgboost 0,9261 / random_forest 7,7869). Das ist eine
%        zweite Messung derselben Groessenordnung und schwaecht R-15 ab. Ein
%        Satz dazu gehoert in 7.3.
%     b) ABBILDUNGEN. Alle zehn trennen nach Zielgroesse und werden NICHT
%        gefiltert. Begruendung: Abbildungen zaehlen nicht zu den 40 bis 60
%        Seiten (Schroeter, 10.08.) und muessen nicht beschrieben werden. Sie
%        tragen die Robustheitspruefung damit kostenlos. Die Bildunterschriften
%        muessen aber sagen, dass die zweite Zielgroesse als Robustheitspruefung
%        mitlaeuft - sonst steht in der Abbildung mehr als im Text, ohne dass
%        der Leser weiss warum.
%        m05_abbildungen.py bleibt unveraendert. Es liest ausschliesslich CSV
%        und trainiert nichts; ein Neulauf koennte keine Zahl aendern - deshalb
%        waere ein Filter zwar risikofrei, ist aber nicht noetig.
%
%     FORMULIERUNG FUER DEN ROBUSTHEITSABSATZ IN 7.1
%     "Die Befunde sind gegenueber der Wahl der Zielgroessenskalierung robust:
%      Richtung und Signifikanzmuster stimmen bei anzahl_einsaetze und
%      einsaetze_je_1000_ew ueberein (Anhang X). Berichtet wird im Folgenden
%      die Anzahl, weil die Expositionsbehandlung bereits im Modell steckt und
%      die Ratenwerte ueber die Folds so stark streuen, dass eine Rangfolge der
%      Verfahren dort nicht belastbar waere."
%     Das ist ein BEFUND. Zwei gleichrangige Tabellen waeren Redundanz - und
%     genau das, was das Gutachten als "umfangreich, aber nicht fokussiert"
%     bemaengelt hat.
%
% B9  DIE OBERGRENZEN DES STRUKTURSTRANGS -- DAMIT WIRD B3 UND B4 ERST LESBAR.
%     Macro-F1 0,33 gegen die 1,0 einer fehlerfreien Vorhersage zu halten ist
%     der falsche Massstab. Zwei Decken begrenzen den Strang, beide entstehen
%     VOR jeder Modellwahl (vorpruefung/v4_decke.py, Entwicklungspanel):
%
%       Mehrheitsklasse (Stufe 1)      0,2233   triviale Baseline
%       Decke B - Stadtteilwissen      0,4572   Modalklasse je Stadtteil bekannt
%       Decke A - Label-Rauschen       0,6404   Klassenwahrsch. exakt bekannt (+/- 0,0163)
%       fehlerfreie Vorhersage         1,0000   nicht erreichbar
%
%     DECKE A entsteht aus der Konstruktion der Zielgroesse. dominante_einsatzart
%     ist der argmax ueber vier Anteile desselben Monats. Zieht man jeden
%     Stadtteil-Monat aus Multinomial(N, p_beobachtet) neu, kippt der argmax in
%     12,5 % der Faelle. Bei 41,6 % der Zeilen liegt der Abstand zwischen Platz
%     eins und zwei unter 0,20, der mittlere Siegeranteil betraegt 0,509.
%
%     DECKE B entsteht aus der Merkmalsstruktur und ist die BINDENDE. 84,6 % der
%     Zeilen tragen die Modalklasse ihres eigenen Stadtteils -- das Label ist
%     fast vollstaendig stadtteilgebunden. Aber von den 29 Entwicklungsstadtteilen
%     haben 25 dieselbe Modalklasse (fehlalarm), 3 technische_hilfe, 1
%     rettung_ems. Stadtteilwissen ist damit arm, und mehr als Stadtteilwissen
%     tragen die zwoelf Praediktoren nicht (siehe Varianzzerlegung, Kap. 8.3).
%
%     AUSSCHOEPFUNG, baselinekorrigiert als (Modell - Mehrheitsklasse) /
%     (Decke - Mehrheitsklasse) -- der Rohquotient waere geschoent, weil der
%     Sockel der Mehrheitsklasse keine Leistung des Modells ist:
%       xgboost        0,3322   46,6 % von Decke B   26,1 % von Decke A
%       random_forest  0,3278   44,7 % von Decke B   25,1 % von Decke A
%
%     ZU SCHREIBEN: Der Abstand zwischen dem besten Verfahren (0,3322) und
%     Decke B (0,4572) ist alles, was Verfahrenswahl und Hyperparametersuche
%     ueberhaupt noch holen koennten -- 0,125 Macro-F1. Das erklaert B4
%     (keine Merkmalsgruppe traegt) und B7 (die Verfahren trennen sich nicht)
%     aus derselben Ursache und macht aus dem schwaechsten Kapitel einen
%     methodischen Befund. Die Decken gehoeren VOR die Ergebnistabelle in 7.2,
%     nicht in die Limitationen -- sonst liest es sich als nachtraegliche
%     Entschuldigung statt als Massstab.
%     Quelle: results/klassifikation/decke.csv · decke_marge.csv ·
%             decke_ausschoepfung.csv · decke.md
%     Hold-out-Gegenstueck in decke_holdout.csv: Decke A 0,6785, Decke B 0,4219.
%
% B5  DIE PUNKTWERTE BEI DER RATE SIND NICHT RANGIERBAR.
%     R2-Streuung ueber die Folds uebersteigt den R2-Mittelwert:
%     random_forest 0,241 +/- 1,247; ridge 0,300 +/- 0,933.
%     Zusaetzlich widersprechen sich RMSE und R2 zwischen den Verfahren.
%     Das ist selbst der Befund und so zu schreiben (Gutachten R6) --
%     keine Rangfolge konstruieren.
%     Quelle: results/regression/menge_mittel.csv.
%
% B6  UEBERANPASSUNG UND AUFWAND -- BELEG FUER UF3/UF4.
%     R2 Training gegen R2 Kreuzvalidierung (anzahl_einsaetze):
%     ridge 0,826/0,514 (Differenz 0,313), xgboost 0,983/0,510 (0,473),
%     random_forest 0,984/0,412 (0,572).
%     Random Forest braucht 7,92 s Training gegen 0,009 s bei ridge --
%     Faktor 860 -- und generalisiert dabei am schlechtesten.
%     Sauberes Effizienzargument, gehoert nach 7.3.
%     Quelle: results/regression/menge_mittel.csv, klassifikation/
%     struktur_mittel.csv.
%
% B7  TRENNSCHAERFE -- WARUM DIE SEKUNDAERVERGLEICHE NICHTS ZEIGEN.
%     Alle paarweisen Verfahrensvergleiche nach Holm p >= 0,117.
%     Die zugehoerigen Effektstaerken liegen bei d = 0,21 bis 0,87, die
%     Trennschaerfe des gepaarten Wilcoxon mit n = 10 damit bei 10 bis 68 %.
%     FORMULIERUNGSAUFLAGE: "Die Verfahren unterscheiden sich nicht" ist
%     nicht gezeigt, sondern NICHT GEMESSEN. Der Unterschied gehoert in den
%     Text, sonst ist es eine Ueberinterpretation eines Nullbefunds.
%     Ergaenzend: p = 0,001953 im Strukturstrang ist das kleinste bei n = 10
%     ueberhaupt erreichbare Wilcoxon-p (10 von 10 gleiches Vorzeichen).
%     Quelle: results/regression/vergleich.csv, klassifikation/vergleich.csv.
%
% B8  HOLD-OUT IST LEICHTER ALS DIE KREUZVALIDIERUNG -- BENENNEN.
%     Extrapolationsanteil 7,6 % (Hold-out) gegen 33,7 % (CV). Deshalb
%     liegen alle Hold-out-Werte ueber den CV-Werten. Ohne diesen Hinweis
%     wirkt die Schlussbewertung geschoent.
%     Quelle: results/regression/holdout.csv, menge_mittel.csv.
% =========================================================================
\subsection{Ergebnisse Einsatzhäufigkeit (Regression)}
\label{sec:ergebnisse-menge}

Keines der drei Verfahren übertrifft die Stufe-2-Baseline nachweisbar --- und keines bleibt nachweisbar hinter ihr zurück. Trennbar ist dagegen der Abstand zwischen den Verfahren: Ridge schlägt beide Baumensembles in beiden Zielgrößen. In der Schlussbewertung kehrt sich diese Rangfolge um.
 
Grundlage sind die 3.960 Zeilen des Entwicklungspanels und die stadtteilweise Kreuzvalidierung aus Abschnitt 5.4: dieselben zwölf Merkmale, dieselben Zeilen, dieselben Folds für jedes Verfahren. Jede Konfiguration wird fünfzigmal gerechnet, fünf Folds in zehn Wiederholungen. Diese fünfzig Läufe sind keine unabhängigen Erhebungen, sondern dieselben 30 Stadtteile in wechselnder Gruppierung: Streuung und Test beruhen deshalb auf den zehn Wiederholungsmitteln, nicht auf den fünfzig Läufen. Hauptmaß ist der RMSE, Referenzpunkt das Poisson-GLM der Stufe 2.
 
\begin{table}[htbp]
    \tabformat
    \caption{Prognosegüte der Einsatzhäufigkeit}
    \label{tab:menge}
    \begin{tabularx}{\textwidth}{@{}Y r r r r r r@{}}
        \toprule
        & \multicolumn{3}{c}{\textbf{Kreuzvalidierung}}
        & \multicolumn{3}{c}{\textbf{Schlussbewertung}} \\
        \cmidrule(lr){2-4} \cmidrule(l){5-7}
        \textbf{Modell} & \textbf{RMSE} & \textbf{MAE} & \textbf{R\textsuperscript{2}}
                        & \textbf{RMSE} & \textbf{MAE} & \textbf{R\textsuperscript{2}} \\
        \midrule
        Poisson-GLM (Stufe 2)      & 32,99 & 21,44 & 0,607  & 28,46 & 20,10 & 0,625 \\
        Ridge Regression           & 30,93 & 20,36 & 0,655  & 30,30 & 21,95 & 0,575 \\
        Random Forest              & 34,96 & 24,07 & 0,498  & 27,46 & 21,13 & 0,651 \\
        XGBoost                    & 37,39 & 25,13 & 0,486  & 25,83 & 19,86 & 0,691 \\
        \midrule
        Gesamtmittelwert (Stufe 1) & 68,09 & 50,34 & $-0{,}865$ & 49,32 & 38,70 & $-0{,}126$ \\
        \bottomrule
    \end{tabularx}
    \tabquelle
\end{table}
 
Zwischen dem besten und dem schwächsten Verfahren liegen 6,5 RMSE, und RMSE und Bestimmtheitsmaß führen hier zur selben Reihenfolge. Aus den Punktwerten allein ist dennoch nichts abzulesen: Die Streuung über die Wiederholungen reicht von 1,6 beim Ridge bis 5,3 beim Random Forest, die Streuung über die Folds von 19,9 bis 23,8 RMSE und übersteigt jeden Verfahrensunterschied um ein Vielfaches. Sie hebt sich weitgehend auf, weil jedes Verfahren in jedem Lauf denselben Testfold sieht --- die Streuung der gepaarten Differenz gegen die Baseline beträgt noch 3,9 bis 7,9. Geprüft werden diese Differenzen mit dem Vorzeichen-Rang-Test\footcite[80]{Wilcoxon1945}, der für den Vergleich mehrerer Lernverfahren auf denselben Daten als Standardverfahren gilt.\footcite[]{Demsar2006}
 
\begin{figure}[H]
    \centering
    \includegraphics[width=\textwidth]{Abbildungen/a11_differenzen.pdf}
    \caption{Gepaarte Differenzen gegen die Stufe-2-Baseline}
    \label{abb:differenzen}
\end{figure}
 
Ridge liegt im Mittel um 2,06 RMSE vor der Baseline und gewinnt acht von zehn Wiederholungen; das Konfidenzintervall schließt die Null jedoch ein (p = 0,275). Random Forest bleibt um 1,98 RMSE zurück (sechs von zehn, p = 0,922), XGBoost um 4,40 (vier von zehn, p = 0,275). Keiner dieser drei Abstände ist trennbar: Ein Prognosegewinn der komplexeren Verfahren lässt sich ebenso wenig belegen wie ein Rückstand.
 
Davon zu trennen ist die Frage, ob sich die drei Verfahren untereinander unterscheiden. Sie bilden gemeinsam mit den Vergleichen für die Einsatzrate eine sekundäre Familie von sechs Tests, deren p-Werte nach dem Verfahren von Holm korrigiert sind.\footcite[]{Holm1979} Hier wird das Bild deutlich: \textbf{Ridge schlägt beide Baumensembles} --- XGBoost in allen zehn Wiederholungen, den Random Forest in neun von zehn. Nur die beiden Baumverfahren untereinander bleiben ununterscheidbar. Damit trägt der Vergleich zwischen den Verfahren eine Aussage, die der Vergleich zur Baseline nicht hergibt.
 
 
\begin{table}[htbp]
\tabformat
\caption{Gepaarter Wilcoxon-Test auf den zehn Wiederholungsmitteln, RMSE}
\label{tab:regression-tests}
\begin{tabularx}{\textwidth}{@{}Y r r r@{}}
\toprule
\textbf{Paarung} & \textbf{Differenz [95-\%-KI]} & \textbf{gewonnen} & \textbf{$p$} \\
\midrule
Ridge gegen Stufe 2         & $+2{,}06$ [$-0{,}76$; $+4{,}88$] & 8/10 & 0,275 \\
Random Forest gegen Stufe 2 & $-1{,}98$ [$-7{,}66$; $+3{,}71$] & 6/10 & 0,922 \\
XGBoost gegen Stufe 2       & $-4{,}40$ [$-9{,}50$; $+0{,}69$] & 4/10 & 0,275 \\
\addlinespace
Ridge gegen Random Forest   & $+4{,}04$ [$+0{,}58$; $+7{,}49$] & 9/10 & \textbf{0,023} \\
Ridge gegen XGBoost         & $+6{,}46$ [$+3{,}53$; $+9{,}40$] & 10/10 & \textbf{0,012} \\
Random Forest gegen XGBoost & $+2{,}43$ [$-0{,}79$; $+5{,}64$] & 7/10 & 0,211 \\
\bottomrule
\end{tabularx}
\tabquelle
\end{table}
 
Das Bild ist gegenüber der Wahl der Zielgröße invariant. Für die Einsatzrate ist ebenfalls kein Verfahren von der Baseline zu trennen (p = 0,084 für Ridge, 0,846 für Random Forest, 0,492 für XGBoost), und auch dort schlägt Ridge XGBoost in allen zehn Wiederholungen nach Korrektur signifikant; der Abstand zum Random Forest verfehlt sie knapp (p = 0,059). Das ist mehr als eine Wiederholung: Bei der Rate wird auf derselben Skala bewertet, auf der alle drei Verfahren angepasst werden, während das Poisson-GLM dort auf Zähldaten schätzt --- der Rückstand besteht also auch dort, wo die Ausrichtung zwischen Verlustfunktion und Gütemaß den Vergleichsverfahren zugutekommt. Berichtet wird gleichwohl die absolute Anzahl, weil Kapitel 1.2 sie wörtlich nennt und die Bestimmtheitsmaße der Rate über die Folds stärker streuen als ihr eigener Mittelwert.
 
Die Schlussbewertung in Tabelle~\ref{tab:menge} bestätigt diese Rangfolge \emph{nicht} --- sie kehrt sie um. Auf den sechs zurückgehaltenen Stadtteilen erreicht XGBoost mit 25,83 RMSE den geringsten Fehler, gefolgt vom Random Forest mit 27,46; beide liegen vor dem Poisson-GLM (28,46), während Ridge mit 30,30 zurückfällt. Für die Einsatzrate bleibt dagegen die Baseline vorn (0,96 gegen 1,06 beim Ridge). In der Kreuzvalidierung liegen die Werte der Rate bei 1,63 RMSE für Ridge, 1,79 für den Random Forest und 1,84 für XGBoost gegen 1,74 der Stufe-2-Baseline. Diese Umkehr ist eine einzelne Messung an sechs Einheiten ohne Streuung und ohne Test; ihre Absolutwerte sind mit denen der Kreuzvalidierung nicht zu verrechnen (Abschnitt~\ref{sec:belastbarkeit}). Berichtet werden beide Auswertungen, und die Umkehr selbst ist der Befund: Welches Verfahren vorn liegt, hängt bei dieser Zahl unabhängiger Einheiten daran, welche Stadtteile im Test stehen.
 
Dass die flexibleren Verfahren die Baseline nicht übertreffen, liegt dabei nicht an einer zu starr gewählten Baseline: Wird sie selbst schrittweise flexibler spezifiziert, verschlechtert sich ihre Güte um ein Vielfaches des Verfahrensunterschieds (Abschnitt~\ref{sec:belastbarkeit}).

\subsection{Ergebnisse Einsatzart (Klassifikation)}
\label{sec:ergebnisse-struktur}
 
Von den 4.751 Stadtteil-Monaten werden 79,6 Prozent von Fehlalarmen dominiert, 15,9 Prozent von technischer Hilfe, 3,1 Prozent von Rettungseinsätzen und 1,5 Prozent von Bränden. Ein Modell, das ausnahmslos Fehlalarm vorhersagt, erreicht damit eine Accuracy von 0,795, ohne aus den Merkmalen etwas gelernt zu haben.\footcite[155]{Manning2008} Hauptmaß ist deshalb Macro-F1, das jeder der vier Klassen dasselbe Gewicht gibt statt jeder einzelnen Zeile\footcite[280-281]{Manning2008}, ergänzt um die Macro-AUROC, die weder von der Entscheidungsschwelle noch von den Klassenanteilen abhängt.\footcite[176]{Hand2001} Accuracy wird mitgeführt, trägt hier aber keine Aussage über die Güte --- sie belohnt das Aufgeben der drei kleinen Klassen. Referenz ist das multinomiale Logit der Stufe 2 aus Abschnitt 5.4, Nullmarke die Mehrheitsklasse.
 
Welche Güte überhaupt erreichbar ist, begrenzen zwei vor jeder Modellwahl
gemessene Decken (Abbildung~\ref{abb:decken}, Tabelle~\ref{tab:struktur}). Die
Unschärfe des Labels --- die dominante Einsatzart wechselt bei erneuter Ziehung
in 13,6 Prozent der Stadtteil-Monate --- setzt die obere; enger bindet das
verfügbare Wissen: Weil die Merkmale bis auf die Saisonterme nahezu stadtteilgebunden sind (Abschnitt 4.2), kann ein Modell im Kern nicht mehr als jedem Stadtteil seine häufigste Klasse zuweisen, und 27 der 30 Entwicklungsstadtteile haben dieselbe. Decke B ist damit eine Näherung von unten; die wahre Grenze liegt geringfügig darüber, weil Kriminalitätsindex und Saisonterme innerhalb eines Stadtteils variieren. Zwischen der Mehrheitsklasse und dieser
Obergrenze liegen 0,193 Macro-F1 --- der gesamte Spielraum für Merkmals- und
Verfahrenswahl.
 
\begin{figure}[htbp]
\centering
\includegraphics[width=\textwidth]{Abbildungen/a12_decken.pdf}
\caption{Erreichbare Güte im Strukturstrang}
\label{abb:decken}
\end{figure}
 
\begin{table}[htbp]
    \tabformat
    \caption{Güte im Strukturstrang und erreichbare Obergrenzen}
    \label{tab:struktur}
    \begin{tabularx}{\textwidth}{@{}Y r r r r@{}}
        \toprule
        & \multicolumn{2}{c}{\textbf{Kreuzvalidierung}}
        & \multicolumn{2}{c}{\textbf{Schlussbewertung}} \\
        \cmidrule(lr){2-3} \cmidrule(l){4-5}
        \textbf{Modell} & \textbf{Macro-F1} & \textbf{Macro-AUROC} & \textbf{Macro-F1} & \textbf{Macro-AUROC} \\
        \midrule
        Mehrheitsklasse (Stufe 1)     & 0,221             & --    & 0,222 & --    \\
        Multinomiales Logit (Stufe 2) & 0,301 $\pm$ 0,015 & 0,725 & 0,350 & 0,797 \\
        Random Forest                 & 0,318 $\pm$ 0,014 & 0,705 & 0,282 & 0,762 \\
        XGBoost                       & 0,301 $\pm$ 0,015 & 0,665 & 0,291 & 0,818 \\
        \midrule
        Decke B (Stadtteilwissen)             & 0,415 & -- & 0,423 & -- \\
        Decke A (Label-Rauschen)              & 0,641 & -- & 0,677 & -- \\
        \bottomrule
    \end{tabularx}
    \tabquelle
\end{table}
 
In der Kreuzvalidierung übertrifft nur der Random Forest die Baseline: um 0,017 Macro-F1, in neun von zehn Wiederholungen und mit $p = 0{,}004$. XGBoost erreicht mit $-0{,}000$ bei $p = 1{,}000$ exakt ihr Niveau. Der sekundäre Vergleich der beiden Baumverfahren untereinander fällt dagegen zugunsten des Random Forest aus ($+0{,}018$, acht von zehn, $p = 0{,}010$). Auch die Macro-AUROC spricht für die Referenz: Sie liegt mit 0,725 über beiden Baumverfahren (0,705 und 0,665).
 
Gemessen am Erreichbaren bleibt der Vorsprung klein: Der Random Forest schöpft
baselinekorrigiert 50,2 Prozent von Decke B aus, XGBoost 41,0; bis zur Decke
fehlen dem besseren noch 0,096 Macro-F1. Die Accuracy der Mehrheitsklasse
erreicht dabei keines von ihnen, wer alle vier Klassen gleich gewichtet, gibt
Accuracy notwendigerweise auf.
 
\begin{table}[htbp]
\tabformat
\caption{Gepaarter Wilcoxon-Test auf den zehn Wiederholungsmitteln, Macro-F1}
\label{tab:klassifikation-tests}
\begin{tabularx}{\textwidth}{@{}Y r r r@{}}
\toprule
\textbf{Paarung} & \textbf{Differenz [95-\%-KI]} & \textbf{gewonnen} & \textbf{$p$} \\
\midrule
Random Forest gegen Stufe 2 & $+0{,}017$ [$+0{,}005$; $+0{,}029$] &  9/10 & \textbf{0,004} \\
XGBoost gegen Stufe 2       & $-0{,}000$ [$-0{,}015$; $+0{,}015$] &  6/10 & 1,000 \\
Random Forest gegen XGBoost & $+0{,}018$ [$+0{,}007$; $+0{,}028$] &  8/10 & \textbf{0,010} \\
\bottomrule
\end{tabularx}
\tabquelle
\end{table}
 
Auf den sechs unberührten Stadtteilen wird der Abstand größer statt kleiner: Das multinomiale Logit liegt dort mit 0,350 deutlich vor beiden Baumverfahren, obwohl es in der Kreuzvalidierung nur 0,301 erreicht. Zwei Erklärungen kommen dafür in Betracht, und dieses Design kann sie nicht trennen. Für Überanpassung spricht, dass die Baumverfahren dort auf 30 statt 24 Stadtteilen trainieren, also mit mehr Daten schlechter werden, während die Baseline besser wird. Dagegen spricht, dass die zurückgehaltenen Stadtteile eine andere Klassenverteilung tragen: Der Anteil brandbestimmter Monate liegt dort bei 4,7 statt 0,8 Prozent, weil die Stratifizierung nach der seltensten Klasse den auf diesem Kriterium rangersten Stadtteil zwangsläufig in die gesperrte Gruppe legt (Abschnitt 5.4). Beide Befunde werden berichtet, keiner als entschieden dargestellt.
 
\begin{figure}[htbp]
\centering
\includegraphics[width=\textwidth]{Abbildungen/a13_umschlag.pdf}
\caption{Kreuzvalidierung und Schlussbewertung im Vergleich}
\label{abb:umschlag}
\end{figure}
 
Sichtbar macht der Hold-out dagegen etwas anderes: Die Rangfolge der beiden Gütemaße fällt dort auseinander. Im Macro-F1 liegt das Logit mit 0,350 vorn, in der Macro-AUROC dagegen XGBoost mit 0,818 vor dem Logit (0,797) und dem Random Forest (0,762). Ein Modell kann die Klassen also gut trennen und sie dennoch schlecht zuordnen --- die AUROC misst die Rangordnung der Wahrscheinlichkeiten, das Macro-F1 die tatsächliche Zuweisung. Hinzu kommt die dünne Besetzung der seltensten Klasse: Je Fold liegen im Mittel 6,6 brandbestimmte Stadtteil-Monate im Test, weshalb Macro-F1 über die Folds mit 0,057 beim Random Forest und 0,077 bei XGBoost streut gegenüber 0,014 und 0,015 über die Wiederholungsmittel. Der Mittelwert über fünfzig Läufe verdeckt diese Instabilität; für eine Anwendung, in der gerade die seltene Klasse zählt, ist sie die praktisch bedeutsamere Information.

Getragen wird dieses Macro-F1 von den beiden häufigen Klassen: Der Random Forest erreicht 0,860 für Fehlalarm und 0,373 für technische Hilfe, aber 0,040 für Rettungseinsätze und 0,000 für Brände, und die seltenste Klasse sagt er in 49 der fünfzig Läufe kein einziges Mal vorher; bei XGBoost fällt das Bild ähnlich aus. Ein Viertel des Leitmaßes ruht damit auf einer Klasse, die praktisch kein Modell vergibt.
 
Zwischen den Merkmalen eines Stadtteils und der Art seiner Einsätze besteht damit ein Zusammenhang, den in der Kreuzvalidierung allein der Random Forest nachweisbar besser nutzt als das lineare Modell. Dass er schwach bleibt, liegt an den Daten: Stadtteile unterscheiden sich in der Menge ihrer Einsätze erheblich, in deren Zusammensetzung kaum.

% LEITFRAGE: Haengen Einsatzart und Stadtteilmerkmale zusammen -- koennen die
% Modelle die ART der Einsaetze vorhersagen?
%
% Zielgroesse: dominante_einsatzart -- vier Klassen je Stadtteil x Monat
% (Brand, Rettung/EMS, Technische Hilfe/Gefahr, Fehlalarm/Good Intent).
% Eine echte Klasse, gebildet als argmax ueber die vier NFIRS-Gruppen des
% Monats; kein gesetzter Schwellwert und keine kuenstliche Einteilung einer
% stetigen Groesse.
%
% Guetemasse: Macro-F1 (Hauptmass) und Macro-AUROC. DIE MESSLATTE IST STUFE 2,
% also das MULTINOMIALE LOGIT (Macro-F1 0,297 +/- 0,014 | AUROC 0,705 |
% Accuracy 0,584) -- nicht die Mehrheitsklasse. Die Mehrheitsklasse (Stufe 1,
% Macro-F1 0,223 | Accuracy 0,806) ist nur die Nullmarke.
% ACCURACY IST HIER KEIN HAUPTMASS -- 79 % der Stadtteil-Monate werden von
% Fehlalarm dominiert, ein Modell das immer "Fehlalarm" sagt erreicht 0,806.
%
% ERGEBNIS, Kreuzvalidierung (03_STAND.md 5.2):
%     Mehrheitsklasse (Stufe 1)   Macro-F1 0,223    -     Acc. 0,806
%     Multinom. Logit (Stufe 2)            0,297  0,705        0,584
%     Random Forest                        0,3278 0,739        0,762
%     XGBoost                              0,3322 0,735        0,744
%   Random Forest +0,0306 und XGBoost +0,0350 gegen Stufe 2, beide 10/10
%   Wiederholungen gewonnen, p = 0,002 -- der kleinste bei n = 10 ueberhaupt
%   erreichbare Wilcoxon-p. Die Richtung ist eindeutig, nicht knapp.
%   Random Forest gegen XGBoost: nicht unterscheidbar (p 0,625).
%
% HOLD-OUT -- DIE RANGFOLGE DREHT SICH, und genau das ist der Befund:
%     Mehrheitsklasse  0,208    -     Acc. 0,711
%     Multinom. Logit  0,327  0,756        0,472
%     XGBoost          0,260  0,611        0,553
%     Random Forest    0,257  0,521        0,570
%   Random-Forest-Macro-AUROC faellt von 0,739 auf 0,521, also praktisch
%   Zufall. EINE RANGFOLGE ZWISCHEN LOGIT UND BAUMVERFAHREN IST IM
%   STRUKTURSTRANG NICHT ZULAESSIG.
%   Kein Widerspruch: Der Hold-out-Wert liegt innerhalb der Spannweite der 50
%   Einzellaeufe (RF 0,229-0,404, XGB 0,230-0,421); 14 bzw. 16 % der Folds
%   lagen darunter. Eine Ziehung aus einer breiten Verteilung, keine Anomalie.
%   WAS DER HOLD-OUT SICHTBAR MACHT (B-40): Beide Baumverfahren sagen die
%   seltenste Klasse dort KEIN EINZIGES MAL vorher -- F1 fuer brand ist 0,000,
%   bei Random Forest liegt die AUROC mit 0,173 unter dem Zufall. Das Logit
%   erkennt dieselbe Klasse mit AUROC 0,895 und klassifiziert lediglich
%   zurueckhaltend. Der CV-Mittelwert verdeckt diese Instabilitaet.
%
% KEIN VERFAHREN ERREICHT DIE ACCURACY DER MEHRHEITSKLASSE (0,806 CV / 0,711
% Hold-out). Das ist Folge der Metrikwahl, kein Fehler -- muss aber offen
% dastehen, sonst wirkt es verschwiegen.
%
% DIE DECKEN GEHOEREN VOR DIE ERGEBNISTABELLE (B9), nicht in die Limitationen
% -- sonst liest es sich als nachtraegliche Entschuldigung statt als Massstab.
% Beide entstehen VOR jeder Modellwahl (vorpruefung/v4_decke.py):
%     Mehrheitsklasse (Stufe 1)        0,2233
%     Decke B, Stadtteilwissen         0,4572   Modalklasse je Stadtteil bekannt
%     Decke A, Label-Rauschen          0,6404   Klassenwahrsch. exakt bekannt
%     fehlerfreie Vorhersage           1,0000   nicht erreichbar
%   Decke A: dominante_einsatzart ist der argmax ueber vier Anteile desselben
%   Monats; zieht man jeden Stadtteil-Monat aus Multinomial(N, p) neu, kippt
%   der argmax in 12,5 % der Faelle. Bei 41,6 % der Zeilen liegt der Abstand
%   zwischen Platz eins und zwei unter 0,20.
%   Decke B ist die BINDENDE: 84,6 % der Zeilen tragen die Modalklasse ihres
%   eigenen Stadtteils, aber von den 29 Entwicklungsstadtteilen haben 25
%   dieselbe Modalklasse (fehlalarm). Stadtteilwissen ist damit arm.
%   AUSSCHOEPFUNG, baselinekorrigiert -- xgboost 46,6 % von Decke B,
%   random_forest 44,7 %. Der Abstand zwischen dem besten Verfahren (0,3322)
%   und Decke B (0,4572) ist ALLES, was Verfahrenswahl und Hyperparametersuche
%   ueberhaupt noch holen koennten: 0,125 Macro-F1. Das macht aus dem
%   schwaechsten Kapitel einen methodischen Befund.
%
% Verfahren: NUR Random Forest und XGBoost -- zwei statt drei. Ridge hat auf
% einer nominalen Zielgroesse keine Entsprechung, die logistische Regression
% wurde aus Fokusgruenden gestrichen. Freigabe Schroeter 03.08.2026, begruendet
% in 6.2, festgehalten als Decision Log #31. Im Text ist zu benennen, dass der
% Klassifikationsstrang damit bewusst schmaler ist als der Regressionsstrang
% -- nicht stillschweigend zwei Verfahren berichten, wo das Expose drei
% verspricht.
%
% Klassenverteilung, gehoert in die Ergebnisdarstellung:
%   Fehlalarm 79,0 % | Techn. Hilfe 16,3 % | Rettung/EMS 3,1 % | Brand 1,5 %
% Die duenne Besetzung von Brand ist als Limitation zu benennen: Je Fold
% liegen 2 bis 13 brand-dominierte Monate im Test.
%                                                          -> UF2

\subsection{Belastbarkeit der Befunde}
\label{sec:belastbarkeit}

Welche Vergleiche signifikant werden und welche nicht, folgt in beiden Strängen der Effektstärke --- und diese Ordnung erklärt zugleich, warum die Referenz von keinem Verfahren zu trennen ist. Im Mengenstrang liegen die Abstände zur Baseline bei Effektstärken von $d = 0{,}25$ bis $d = 0{,}62$\footcite[]{Cohen1988}, die Abstände zwischen den Verfahren dagegen bei $d = 0{,}54$ bis $d = 1{,}57$; im Strukturstrang trennt den Random Forest von der Referenz $d = 1{,}02$ und von XGBoost $d = 1{,}16$, während XGBoost und Referenz mit $d = 0{,}01$ praktisch zusammenfallen. Für Werte um $d = 0{,}5$ erreicht der gepaarte Vorzeichen-Rang-Test bei zehn Paaren in der Simulation eine Trennschärfe zwischen 12 und 30 Prozent; ein Unterschied dieser Größe wäre also in sieben bis neun von zehn Fällen unentdeckt geblieben.\footcite[485]{Altman1995} Genau dort liegen sämtliche Abstände zur Referenz --- ihre Nichtsignifikanz ist deshalb keine Aussage über Gleichheit, sondern über die Auflösung dieses Designs. Signifikant werden ausschließlich die Vergleiche mit $d \geq 0{,}84$. Nach oben ist die Skala ebenso begrenzt: Der kleinste bei zehn Paaren erreichbare p-Wert ist 0,002, und er entsteht bereits, wenn alle zehn Wiederholungen dasselbe Vorzeichen tragen --- er markiert die Auflösungsgrenze und nicht besonders starke Evidenz.

Die beiden Stränge bilden dabei getrennte Testfamilien, weil sie verschiedene
Teilfragen beantworten: Ein Zufallstreffer im Mengenstrang macht den
Strukturstrang nicht falsch. Folgenlos ist die Festlegung ohnehin --- würde man
beide zu einer Familie aus sieben Tests zusammenlegen, bliebe jede der sieben
Entscheidungen dieselbe, der Strukturvergleich mit einem korrigierten $p$ von
$0{,}039$.

Dass zehn Paare zur Verfügung stehen und nicht fünfzig, ist selbst eine Festlegung. Die fünfzig Läufe sind dieselben 30 Stadtteile in wechselnder Gruppierung; ihre Trainingsmengen überlappen einander, und ein Test über sie hinweg behandelte abhängige Ziehungen als unabhängige, also Pseudoreplikation.\footcite[187]{Hurlbert1984} Nadeau und Bengio zeigen für genau diese Bauform, dass die daraus geschätzte Streuung die tatsächliche Unsicherheit grob unterschätzt, weil sie allein die Testbeispiele abbildet und nicht die Wahl der Trainingsmenge.\footcite[239-240]{Nadeau2003} Ein Test, der die Korrelation zwischen den Läufen nicht kennt, setzt sie stillschweigend auf null.\footcite[254]{Nadeau2003} Geprüft wird deshalb auf den zehn Wiederholungsmitteln, und die kleine Trennschärfe ist der Preis dafür.

Die Grenze dahinter gilt für beide Stränge und ist die Zahl unabhängiger Einheiten: rund 39 statt 4.752 Beobachtungen (Abschnitt~\ref{sec:rahmen-baselines}). Belastbar ist unter dieser Grenze ein Abstand, der in neun oder zehn von zehn Wiederholungen dasselbe Vorzeichen trägt; nicht belastbar ist eine Rangfolge, deren Abstände kleiner sind als die Streuung zwischen den Folds.

Kreuzvalidierung und Schlussbewertung bilden dabei keine gemeinsame Zahlenreihe: Trainiert wird dort auf 30 statt 24 Stadtteilen, im Strukturstrang trägt die zurückgehaltene Gruppe eine andere Klassenverteilung (Abschnitt~\ref{sec:ergebnisse-struktur}), und vor allem wird die Kreuzvalidierung fünfzigmal gemessen, die Schlussbewertung einmal, an sechs Einheiten, ohne Streuung und ohne Test. Dass die Rangfolge im Mengenstrang kippt und im Strukturstrang deutlicher ausfällt, ist mit diesem Design vereinbar und begründet für sich genommen keine der beiden Reihenfolgen.

Dass die flexibleren Verfahren die Baseline nicht übertreffen, lässt sich im Mengenstrang schließlich unabhängig von ihnen prüfen, indem nicht das Verfahren gewechselt, sondern die Baseline selbst schrittweise flexibler spezifiziert wird --- eine Prüfung auf ausgelassene Nichtlinearität in der Bauart des in Abschnitt~\ref{sec:eignung} eingeführten RESET-Ansatzes. Träfe die Vermutung zu, die Baseline sei zu starr, müsste die Güte mit jeder Erweiterung steigen. Das Gegenteil tritt ein: Der RMSE steigt von 33,0 bei zwölf Termen auf 74,6 mit quadratischen Termen, 257,0 mit Interaktionen und 4.203,3 mit beidem, bei fünfzig von fünfzig konvergierten Läufen. Zwischen dem besten und dem schwächsten der vier verglichenen Modelle liegen 6,5 RMSE, zwischen der sparsamsten und der reichhaltigsten Spezifikation dagegen 4.170,3 (Abbildung~\ref{abb:spezifikation}). Belegt ist das für unpenalisierte Erweiterungen; auf regularisierte und getunte Baumensembles überträgt es sich nicht unmittelbar, für sie gilt der eigene Vergleich aus Abschnitt~\ref{sec:ergebnisse-menge}.Eine zu knapp bemessene Suche scheidet als Erklärung ebenfalls aus (Abschnitt~\ref{sec:tuning-befunde}).

\begin{figure}[H]
    \centering
    \includegraphics[width=\textwidth]{Abbildungen/a3_spezifikation.pdf}
    \caption{Spannweite der Verfahrenswahl gegenüber der Spannweite der
             Spezifikationswahl}
    \label{abb:spezifikation}
\end{figure}

\subsection{Trainings- und Inferenzaufwand}
\label{sec:aufwand}

Der Abstand ist vor der Messung absehbar. Die Ridge Regression löst ein einziges
Gleichungssystem, mit einem Aufwand von $\mathcal{O}(np^{2} + p^{3})$. Bei zwölf
Merkmalen fällt der zweite Term nicht ins Gewicht. Die beiden Ensembles wiederholen ihre
Suche nach der besten Teilung für jeden Baum, ihr Aufwand wächst deshalb proportional
zur Baumzahl.\footcite[790]{Chen2016} Der Unterschied zwischen den Verfahren ist damit
kein Unterschied in der Ordnung, sondern ein konstanter Faktor und dieser Faktor ist
die Zahl der Bäume.

Ridge passt einen Fold in 0,008 Sekunden an, XGBoost braucht das 296-Fache, Random Forest das 328-Fache. Das ist der einzige Vergleich dieser Arbeit, in dem sich die Verfahren um Größenordnungen trennen --- und damit der einzige, in dem eine Rangfolge belastbar ist. Alle Zeiten stammen allerdings aus einem Durchgang auf derselben Maschine ohne Nebenlast. Je Konfiguration liegt eine einzelne Messung vor und keine Wiederholungsmessung; die Werte sind als Größenordnung zu lesen, nicht als exakte Kennzahl.
 
\begin{table}[htbp]
    \tabformat
    \caption{Trainings- und Inferenzzeit je Fold, einkernig gemessen}
    \label{tab:aufwand_zeit}
    \begin{tabularx}{\textwidth}{@{}Y r r r@{}}
        \toprule
        & \textbf{Ridge} & \textbf{XGBoost} & \textbf{Random Forest} \\
        \midrule
        Training, Einsatzhäufigkeit & 0,008\,s & 2,40\,s  & 2,66\,s \\
        Training, Rate              & 0,008\,s & 2,39\,s  & 2,66\,s \\
        Training, Einsatzart        & --       & 2,46\,s  & 2,47\,s \\
        \addlinespace
        Inferenz, Einsatzhäufigkeit & 0,003\,s & 0,021\,s & 0,078\,s \\
        \bottomrule
    \end{tabularx}
    \tabquelle
\end{table}
 
Der Abstand folgt daraus, wie viel je Fold angepasst wird: Ridge ein einziges Modell, die beiden Ensembles mehrere hundert Bäume. Untereinander liegen die Ensembles nahe beieinander, weil sich ihre Konfigurationen ausgleichen --- was der Random Forest an Tiefe hat, hat XGBoost an Baumzahl. In der Inferenz trennen sie sich um das Vierfache, allerdings auf unkritischem Niveau: Auch das langsamste Verfahren sagt einen vollständigen Testfold in unter einer Zehntelsekunde vorher.
 
\begin{figure}[H]
    \centering
    \includegraphics[width=\textwidth]{Abbildungen/a4_laufzeit_guete.pdf}
    \caption{Trainingsaufwand gegenüber Prognosegüte}
    \label{abb:laufzeit_guete}
\end{figure}
 
Parallelisierung gleicht diesen Abstand nicht aus. Random Forest baut seine Bäume unabhängig voneinander auf\footcite[5--6]{Breiman2001} und lässt sich deshalb ohne Abstimmung verteilen, was bei der Einsatzhäufigkeit einen Gewinn von 2,06 ergibt; Ridge profitiert kaum (1,08), und bei XGBoost liegt der Wert mit 0,84 unter eins --- der Fit über alle vier logischen Kerne dauert länger als der einkernige. Gemessen ist der Effekt, nicht seine Ursache.
 
Der größte Einzelposten ist ohnehin nicht das Training eines Modells, sondern die Suche davor: 25 Suchläufe mit je 100 Ziehungen über vier innere Folds ergeben 10.000 Modellanpassungen und rund 137 Minuten reine Suchzeit, davon 36 Sekunden auf Ridge. Der Aufwandsnachteil der Baumverfahren entsteht damit nicht erst beim finalen Fit, sondern vervielfacht sich über die Modellwahl.
 
Im Aufwand trennen sich die Verfahren eindeutig, um Größenordnungen statt um Prozentpunkte --- und anders als in der Güte deckt sich diese Rangfolge im Mengenstrang mit der Prognoseleistung: Ridge ist zugleich das schnellste und, gegenüber beiden Ensembles, das nachweisbar bessere Verfahren. Es bleibt allerdings auch das teuerste Verfahren hinter einer Referenz zurück, die ohne jeden Hyperparameter auskommt und in Sekundenbruchteilen angepasst ist. Im Strukturstrang, wo Ridge nicht zur Verfügung steht, sind die beiden Ensembles im Aufwand praktisch gleichauf; die Wahl entscheidet dort allein die Güte, und die spricht für den Random Forest.

% Laufzeiten, Tuning-Aufwand.                             -> UF3
%
% ===================================================================
% STAND 16.08.2026. Zahlen: docs/03_STAND.md Abschnitt 5.4.
% Abbildungen A4 (Aufwand gegen Guete) und A9 (Parallelisierung).
% ===================================================================
%
% MESSUMGEBUNG zuerst nennen, sonst sind die Zahlen nicht lesbar:
% Intel Core i5-7300U @ 2,60 GHz, 2 physische / 4 logische Kerne, 7,8 GB RAM,
% Windows 10 Pro, Python 3.14.0. Ein Durchgang, ohne Nebenlast.
%
% TRAININGSZEIT JE FOLD, EINKERNIG (der zwischen den Verfahren vergleichbare
% Wert -- Begruendung fuer einkernig: siehe 6.1):
%       Menge, anzahl_einsaetze   Ridge 0,009 s | XGBoost 0,98 s | RF 7,92 s
%       Menge, Rate               Ridge 0,009 s | XGBoost 0,93 s | RF 7,79 s
%       Struktur                              XGBoost 2,07 s | RF 3,57 s
%       Inferenz, Menge           Ridge 0,003 s | XGBoost 0,011 s | RF 0,183 s
%   => Ridge ist 106-mal schneller als XGBoost und 861-mal schneller als
%      Random Forest. KEINE Faktoren aus einer anderen Konfiguration nennen
%      (#52). DAS IST DIE BELASTBARSTE AUSSAGE DER ARBEIT, weil hier
%      Groessenordnungen liegen und nicht Prozentpunkte -- und weil UF3 gegen
%      R-1 immun ist: Im Aufwand trennen sich die Verfahren eindeutig.
%   Einordnung dazu: Ridge ist bei besserer Guete als XGBoost und nicht
%   unterscheidbarer gegenueber Random Forest 106- bzw. 861-mal schneller.
%
% PARALLELISIERUNGSGEWINN (einkernige Zeit / parallele Zeit):
%       Menge      Ridge 1,33 / 1,23 | RF 2,23 / 2,18 | XGBoost 0,73 / 0,72
%       Struktur                       RF 1,66        | XGBoost 0,77
%   BEI XGBOOST LIEGT DER GEWINN UNTER 1 -- der Fit ueber alle Kerne dauert
%   LAENGER als der einkernige. Auf zwei physischen Kernen uebersteigt der
%   Verwaltungsaufwand der Threads den Gewinn. Random Forest profitiert, weil
%   seine Baeume unabhaengig sind; Ridge hat als geschlossene Loesung nichts zu
%   verteilen (Gewinn um 1 erwartbar, kein Messfehler). Abbildung A9.
%
% EINSCHRAENKUNG, die zwingend dazugehoert (R-15, R-16):
%   - Der Parallelisierungsgewinn ist staerker an diese Maschine gebunden als
%     die Einkern-Zeiten. Zwei physische Kerne mit Hyperthreading sind ein
%     Grenzfall; auf acht echten Kernen fielen die Faktoren anders aus.
%   - Bei einer U-Serie-CPU ist thermische Drosselung ueber eine Laufzeit von
%     mehreren Stunden nicht auszuschliessen.
%   - UF3 beruht auf EINER Zeitmessung je Konfiguration, nicht auf einer
%     Wiederholungsmessung.
%   - XGBoost ist nicht threaddeterministisch: Bei anderer Kernzahl weichen
%     die Vorhersagen ab (B-24). Die berichteten Guetemasse stammen aus dem
%     einkernigen Fit und sind davon unberuehrt -- das ist zu sagen, sonst
%     liest es sich wie ein Reproduzierbarkeitsproblem der Ergebnisse.
%
% ZUM TUNING-AUFWAND: 25 Suchlaeufe (15 Regression, 10 Klassifikation) mal
% 100 Ziehungen mal 4 innere Folds = 10.000 Modellanpassungen, 139 Minuten
% Suchzeit. Das ist die Kennzahl fuer die Komplexitaet des "V" in E-V-A
% (Auflage 10.08.2026, in 01_VORGABEN.md noch als offen vermerkt).

\subsection{Interpretierbarkeit: SHAP-Analyse}

Die Merkmalsgruppe, auf die ein Modell die meiste Aufmerksamkeit legt, ist nicht die, deren Wegfall am meisten kostet. Die erste Unterfrage wird deshalb in zwei Schritten beantwortet: Die Attribution misst, wie ein angepasstes Modell seine Aufmerksamkeit über die Merkmale verteilt, die Ablation, was fehlt, wenn eine Faktorgruppe weggelassen und das Modell ohne sie neu angepasst wird.\footcite[4]{Hooker2021} Attribuiert wird nur dort, wo ein Modell überhaupt etwas erklärt: im Strukturstrang über TreeSHAP für den Random Forest, der als einziger die Stufe-2-Baseline schlägt\footcite[4]{Lundberg2017}, im Mengenstrang über die standardisierten Koeffizienten des Poisson-GLM, weil dort kein Verfahren die Referenz übertrifft. Ein Vergleich der Attribution zwischen den beiden Baumverfahren ist damit nicht möglich; die Ablation im Anschluss deckt alle vier Modelle ab und trägt die Aussage allein. Beide Größen werden auf die fünf Faktorgruppen verdichtet, weil Beitragsmaße bei korrelierten Prädiktoren einzelne von ihnen bevorzugen\footcite[Abstract]{Strobl2008} und auch TreeSHAP die Abhängigkeit zwischen den Merkmalen nur teilweise abbildet\footcite[2]{Aas2021}; innerhalb eines Strangs sind sie auf die Summe eins normiert und nur dort vergleichbar.
 
\begin{figure}[H]
    \centering
    \includegraphics[width=\textwidth]{Abbildungen/a15_attribution_ablation.pdf}
    \caption{Attribution und Ablation je Faktorgruppe}
    \label{abb:attribution_ablation}
\end{figure}
 
Die Attribution verteilt sich in beiden Strängen auf alle drei Faktorgruppen: In der Menge binden die sozioökonomischen Merkmale 41,4 und die baulichen 34,0 Prozent der Koeffizientenmasse, der Kriminalitätsindex dagegen nur 10,7. In der Struktur verteilt der Random Forest gleichmäßiger --- 29,8 Prozent auf die baulichen, 28,7 auf die sozioökonomischen und 18,5 auf die kriminalitätsbezogenen Merkmale. Das ist eine Aussage darüber, worauf die Modelle ihre Aufmerksamkeit legen --- nicht darüber, was eine Gruppe wert ist.
 
Diese zweite Frage stellt die Ablation, und ihr Ergebnis steht quer zur Attribution --- im Mengenstrang sogar diametral. Dort ist der Kriminalitätsindex unverzichtbar, obwohl er die geringste Aufmerksamkeit bindet: Sein Weglassen kostet 6,10 RMSE, rund 18 Prozent, in zehn von zehn Wiederholungen. Umgekehrt \emph{verbessern} die beiden Gruppen die Prognose durch ihr Weglassen, auf die das Modell drei Viertel seiner Masse legt --- die baulichen Merkmale um 8,04 RMSE in zehn von zehn, die sozioökonomischen um 4,95. In der Struktur ist er praktisch wertlos; dort tragen in allen drei Modellen die sozioökonomischen Merkmale am meisten. Zwei Einschränkungen gehören dazu: Auf der Ablation wird nicht getestet, weil die Testfamilien vorab festliegen und weitere Tests in die Holm-Korrektur eingingen; und bei den Baumverfahren stammen die Hyperparameter aus dem vollen Merkmalssatz, sodass die gemessene Verschlechterung dort auch eine nicht mehr passende Einstellung enthält. Das Poisson-GLM ist davon frei, weil es keinen Hyperparameter besitzt --- der größte gemessene Effekt ist damit zugleich der unbelastete.
 
\begin{table}[htbp]
    \tabformat
    \caption{Güteänderung beim Weglassen einer Faktorgruppe}
    \label{tab:ablation}
    \begin{tabularx}{\textwidth}{@{}Y r r r r@{}}
        \toprule
        \textbf{weggelassen} & \textbf{Poisson-GLM} & \textbf{Logit}
                             & \textbf{Random Forest} & \textbf{XGBoost} \\
                             & \textbf{$\Delta$ RMSE} & \multicolumn{3}{c}{\textbf{$\Delta$ Macro-F1}} \\
        \midrule
        kriminalitätsbezogen & +6,10 (10)  & $-0{,}0048$ (2)  & +0,0020 (7)  & +0,0001 (4)  \\
        sozioökonomisch      & $-4{,}95$ (1)   & +0,0439 (10) & +0,0234 (10) & +0,0381 (10) \\
        Größenkontrolle      & $-5{,}02$ (0)   & +0,0215 (9)  & +0,0147 (10) & +0,0080 (6)  \\
        baulich              & $-8{,}04$ (0)   & $-0{,}0315$ (0)  & $-0{,}0385$ (0)  & $-0{,}0268$ (1)  \\
        Saison               & +0,28 (10)  & +0,0083 (9)  & +0,0106 (10) & +0,0053 (7)  \\
        \midrule
        Ausgangswert (absolut)        & 32,99 & 0,3010 & 0,3184 & 0,3008 \\
        \bottomrule
    \end{tabularx}
    \tabquelle
\end{table}
 
Drei am Datensatz messbare Mechanismen erklären den Gegensatz. Erstens die zeitliche Auflösung: Der Kriminalitätsindex ist das einzige Strukturmerkmal, das sich von Monat zu Monat bewegt, während die sozioökonomischen Merkmale höchstens vier und die baulichen genau einen Wert je Stadtteil annehmen (Abschnitt~\ref{sec:eda}); monatlich variieren sonst nur die Saisonterme, die tragen kein Stadtteilprofil. Zweitens die Kollinearität:Der Index korreliert auf dem Entwicklungspanel mit dem Risikogewerbeanteil zu 0,779 und mit Leerstands- und Armutsquote zu 0,663 und 0,611, ist also ein Sammelindikator; gegeben ihn sind die übrigen Merkmale redundant, und bei 24 Trainingsstadtteilen kosten redundante Prädiktoren mehr Varianz, als sie an Verzerrung abbauen.\footcite[38]{Hastie2009} Drittens der Offset: Das Poisson-GLM führt den Logarithmus der Bevölkerung ohnehin mit, der Prädiktor ist dort ein redundanter freier Parameter, während die Klassifikationsmodelle die Größe als Signal brauchen. Am deutlichsten zeigen es die baulichen Merkmale, die in allen vier Modellen schaden: drei je Stadtteil konstante Koeffizienten, angepasst an die Zwischenvarianz von 24 Stadtteilen. Dass nur ein Jahrgang der Land-Use-Daten vorliegt, ist damit zugleich eine Limitation der Datenlage.
 
Für die erste Unterfrage heißt das: Die Faktorgruppen tragen dieselbe Information mehrfach, welche von ihnen nützt, hängt an der zeitlichen Auflösung der Merkmale und an der Form des Modells. Ein auf den Kriminalitätsindex zusammengestrichener Merkmalssatz träfe deshalb nur den Mengenstrang; im Strukturstrang trägt gerade dieser Index am wenigsten. Der Merkmalssatz bleibt unverändert, weil er in Abschnitt 5.3 vor Kenntnis der Ergebnisse festgelegt wurde und für alle Verfahren derselbe ist --- ihn nachträglich an den Ergebnissen auszurichten, wäre eine ergebnisgetriebene Spezifikationswahl.
% Einflussstärke der Faktorgruppen.                       -> UF1
%
% ===================================================================
% STAND 16.08.2026 -- ERGEBNISSE LIEGEN VOR. Zahlen: docs/03_STAND.md
% Abschnitt 5.6, Analyse: docs/07_BEFUNDE.md B-47. Abbildung A6.
% ===================================================================
%
% DIESER ABSCHNITT HAT ZWEI TEILE, und der zweite widerspricht dem ersten.
% Genau das ist der Befund -- nicht als Widerspruch schreiben, sondern als
% Unterscheidung zwischen zwei verschiedenen Fragen.
%
% TEIL 1 -- ATTRIBUTION: Wie verteilt ein Modell seine Aufmerksamkeit?
%   Quelle: results/shap/gruppen.csv, faktorgruppen_menge.csv (Abbildung A6).
%   Menge (Poisson-GLM, standardisierte Koeffizienten):
%       kriminalitaetsbezogen 36,2 % | baulich 25,6 % | soziooekonomisch 23,2 %
%       Groessenkontrolle 11,2 % | Saison 3,9 %
%   Struktur (SHAP): RF soziooekonomisch 34,9 % / baulich 29,2 % / krimi 16,7 %
%                    XGB soziooekonomisch 50,3 % / krimi 14,8 % / baulich 13,4 %
%   Aussage: Alle drei Faktorgruppen des Exposes tragen bei, keine ist
%   bedeutungslos. Die TREIBER unterscheiden sich zwischen den Straengen --
%   Menge vom Kriminalitaetsindex, Struktur von der Soziooekonomie.
%   WICHTIG: Menge = Koeffizienten, Struktur = SHAP. Zwei verschiedene
%   Groessen, beide auf Summe 1 normiert. Das MUSS dastehen (siehe A6-Fusszeile).
%   Warum die Menge aus der Baseline kommt: Kein Vergleichsverfahren schlaegt
%   sie, Beitraege eines unterlegenen Modells waeren erklaertes Rauschen.
%
% TEIL 2 -- ABLATION: Was ist eine Gruppe WERT?
%   Quelle: results/shap/ablation_faktorgruppen_mittel.csv
%   Jede Gruppe einmal weggelassen, positive Zahl = schlechter ohne sie.
%       weggelassen              Menge(RMSE)  Logit     RF        XGB
%       kriminalitaetsbezogen      +24,27    +0,0008  +0,0004  +0,0020
%       Groessenkontrolle           -6,86    +0,0239  +0,0176  +0,0158
%       soziooekonomisch            -6,33    -0,0131  +0,0032  +0,0180
%       baulich                     -5,44    -0,0112  -0,0041  -0,0058
%       Saison                      +0,24    +0,0036  +0,0045  +0,0083
%
%   Attribution sagt, WIE das Modell schaut. Ablation sagt, WAS FEHLT, wenn
%   man es wegnimmt. Ein Merkmal kann viel Masse binden und ersetzbar sein.
%
% DIE DREI MECHANISMEN -- so und in dieser Reihenfolge schreiben, sie
% erklaeren den Widerspruch VOLLSTAENDIG und machen ihn ungefaehrlich:
%
%   (1) ZEITLICHE AUFLOESUNG. Eindeutige Werte je Stadtteil ueber 132 Monate:
%       log_kriminalitaetsindex 128,6 -- ACS-Merkmale und log_bevoelkerung
%       je 4 (fuenf Jahrgaenge mit Publikationsversatz) -- die drei baulichen
%       Merkmale 1, also KONSTANT, weil Land-Use-Snapshot 2020.
%       Der Kriminalitaetsindex ist das EINZIGE monatlich variierende Merkmal.
%
%   (2) KOLLINEARITAET. Er korreliert mit anteil_risikogewerbe_pct +0,739,
%       leerstandsquote_pct +0,636, armutsquote_pct +0,622,
%       anteil_wohngebaeude_pct -0,603. Er ist ein SAMMELINDIKATOR, der die
%       soziooekonomische und bauliche Information mittraegt -- nur besser
%       aufgeloest. Gegeben ihn sind die uebrigen redundant, und bei 23
%       Trainingsstadtteilen kosten redundante Praediktoren mehr Varianz, als
%       sie Verzerrung abbauen. (Passt zum hoechsten VIF 12,29.)
%
%   (3) DER OFFSET. Groessenkontrolle schadet in der Menge (-6,86) und ist in
%       der Struktur das wertvollste Merkmal (+0,016 bis +0,024). Das
%       Poisson-GLM hat log(Bevoelkerung) ohnehin als OFFSET -- der Praediktor
%       ist dort ein redundanter freier Parameter. Die Klassifikationsmodelle
%       haben keinen Offset und brauchen die Groesse als Signal.
%
% DER KLARSTE EINZELBEFUND, eigener Absatz:
%   Die BAULICHEN Merkmale schaden in ALLEN VIER Modellen und BEIDEN
%   Straengen. Es sind genau die drei, die aus dem Snapshot eines einzelnen
%   Jahres stammen und je Stadtteil konstant sind -- drei Koeffizienten,
%   gefittet auf 23 Stadtteile Zwischenvarianz. Das ist gleichzeitig eine
%   Limitation der Datenlage (nur ein Land-Use-Jahrgang verfuegbar) und ein
%   methodischer Befund.
%
% SO IST UF1 ZU BEANTWORTEN -- und NICHT anders:
%   FALSCH: "Nur der Kriminalitaetsindex traegt." Im Strukturstrang widerlegt,
%           dort ist er praktisch wertlos (+0,0004 bis +0,0020).
%   RICHTIG: Die drei Faktorgruppen des Exposes tragen dieselbe Information
%            MEHRFACH; welche davon nuetzt, haengt an der zeitlichen
%            Aufloesung und an der Modellform.
%
% ZWEI SAETZE, DIE NICHT FEHLEN DUERFEN:
%   - Kein Signifikanztest auf der Ablation. Die Testfamilien sind mit
%     Decision Log #38 festgelegt; weitere Tests wuerden die
%     Holm-Korrekturstruktur beruehren. Berichtet werden Mittelwert, Streuung
%     ueber die zehn Wiederholungsmittel und die Zahl der Wiederholungen mit
%     Verschlechterung (bei kriminalitaetsbezogen: 10 von 10).
%   - Aus der Ablation folgt NICHT, den Merkmalssatz zu kuerzen. Er kommt aus
%     dem Expose und ist durch die Fairness-Regel gebunden; nachtraeglich zu
%     kuerzen waere eine ergebnisgetriebene Spezifikationswahl.

% ---- ÜBERNAHME AUS EXPOSÉ (Kap. 3, Abs. 4) -- Satzfragment ----

% „Zur inhaltlichen Einordnung werden die Modelle mittels SHAP-Werten analysiert, einem
% fundierten Erklärungsansatz für komplexe Modelle."\footcite[17]{Lundberg2017}
% -> Ankündigung, kein Ergebnis. Der Beleg ist verwertbar, der Abschnitt selbst
%    entsteht vollständig neu aus Ihren Analysen.
% ---- ENDE ÜBERNAHME ----

\subsection{Evaluation des Analyseprozesses}
% Rückblick auf den CRISP-DM-Durchlauf; Abgrenzung der Deployment-Phase.

 
Das in Abschnitt 3.2 gesetzte Erfolgskriterium, das Übertreffen der Stufe-2-Baseline, ist im Mengenstrang nicht erfüllt, im Strukturstrang in der Kreuzvalidierung erfüllt und auf dem Hold-out nicht bestätigt. Die Phase Evaluation umfasst auch den Rückblick auf den Analyseprozess und die Festlegung des weiteren Vorgehens.\footcite[30-31]{Chapman2000}
 
Zu diesem Ergebnis führt kein linearer Durchlauf: Vier Befunde der Modellierung haben die Aufbereitung verändert, die Korrektur der Randmonate, die Berücksichtigung des Publikationsversatzes, die einheitliche Expositionskontrolle und die doppelte Stratifizierung der Folds, ohne die einzelne Testmengen keinen brandbestimmten Monat enthielten. Solche Rücksprünge zwischen Data Preparation und Modelling sieht das Vorgehensmodell ausdrücklich vor.\footcite[13]{Chapman2000}

Die Trennung zwischen Suche und Bewertung ist nicht nur konstruktiv gesichert, sondern gemessen. Getunt wird einmal, auf der ersten Wiederholung; in den übrigen neun werden die Folds neu gemischt, sodass dort im Mittel 78 Prozent der Teststadtteile in jener Menge liegen, auf der die Parameter gefunden wurden. Wäre der übertragene Parametersatz ein Leck, müsste der Vorsprung der getunten Verfahren in diesen neun Wiederholungen größer ausfallen als in der ersten. Gemessen ist in allen acht Vergleichen das Gegenteil: Der Abstand fällt dort um 0,098 bis 0,317 Fold-Streuungen kleiner aus. Die Signatur, die ein Leck hinterlassen würde, ist damit nicht vorhanden — und sie wäre sichtbar, weil die Prüfung genau auf sie gerichtet ist.
 
Sämtliche Gütemaße, Hyperparameter, Baselines, SHAP-Beiträge und die Spezifikationsgegenprobe beruhen auf einer festen Startbelegung des Zufallszahlengenerators und reproduzieren exakt. Nicht reproduzierbar sind allein die Laufzeitmessungen.
 
Für die Phase Deployment folgt daraus kein Anlass: Wo die Referenz nicht übertroffen wird, ist der produktive Einsatz eines aufwendigeren Verfahrens nicht begründbar.
% ###########################################################################
% KAPITEL 8 -- VOLLSTAENDIGE NOTIZEN, SCHREIBANLEITUNG UND TABELLEN
% Stand 26.08.2026
%
% EINBAU: ersetzt in main.tex den Bereich von \section{Diskussion} bis
% unmittelbar VOR \section{Fazit und Ausblick} (Stand 26.08.: Zeilen
% 3785-4044). Die beiden Tabellen sind fertiger, kompilierbarer Code; alles
% Uebrige steht als Kommentar und kostet keine Seite.
%
% EINE STELLE AUSSERHALB DIESES BLOCKS IST ANZUFASSEN -- sonst kompiliert die
% Bildunterschrift von Tabelle 8.1 mit einem ?? und einer Warnung:
% Unter \subsection{Forschungsluecke und abgeleitete Erwartungen} (main.tex
% Z. 845) fehlt bisher als einziger Unterabschnitt ein Label. Direkt darunter
% einfuegen:
%     \label{sec:forschungsluecke}
%
% Alle Zahlen in dieser Datei stammen aus results/ des berichteten Laufs
% (#52) und sind am 26.08.2026 gegen die CSV-Dateien geprueft. Wo eine Zahl
% NICHT aus results/ stammt, steht es dabei.
% ###########################################################################


% ===========================================================================
% BLOCK 0 -- DAS SEITENBUDGET. VERBINDLICH.
% ===========================================================================
% Kapitel 8 darf hoechstens 2,5 Seiten belegen. Die folgende Aufteilung ist
% NICHT geschaetzt, sondern am 26.08.2026 mit der Seitengeometrie dieses
% Dokuments nachgemessen (geometry 4/2/4/2 cm, 12pt, onehalfspacing):
%
%   Satzspiegelhoehe        674,3 pt       Zeilenabstand 18,0 pt
%   -> 37,5 Zeilen je Seite, rund 550 Woerter je Seite Fliesstext
%
%   Tabelle 8.1 mit Unterschrift und Float-Abstand      0,25 S.  (gemessen)
%   Tabelle 8.2 mit Unterschrift und Float-Abstand      0,54 S.  (gemessen)
%   1 Ueberschrift Ebene 1 und 3 Ebene 2                0,20 S.
%                                                     ---------
%   fuer Fliesstext bleiben                             1,51 S. = 850 Woerter
%
% VERTEILUNG DER 850 WOERTER -- das ist das eigentliche Budget:
%
%   8.1  3 Absaetze   320 Woerter   (rund 105 je Absatz)
%   8.2  3 Absaetze   280 Woerter   (rund  95 je Absatz)
%   8.3  2 Absaetze   250 Woerter   (rund 125 je Absatz)
%
% Das Schreibrezept verlangt 70 bis 120 Woerter je Absatz -- die Aufteilung
% liegt genau in diesem Fenster. Wer einen Absatz laenger schreibt, nimmt
% ihn einem anderen weg; einen neunten Absatz gibt es nicht.
%
% DREI REGELN, DIE DAS BUDGET TRAGEN -- ohne sie wird es nichts:
%
% 1. JEDE ZAHL STEHT GENAU EINMAL. Entweder in einer Tabelle oder im
%    Fliesstext, nie in beiden. Der Text nennt die RICHTUNG und den GRUND,
%    die Tabelle nennt den WERT. Falsch: "Der Random Forest generalisiert am
%    schlechtesten (R2 0,412 gegen 0,984 im Training)." Richtig: "Der Random
%    Forest generalisiert am schlechtesten und rechnet am laengsten
%    (Tabelle~\ref{tab:erwartung})."
%
% 2. KEINE WIEDERHOLUNG AUS KAPITEL 7. Kapitel 7 berichtet die Ergebnisse,
%    Kapitel 8 deutet sie. Jeder Satz, der eine Zahl aus Kapitel 7 erneut
%    nennt, ohne eine neue Folgerung daraus zu ziehen, wird gestrichen.
%
% 3. LIMITATIONEN GEHEN IN DIE TABELLE, NICHT IN DEN FLIESSTEXT. Der
%    Fliesstext von 8.3 hat zwei Aufgaben und sonst keine: die Groesse der
%    effektiven Stichprobe einordnen und die gemessene Fairnesspruefung
%    berichten. Alles andere ist eine Tabellenzeile.
%
% WENN ES AM ENDE DOCH ZU LANG IST, in dieser Reihenfolge kuerzen:
%   1. Tabelle 8.2 von 13 auf 10 Zeilen -- zuerst fallen "Berichtsumfang",
%      "Schluss auf Personen" und "Gebietszuschnitt" (spart 0,12 S.).
%   2. 8.2 von drei auf zwei Absaetze -- die drei Mechanismen auf zwei
%      Halbsaetze verdichten (spart 0,17 S.).
%   3. Erst danach an 8.1 oder am Fairnessabsatz kuerzen. Beide tragen
%      Bewertungspunkte, die sonst nirgends anfallen.
%
% BLINDABSATZ: zwischen \section{Diskussion} und \subsection 8.1 steht kein
% Text (Vorgabe Schroeter).
%
% FORMALIA (Schroeter 10.08.): keine Kursivschrift, keine
% Anfuehrungszeichen, keine Unterstreichungen, keine Zwischenfazits.
% ===========================================================================


% ===================================================
\section{Diskussion}
\label{sec:diskussion}


% ===========================================================================
% SCHREIBANLEITUNG 8.1 -- VERGLEICHENDE BEWERTUNG DER ALGORITHMEN
% 3 Absaetze, zusammen hoechstens 320 Woerter, danach Tabelle 8.1.
% Beantwortet: Was war zu erwarten, was ist eingetreten, woran liegt es?
% ===========================================================================
%
% ABSATZ 1 -- DIE WIDERLEGTE ERWARTUNG. Gehoert an den ANFANG, nicht ans
% Ende. Abschnitt 2.8 leitet aus Grinsztajn et al. (2022) ab, dass
% baumbasierte Verfahren auf tabellarischen Daten dieser Groessenordnung am
% staerksten abschneiden. Eingetreten ist das Gegenteil: Im Mengenstrang
% schlaegt kein Verfahren die Referenz der Stufe 2, und im Strukturstrang
% kehrt sich das Ergebnis der Kreuzvalidierung auf dem Hold-out um. Das ist
% der interessanteste Befund der Arbeit.
% Erster Satz nennt den Befund, nicht die Erwartung.
%
% ABSATZ 2 -- DIE ERKLAERUNG, UND SIE IST NICHT DIE VERFAHRENSWAHL.
% Das ist die Antwort auf UF4 und das staerkste Argument des Kapitels:
% Zwischen den drei Verfahren liegen 2,2 RMSE, zwischen zwei Spezifikationen
% desselben Verfahrens 19 bis 29. Der Grund ist die multiplikative
% Groessenstruktur: Baumverfahren geben je Blatt einen festen Wert aus und
% koennen "Einsaetze = Bevoelkerung mal Risiko" nicht abbilden; Verfahren mit
% Log-Verknuepfung bekommen die Multiplikation geschenkt.
% ACHTUNG, UEBERHOLTE FASSUNG: Frueher stand hier von einer
% "Spezifikationsasymmetrie" zulasten der Vergleichsverfahren. Seit #43
% modellieren ALLE Verfahren die Rate und multiplizieren zurueck -- die
% Asymmetrie ist konstruktiv beseitigt, nicht nur benannt. NICHT mehr so
% schreiben.
% ZWEITE UEBERHOLTE FASSUNG, ebenfalls nicht mehr schreiben: Der frueher
% vermerkte "schwaechere Rest" der Asymmetrie ist gemessen und existiert
% nicht. B-19: Der Offsetvorteil der Referenz betraegt ueber alle 50 Laeufe
% gepaart -0,0017 RMSE bei der Anzahl und -0,0000 bei der Rate; das negative
% Vorzeichen heisst, die Variante OHNE Offset ist minimal besser. Grund:
% log_bevoelkerung steht ohnehin in PRAEDIKTOREN, der Offset legt nur einen
% Ausgangspunkt fest, den ein frei geschaetzter Koeffizient verschiebt.
% Wenn der Punkt ueberhaupt vorkommt, dann als Fussnote mit dieser Zahl.
%
% ABSATZ 3 -- WARUM FLEXIBILITAET HIER NICHT ZAHLT, UND WAS DAS KOSTET.
% Zwei Haelften, beide gehoeren hierher:
% (a) Die Gegenprobe zur Spezifikation (B-41). In-sample ist Nichtlinearitaet
%     klar nachweisbar; ergaenzt man genau die fehlenden Terme und wertet im
%     selben Rahmen aus, faellt die Prognose um den Faktor drei bis fuenf:
%     linear (12 Terme) RMSE 33,98 -- + Quadrate (22) 101,11 --
%     + Interaktionen (57) 121,63 -- + beides (67) 180,86. Alle 200
%     Anpassungen konvergiert, also kein numerisches Artefakt.
%     Ordnet man die Verfahren nach ungebremster Flexibilitaet, faellt R2
%     monoton: GLM 0,542 -- XGBoost 0,532 -- Ridge 0,511 -- Random Forest
%     0,402 -- Poisson mit Interaktionen, unbestraft, -6,25. Was die
%     Ensembles rettet, ist ihre Regularisierung, nicht ihre
%     Ausdrucksfaehigkeit. Bei 29 unabhaengigen Einheiten ist das das
%     Bias-Varianz-Dilemma in Reinform.
%     DAS IST DER UEBERTRAGBARE METHODISCHE BEITRAG DER ARBEIT: In-sample
%     nachweisbare Nichtlinearitaet ist KEIN hinreichender Grund fuer
%     flexible Verfahren, wenn die Diagnostik geklumpte Daten als unabhaengig
%     behandelt. Ein Satz genuegt, aber er muss stehen.
% (b) Der Preis in Rechenzeit. Der Random Forest braucht 861-mal laenger als
%     Ridge und generalisiert dabei am schlechtesten -- Groessenordnungen
%     statt Prozentpunkte, und damit die belastbarste Aussage des Kapitels.
%     Die Zahlen stehen in der Tabelle, nicht im Satz.
%
% WAS AUS PLATZGRUENDEN NICHT IN DEN FLIESSTEXT KOMMT, ABER IN DIE TABELLE:
% - Unterscheidbarkeit (R-1): keine einzige Verfahrenspaarung ist trennbar.
%   Formulierungsauflage B7 -- NICHT "nicht gezeigt", sondern "nicht
%   gemessen". Sekundaertests Menge: alle sechs nach Holm nicht signifikant;
%   Struktur: RF gegen XGBoost -0,0045, 4 von 10 gewonnen, p 0,625.
% - Ueberanpassung (#51, B6): die R2-Paare Training gegen Kreuzvalidierung.
%   Zu benennende Grenze: Ein Wald mit min_samples_leaf = 1 interpoliert
%   seine Trainingsdaten konstruktionsbedingt; ein hohes Trainings-R2 belegt
%   dort fuer sich genommen keine krankhafte Ueberanpassung. Der Abstand ist
%   zwischen KONFIGURATIONEN und gegen die linearen Modelle zu lesen, nicht
%   als Rangfolge zwischen den Verfahren.
% ===========================================================================

\subsection{Vergleichende Bewertung der Algorithmen}
\label{sec:diskussion-verfahren}

Die Einsatzhäufigkeit eines Stadtteils, den das Modell nicht kennt, lässt sich aus den erhobenen Merkmalen vorhersagen: Der Fehler des Poisson-GLM beträgt weniger als die Hälfte dessen, was der Gesamtmittelwert erreicht, und auf den sechs zurückgehaltenen Stadtteilen erklärt es knapp zwei Drittel der Varianz. Das Erfolgskriterium aus Abschnitt~\ref{sec:zielgroessen} misst nicht dagegen, sondern gegen eine Referenz, die dieselben Merkmale bereits nutzt; sein Verfehlen im Mengenstrang sagt deshalb etwas über die Verfahren und nichts über die Merkmale. Über die Verfahren sagt es dies: Kein Abstand zur Referenz ist trennbar, wohl aber der Abstand der Verfahren untereinander --- und dort liegt in beiden Zielgrößen das sparsamste vorn. Die Erwartung aus Abschnitt~\ref{sec:forschungsluecke} bestätigt sich allein im Strukturstrang, dort nur für den Random Forest und nur in der Kreuzvalidierung; auf dem Hold-out kehrt sich die Reihenfolge um (Tabelle~\ref{tab:erwartung}).

Der Hebel liegt nicht bei der Verfahrenswahl, sondern bei der Spezifikation, und die Gegenüberstellung der beiden Spannweiten in Abschnitt~\ref{sec:belastbarkeit} zeigt, um wie viel: Was zwischen den drei Verfahren liegt, sind nicht einmal zwei Promille dessen, was zwischen der sparsamsten und der reichhaltigsten Spezifikation liegt. Dahinter steht die multiplikative Struktur der Zielgröße: Einsätze entstehen als Produkt aus Bevölkerung und Risiko, ein Modell mit logarithmischer Verknüpfung bildet dieses Produkt ohne weiteres Zutun ab, ein Baum gibt je Blatt einen festen Wert aus. Benachteiligt sind die Vergleichsverfahren dadurch nicht: Ohne die Expositionsbehandlung verlieren Random Forest und XGBoost 15,9 und 14,4 RMSE, während das Poisson-GLM ohne den Bevölkerungsprädiktor sogar besser wird.

Für die vierte Unterfrage, die Implikationen für vergleichbare tabellarische Prognoseaufgaben, folgt daraus eine Regel: Nachweisbare Nichtlinearität im Anpassungsdatensatz begründet noch keine flexiblen Verfahren, solange die Diagnostik geklumpte Beobachtungen wie unabhängige behandelt. Die Eignungsprüfung urteilt über Tausende von Zeilen, entschieden wird die Frage aber über wenige Dutzend voneinander unabhängige Einheiten. Belastbar getrennt haben sich die Verfahren im Aufwand --- dort um Größenordnungen --- und im Mengenstrang zusätzlich in der Güte, und beide Trennungen weisen in dieselbe Richtung: zugunsten des sparsamsten und zulasten des flexibelsten, das zugleich am schwächsten generalisiert. Ein Prognosegewinn, der den Mehraufwand aufwöge, ist in keinem der beiden Stränge nachgewiesen --- die Prognose, der man vertrauen würde, kommt vom billigsten Modell im Feld.

\begin{table}[htbp]
\tabformat
\caption{Erwartung aus Abschnitt~\ref{sec:forschungsluecke} gegen gemessenen Befund}
\label{tab:erwartung}
\begin{tabularx}{\textwidth}{@{}L{2.4cm} L{3.3cm} Y@{}}
\toprule
\textbf{Frage} & \textbf{Erwartung} & \textbf{Befund} \\
\midrule
Menge gegen Referenz & Ensembles übertreffen sie & kein Verfahren übertrifft sie nachweisbar --- und keines bleibt nachweisbar zurück; Ridge liegt numerisch vorn, beide Ensembles dahinter \\
\addlinespace
Struktur gegen Referenz & Ensembles übertreffen sie & in der Kreuzvalidierung nur der Random Forest ($+0{,}017$), XGBoost nicht ($-0{,}000$) --- auf dem Hold-out beide darunter, $-0{,}069$ und $-0{,}059$ Macro-F1 \\
\addlinespace
Verfahren untereinander & Rangfolge erkennbar & teilweise bestätigt, aber gegen die Erwartung: Ridge schlägt beide Ensembles in beiden Zielgrößen, im Strukturstrang schlägt der Random Forest XGBoost \\
\addlinespace
Aufwand & Ensembles teurer & bestätigt, um Größenordnungen: das langsamste Verfahren braucht je Anpassung das 328-Fache des schnellsten \\
\addlinespace
Überanpassung & Ensembles anfälliger & bestätigt: $R^2$ Training gegen Kreuzvalidierung 0,842 zu 0,655 bei Ridge, 0,975 zu 0,486 bei XGBoost, 0,965 zu 0,498 beim Random Forest \\
\addlinespace
Faktoren inhaltlich & Armut, Einkommen und Altbau als stärkste Prädiktoren & Richtung
bestätigt (Tabelle~\ref{tab:zusammenhaenge}), Rang nicht: stärkster Prädiktor ist der
Kriminalitätsindex \\
\bottomrule
\end{tabularx}
\tabquelle
\end{table}


% ===========================================================================
% SCHREIBANLEITUNG 8.2 -- ERKLAERUNGSBEITRAG DER UNTERSUCHTEN FAKTOREN
% 3 Absaetze, zusammen hoechstens 280 Woerter, KEINE eigene Tabelle.
% Beantwortet UF1. Bis 26.08.2026 hatte dieser Abschnitt als einziger
% ueberhaupt keine Schreibanleitung.
% ===========================================================================
%
% SCHNITT GEGEN 7.4 -- OHNE IHN DOPPELN SICH DIE BEIDEN ABSCHNITTE:
%   7.4 zeigt die Zahlen: SHAP-Attribution und Ablation je Faktorgruppe,
%       mit der Tabelle. 8.2 nennt KEINE dieser Zahlen erneut.
%   8.2 zieht die Folgerung: was die Zahlen ueber die Faktoren aussagen,
%       woran ihre Uneinheitlichkeit liegt und was sie NICHT hergeben.
% Ein Rueckverweis auf die Tabelle in 7.4 ersetzt jede Wiederholung.
%
% ABSATZ 1 -- DIE ANTWORT AUF UF1, UND SIE IST NICHT DIE NAHELIEGENDE.
% FALSCH: "Nur der Kriminalitaetsindex traegt." Im Strukturstrang widerlegt,
%         dort ist er praktisch wertlos (+0,0004 bis +0,0020 Macro-F1).
% RICHTIG: Die drei Faktorgruppen des Exposes tragen dieselbe Information
%          MEHRFACH; welche davon nuetzt, haengt an der zeitlichen Aufloesung
%          und an der Modellform.
% Praezisierung, ohne die die Aussage angreifbar ist: Im Querschnitt
% korrelieren mehrere Merkmale aehnlich stark mit der Zielgroesse
% (anteil_risikogewerbe_pct +0,72, anteil_wohngebaeude_pct -0,79,
% leerstandsquote_pct +0,72 gegenueber +0,66 beim Kriminalitaetsindex). Der
% korrekte Satz lautet: GEGEBEN den Kriminalitaetsindex tragen die uebrigen
% Bloecke nichts Eigenstaendiges mehr bei -- nicht "nur Kriminalitaet haengt
% zusammen".
%
% ABSATZ 2 -- DIE DREI MECHANISMEN, die die Uneinheitlichkeit vollstaendig
% erklaeren. Verdichtet auf drei Halbsaetze, die Zahlen stehen in 7.4:
%   (1) ZEITLICHE AUFLOESUNG. Eindeutige Werte je Stadtteil ueber 132 Monate:
%       Kriminalitaetsindex 128,6 -- ACS-Merkmale und log_bevoelkerung je 4
%       (fuenf Jahrgaenge mit Publikationsversatz) -- die drei baulichen
%       Merkmale 1, also KONSTANT (Land-Use-Snapshot 2020). Der
%       Kriminalitaetsindex ist das einzige monatlich variierende Merkmal.
%   (2) KOLLINEARITAET. Er korreliert mit anteil_risikogewerbe_pct +0,739,
%       leerstandsquote_pct +0,636, armutsquote_pct +0,622,
%       anteil_wohngebaeude_pct -0,603 und ist damit ein SAMMELINDIKATOR, der
%       die soziooekonomische und bauliche Information mittraegt, nur besser
%       aufgeloest. Bei 23 Trainingsstadtteilen kosten redundante Praediktoren
%       mehr Varianz, als sie Verzerrung abbauen (hoechster VIF 12,29).
%   (3) DER OFFSET. Die Groessenkontrolle schadet in der Menge und ist in der
%       Struktur das wertvollste Merkmal. Das Poisson-GLM hat
%       log(Bevoelkerung) ohnehin als Offset -- der Praediktor ist dort ein
%       redundanter freier Parameter. Die Klassifikationsmodelle haben keinen
%       Offset und brauchen die Groesse als Signal.
%
% ABSATZ 3 -- DIE EINSCHRAENKUNG, DIE DIE ANTWORT TRAEGT (R-18).
% Das ist der Absatz, der bisher NIRGENDS vorgesehen war, und er ist der
% wichtigste des Abschnitts: log_kriminalitaetsindex ist das einzige Merkmal,
% dessen Weglassen die Prognose messbar verschlechtert (+24,3 RMSE, 10 von 10
% Wiederholungen). Der gesamte inhaltliche Befund zu UF1 haengt an dieser
% einen Spalte -- und sie stammt aus SFPD Incident Reports, also aus
% polizeilich REGISTRIERTEN Delikten. Registrierte Kriminalitaet ist das
% Produkt aus Tatgeschehen, Anzeigebereitschaft und Kontrollintensitaet; wo
% mehr gestreift wird, wird mehr registriert.
% FORMULIERUNGSAUFLAGE, verbindlich:
%   NICHT: "Kriminalitaet erklaert Feuerwehreinsaetze."
%   SONDERN: "Die registrierte Kriminalitaetsbelastung eines Stadtteils ist
%             der staerkste verfuegbare Indikator." Schwaecher formuliert,
%             aber haltbar.
% Belegbar mit Lum & Isaac (2016), Significance 13(5), 14-19 -- kanonischer
% Nachweis der Rueckkopplung zwischen Polizeieinsatz und registrierten
% Delikten. QUELLE IST NOCH NICHT GEPRUEFT, siehe Block Z.
%
% EIN ZWEITER BEFUND, der in einen Halbsatz passt und viel traegt:
% Die BAULICHEN Merkmale schaden in ALLEN VIER Modellen und BEIDEN Straengen.
% Es sind genau die drei, die aus dem Snapshot eines einzelnen Jahres stammen
% und je Stadtteil konstant sind -- drei Koeffizienten, gefittet auf 23
% Stadtteile Zwischenvarianz. Gleichzeitig eine Limitation der Datenlage und
% ein methodischer Befund. Die Zeile dazu steht in Tabelle 8.2.
%
% ZWEI SAETZE, DIE NICHT FEHLEN DUERFEN (ggf. als Fussnote):
% - Kein Signifikanztest auf der Ablation. Die Testfamilien sind mit #38
%   festgelegt; weitere Tests beruehrten die Holm-Struktur. Berichtet werden
%   Mittelwert, Streuung ueber die zehn Wiederholungsmittel und die Zahl der
%   Wiederholungen mit Verschlechterung.
% - Aus der Ablation folgt NICHT, den Merkmalssatz zu kuerzen. Er kommt aus
%   dem Expose und ist durch die Fairness-Regel gebunden; nachtraeglich zu
%   kuerzen waere eine ergebnisgetriebene Spezifikationswahl.
% ===========================================================================

\subsection{Erklärungsbeitrag der untersuchten Faktoren}
\label{sec:diskussion-faktoren}

Die inhaltliche Erwartung aus Abschnitt 2.8 bestätigt sich in ihrer Richtung, nicht in ihrem Rang: Armutsquote und Altbauanteil hängen wie erwartet positiv, das Haushaltseinkommen negativ mit der Einsatzhäufigkeit zusammen (Tabelle 2); als einflussreichster Prädiktor tritt jedoch nicht Armut oder Altbaubestand hervor, sondern der Kriminalitätsindex, und die baulichen Merkmale verschlechtern die Prognose in allen vier Modellen. Beantwortet ist damit auch die dort offen gelassene Frage nach ihrer eigenständigen Erklärungskraft, aber enger, als sie gestellt war. Im Querschnitt hängen mehrere Merkmale ähnlich stark mit der Einsatzhäufigkeit zusammen: der Risikogewerbeanteil zu $+0{,}68$, der Leerstandsanteil zu $+0{,}52$ und der Wohngebäudeanteil zu $-0{,}38$ gegenüber $+0{,}63$ beim Kriminalitätsindex. Eigenständig ist sein Beitrag nicht deshalb, weil die übrigen Merkmale nicht zusammenhingen, sondern weil sie neben ihm nichts Zusätzliches tragen.

Statistisch geprüft ist der Zusammenhang dort, wo die Datenstruktur es zulässt: Für die Einsatzart weisen Kruskal-Wallis-Tests auf den Trainingsstadtteilen des ersten Folds alle zehn Strukturmerkmale als klassentrennend aus (H zwischen 30,6 und 459,6, p < 0,001); für die Einsatzhäufigkeit verbietet der Designeffekt von 122 einen Test auf Zeilenebene, weshalb dort Effektstärken und die gepaarten Verfahrensvergleiche an seine Stelle treten (Abschnitt 7.3).

Der inhaltliche Befund ruht damit auf einer einzigen Spalte, und diese Spalte misst nicht, was ihr Name nahelegt: Sie stammt aus polizeilichen Meldungen und ist das Ergebnis aus Tatgeschehen, Anzeigebereitschaft und Kontrollintensität. Polizeiliche Register sind weder eine Vollerhebung der Delikte noch eine repräsentative Stichprobe daraus; sie bilden eine Wechselwirkung zwischen Kriminalität, Einsatzstrategie und dem Verhältnis zwischen Polizei und Anwohnern ab.\footcite[15-16]{LumIsaac2016} Belegt ist deshalb nicht, dass Kriminalität Feuerwehreinsätze erklärt, sondern dass die registrierte Kriminalitätsbelastung eines Stadtteils der stärkste verfügbare Einzelindikator für seine Einsatzlast ist.

Übertragen auf eine andere Stadt liefert dieser Zusammenhang eine Rangfolge und keinen Zahlenwert: Welche Stadtteile über- und welche unterdurchschnittlich belastet sind, ist aus den Merkmalen bestimmbar, die Einsatzzahl eines einzelnen Stadtteilmonats dagegen nur mit einem mittleren absoluten Fehler von rund 20 Einsätzen. Für die Einsatzart trägt die Übertragung weniger weit, weil dort bereits hier nur etwa die Hälfte dessen ausgeschöpft wird, was mit stadtteilgebundenen Merkmalen überhaupt erreichbar ist. Zwei Voraussetzungen gelten in beiden Fällen: ein Merkmal, das monatlich fortgeschrieben wird, weil nur ein solches die Bewegung innerhalb eines Stadtteils abbildet, und eine Meldepraxis, die mit der hiesigen vergleichbar ist, weil ein polizeilich erhobener Indikator die Kontrollintensität seiner Stadt mitträgt.


% ===========================================================================
% SCHREIBANLEITUNG 8.3 -- LIMITATIONEN UND UEBERTRAGBARKEIT
% 2 Absaetze, zusammen hoechstens 250 Woerter, danach Tabelle 8.2.
% Der Fliesstext hat GENAU ZWEI Aufgaben. Alles andere ist Tabellenzeile.
% ===========================================================================
%
% ABSATZ 1 -- DIE EFFEKTIVE STICHPROBE. Die Zahl lautet: rund 38 effektive
% Einheiten bei 4.620 Zeilen (ICC 0,926, Designeffekt 122). Ohne Saison und
% Kriminalitaetsindex nehmen die Merkmale nur 140 verschiedene Werte an.
% (Diese Zahlen stehen NICHT in results/ -- sie sind eine Varianzzerlegung
% der Praediktoren, B-49. Reproduktionsweg dort.)
%
% DER TON ENTSCHEIDET HIER MEHR ALS DIE ZAHL. Drei Fassungen:
%   FALSCH, defensiv: "Leider standen nur 35 Stadtteile zur Verfuegung."
%       Liest sich als Entschuldigung und laedt zur Nachfrage ein.
%   FALSCH, offensiv: "Viele Arbeiten mit 35 Gebietseinheiten berichten
%       p-Werte auf Basis der Zeilenzahl und merken es nicht."
%       Ein Seitenhieb auf die Literatur OHNE BELEG. Angreifbar, nicht
%       souveraen. NICHT SO SCHREIBEN.
%   RICHTIG, sachlich: Die Groessenordnung als Eigenschaft des Gegenstands
%       benennen, die eigene Behandlung dagegenstellen, beides belegen.
%
% GERUEST, KEINE VORLAGE -- eigener Satzbau, eigene Reihenfolge:
%   "Flaechenbezogene Untersuchungen arbeiten typischerweise mit wenigen
%   Dutzend Gebietseinheiten; die effektive Stichprobe liegt hier bei rund
%   38. Entscheidend ist weniger diese Groesse als der Umgang mit ihr:
%   Getestet wird auf den zehn Wiederholungsmitteln statt auf den 50
%   Einzellaeufen, weil letztere Pseudoreplikation waeren; validiert wird
%   stadtteilweise statt zufaellig, weil zufaelliges k-Fold bei
%   Abhaengigkeitsstrukturen zu optimistische Schaetzungen liefert. Die
%   berichtete Streuung quantifiziert die Abhaengigkeit von der
%   Fold-Zuteilung und ist damit eine UNTERGRENZE der Unsicherheit ueber die
%   Grundgesamtheit der Stadtteile."
%
% ERHALTEN BLEIBEN MUESSEN nur die pruefbaren Behauptungen:
%   - die Zahlen: n_eff rund 38, ICC 0,926, Designeffekt 122, 140
%     Merkmalsvektoren
%   - Hurlbert (1984) an der Stelle "warum n = 10 und nicht 50"
%   - Roberts et al. (2017) an der Stelle "warum stadtteilweise, nicht
%     zufaellig" -- diese Quelle IST geprueft und steht bereits in 6.1
%     zitiert, Seite 925
%   - die Richtung: die berichtete Streuung ist eine UNTERGRENZE
%   - kein unbelegter Seitenhieb auf andere Arbeiten
%
% FALLE BEIM UMFORMULIEREN: Aus "Untergrenze der Unsicherheit" wird leicht
% "das Konfidenzintervall ist konservativ" oder "vorsichtig geschaetzt". Das
% ist das GEGENTEIL und waere falsch. Konservativ hiesse zu breit; hier ist
% es zu SCHMAL, weil alle zehn Wiederholungen dieselben 29 Stadtteile teilen
% (B-50). Sichere Wendungen: "eher zu schmal als zu breit", "die
% tatsaechliche Unsicherheit liegt darueber".
%
% EBENFALLS IN DIESEN ABSATZ, ein Satz, weil er der naheliegenden
% Pruefungsfrage die Spitze nimmt: Eine Ausweitung des Zeitraums wuerde die
% Aussagekraft NICHT erhoehen -- die Beschraenkung liegt bei 35
% Analyseeinheiten, nicht bei der Zahl der Beobachtungen. Ein zusaetzliches
% Jahr braechte 420 Zeilen und null zusaetzliche Stadtteile; unter einem
% Stadtteil-Split besteht die Testmenge aus Stadtteilen.
%
% ------------------------------------------------------------------------
% ABSATZ 2 -- DIE GEMESSENE FAIRNESSPRUEFUNG (R-24, B-52).
% Bis 26.08.2026 war dieser Absatz NIRGENDS vorgesehen. Er gehoert hierher,
% weil das Gutachten nach R6/R9 nicht bewertet, ob Probleme auftreten,
% sondern ob sie methodisch geloest wurden -- und ein gemessener und
% verneinter Verdacht ist staerker als ein ungeprueftes Bedenken.
%
% Der naheliegende Vorwurf: Das Modell ist fuer arme Stadtteile ungenauer,
% benachteiligt sie also. Geprueft am 17.08.2026 ueber alle 50 Laeufe, indem
% je Fold das Sozialprofil der Teststadtteile gegen die Fold-Guete gestellt
% wurde.
%   ABSOLUT besteht der Zusammenhang: Spearman-Rho zwischen Armutsquote und
%   RMSE +0,35 bis +0,60 (p < 0,015), ueber alle Verfahren und beide
%   Zielgroessen, die Poisson-Referenz eingeschlossen.
%   RELATIV ZUM NIVEAU verschwindet er vollstaendig: setzt man den RMSE ins
%   Verhaeltnis zur mittleren Einsatzzahl des Folds, liegt Rho zwischen
%   -0,006 und +0,275, und kein Wert ist auf 5 Prozent signifikant
%   (kleinstes p = 0,053).
%   LESART: Der absolute Fehler ist in aermeren Stadtteilen groesser, weil
%   dort mehr Einsaetze stattfinden. Die relative Genauigkeit ist ueber das
%   Sozialgefaelle hinweg gleich. Es liegt kein Hinweis auf systematische
%   Benachteiligung vor.
%
% ZWEI EHRLICHKEITEN, die dazugehoeren und nicht wegfallen duerfen:
%   1. Der relative Fehler ist auf allen Stufen hoch -- im Mittel 0,48 bei
%      der Anzahl und 0,66 bis 0,69 bei der Rate. Das Modell ist nirgends
%      praezise, nicht nur in armen Stadtteilen.
%   2. Es ist eine Pruefung auf FehlerGLEICHHEIT, nicht auf
%      Verteilungsgerechtigkeit eines hypothetischen Einsatzes. Die waere
%      eine andere Frage und wird hier nicht beantwortet.
%
% WICHTIG: Diese Zahlen liegen NICHT in results/. Reproduktionsweg in B-52:
% je (Wiederholung, Fold) die mittlere Armutsquote der Teststadtteile
% berechnen und gegen die Fold-Guete stellen (Spearman, 50 Laeufe).
% Vor der Verwendung im Text einmal nachrechnen und die Zahlen ablegen.
%
% ANSCHLUSS IN EINEM HALBSATZ, wenn der Platz reicht (R-20): Die Abgrenzung
% der Deployment-Phase ist bisher nur eine Umfangsentscheidung; sie ist
% zusaetzlich als ethische Grenze zu formulieren. Das unterscheidet die
% Arbeit von der Predictive-Policing-Kritik, statt ihr auszuweichen: Ein
% zusaetzlicher Loeschzug NUETZT einem Stadtteil, zusaetzliche Streifen
% BELASTEN ihn. Zuteilung nach vorhergesagtem Bedarf ist in der
% Notfallversorgung nicht per se ungerecht, sondern das Ziel.
% Wenn der Platz nicht reicht: in den Ausblick, Kapitel 9.
% ===========================================================================

\subsection{Limitationen und Übertragbarkeit}
\label{sec:limitationen}

Die Grenze, die alle übrigen bindet, ist die Zahl unabhängiger Einheiten. Dass
36 Stadtteile 4.752 Zeilen ergeben, hilft nur dort, wo sich die Merkmale
innerhalb eines Stadtteils bewegen, und mit Ausnahme des Kriminalitätsindex tun
sie das kaum; ein längerer Beobachtungszeitraum ändert daran nichts, denn ein
zusätzliches Jahr brächte 432 Zeilen und keinen einzigen weiteren Stadtteil. Wie
unter der in Abschnitt~\ref{sec:belastbarkeit} bezifferten effektiven Stichprobe
geprüft werden darf, ist dort ausgeführt; die berichtete Streuung misst die
Abhängigkeit von der Fold-Zuteilung und ist als Unsicherheit über die
Grundgesamtheit der 36 Stadtteile eher zu schmal als zu breit.

Damit fällt auch die Abwägung, die Abschnitt~\ref{sec:rahmen-baselines} an dieser Stelle ankündigt. Der gewählte Rahmen prüft, was die Forschungsfrage stellt: die Vorhersage für einen Stadtteil, den das Modell nicht kennt. Der ursprünglich angekündigte Zeitschnitt hätte stattdessen die Fortschreibung bereits bekannter Stadtteile gemessen und wäre an der räumlichen Konzentration der Varianz gescheitert. Was der Wechsel kostet, ist genau diese Fortschreibung: Ob die geschätzten Zusammenhänge in einigen Jahren noch gelten, prüft diese Arbeit nicht.

Eine naheliegende Sorge lässt sich messen, aber nur teilweise ausräumen. Stellt
man je Fold das Sozialprofil der Teststadtteile gegen die dort erreichte Güte,
besteht der Zusammenhang absolut für jedes Verfahren: Zwischen Armutsquote und
RMSE liegt der Rangkorrelationskoeffizient zwischen $+0{,}39$ und $+0{,}81$, das
größte $p$ beträgt 0,0052. Das allein besagt wenig, denn in ärmeren Stadtteilen
finden mehr Einsätze statt, und ein größerer absoluter Fehler bei gleicher
relativer Genauigkeit ist keine Benachteiligung. Ins Verhältnis zur mittleren
Einsatzzahl des Folds gesetzt, verschwindet er bei den beiden Baumverfahren, bei
den beiden linearen Modellen dagegen nicht: Das Poisson-GLM behält in beiden
Zielgrößen einen signifikanten Rest, Ridge bei der Einsatzhäufigkeit einen
schwächeren. Ausgerechnet die Modelle, die im Mittel am genauesten
prognostizieren, verteilen ihren Fehler am ungleichmäßigsten über das
Sozialprofil. Geprüft ist damit die Gleichheit der Fehler, nicht die
Gerechtigkeit einer daraus abgeleiteten Zuteilung --- und für die linearen
Modelle ist sie nicht gegeben.

\begin{table}[htbp]
\tabformat
\caption{Grenzen der Untersuchung und ihre Wirkung auf die Aussagekraft}
\label{tab:limitationen}
\begin{tabularx}{\textwidth}{@{}L{2.6cm} L{4.8cm} Y@{}}
\toprule
\textbf{Grenze} & \textbf{Beleg} & \textbf{Wirkung auf die Aussage} \\
\midrule
Extrapolation & 36,6 \% der Testzeilen über alle fünfzig Läufe außerhalb des gelernten Bereichs, im Hold-out 34,8 \% & begrenzt die Übertragung auf unbekannte Stadtteile in beiden Auswertungen gleichermaßen \\
\addlinespace
Parameterwahl & getunt wird einmal, auf der ersten Wiederholung; die Referenzen sind ungetunt & schwacher Informationsübertrag, alle Verfahren gleichermaßen betroffen \\
\addlinespace
Seltene Klasse & 33 brandbestimmte Monate, 2 bis 13 je Testfold & ein Viertel des Macro-F1 ruht auf wenigen Fällen \\
\addlinespace
Zeitliche Auflösung & ACS jährlich mit Publikationsversatz, Flächennutzung als Momentaufnahme & die baulichen Merkmale verschlechtern alle vier Modelle \\
\addlinespace
Gebietszuschnitt & 36 Analysis Neighborhoods statt 242 Census Tracts & eine andere Einteilung kann eine andere Rangfolge liefern \\
\addlinespace
Verfahrenszahl & zwei statt drei Verfahren in der Klassifikation & verglichen werden zwei einander ähnliche Baumverfahren \\
\addlinespace
Expositions\-gewichtung & die Referenz gewichtet über den Offset nach Einwohnerzahl, die
drei Verfahren nicht & betrifft beide Auswertungen; ungewichtete Raten- und
expositionsgewichtete Zählskala führen zur selben Rangfolge \\
\bottomrule
\end{tabularx}
\tabquelle
\end{table}

Die übrigen Grenzen sind zusammengefasst (Tabelle~\ref{tab:limitationen}); drei
von ihnen bezeichnen zugleich den nächsten Arbeitsschritt: die feinere
Gliederung der Census Tracts, ein in gleicher Frequenz fortgeschriebener
Merkmalssatz und eine eigene Hyperparametersuche je Wiederholung. Dass die
Deployment-Phase ausgeklammert bleibt, ist dabei nicht allein eine Frage des
Umfangs: Ein Modell, das Einsatzlast vorhersagt, verteilt Hilfe und nicht
Kontrolle --- ein zusätzlicher Löschzug nützt einem Stadtteil, zusätzliche
Streifen belasten ihn ---, und Zuteilung nach vorhergesagtem Bedarf ist in der
Notfallversorgung deshalb kein Fehler, sondern das Ziel, solange der Bedarf
nicht über einen Indikator geschätzt wird, dessen Höhe von der Zuteilung selbst
abhängt.

Die Blöcke folgen dabei der Verwaltungsgrenze und nicht der Reichweite der räumlichen Abhängigkeit — ein Puffer um die Teststadtteile ließe im Mittel 8,6 statt 24 Trainingsstadtteile übrig —, sodass der berichtete Fehler eher zu klein als zu groß ist.

Die Exposition ist die Wohnbevölkerung, und für die Innenstadt ist sie unvollständig: Financial District/South Beach hat mit 20.044 Einwohnern eine Wohnbevölkerung im Median aller Stadtteile, aber mit 10,40 Einsätzen je 1.000 Einwohner die höchste Rate der Stadt — das Vierfache des Medians. Wer dort tagsüber ist, steht in keinem der verwendeten Register.


% ###########################################################################
% BLOCK Z -- WAS NICHT IN KAPITEL 8 GEHOERT, WAS OFFEN IST, WAS FEHLT
% Kein Text, nur Merkposten. Vor der Abgabe abarbeiten.
% ###########################################################################
%
% ---------------------------------------------------------------------------
% Z1. AUSDRUECKLICHE GEGENANWEISUNGEN -- NICHT NACH KAPITEL 8
% ---------------------------------------------------------------------------
% - DIE DECKEN (B-48). Gehoeren VOR die Ergebnistabelle in 7.2, nicht in die
%   Limitationen -- sonst lesen sie sich als nachtraegliche Entschuldigung
%   statt als Massstab. Werte: Mehrheitsklasse 0,2233 -- Decke B
%   (Stadtteilwissen) 0,4572 -- Decke A (Label-Rauschen) 0,6404 -- fehlerfrei
%   1,0. Hold-out: Decke A 0,6785, Decke B 0,4219.
% - B-45 (weitere Suchraeume verschlechtern die Prognose) traegt den Vermerk
%   NUR INTERN, geht NICHT in die Arbeit. Nach #52 wird genau ein Lauf
%   berichtet, kein Vorher-Nachher. Uebernehmbar ist allein die allgemeine
%   Lehre, und die ist B-41 (steht in 8.1).
% - Der Seitenhieb auf die Literatur zur effektiven Stichprobe: ohne Beleg,
%   ausdruecklich nicht zu schreiben.
% - Die Extrapolation erklaert KEINEN Verfahrensunterschied (B-31, gemessen
%   und widerlegt: der Zusammenhang zwischen Extrapolationsanteil und RMSE
%   ist bei RIDGE am staerksten, Spearman +0,302, beim Random Forest am
%   schwaechsten, +0,108 bei p 0,456). Sie bleibt eine Eigenschaft des
%   Validierungsrahmens. Diese Erklaerung darf in 8.1 NICHT auftauchen.
%
% ---------------------------------------------------------------------------
% Z2. AUS PLATZGRUENDEN GESTRICHEN -- bewusst, nicht vergessen
% ---------------------------------------------------------------------------
% Diese Punkte sind fuer Kapitel 8 vorgemerkt, passen aber nicht in 2,5
% Seiten. Wenn Platz entsteht, in dieser Reihenfolge aufnehmen:
%   1. Die Anpassungsskala (Alle drei Vergleichsverfahren werden auf der Rate
%      angepasst, das Poisson-GLM auf der Anzahl mit Offset; gemessen wird
%      symmetrisch, angepasst nicht). Zwei Gegenproben liegen vor: im
%      Ratenstrang haben die Baeume die Ausrichtung zwischen Verlust und
%      Metrik und verlieren trotzdem (Hold-out RMSE GLM 0,977 gegen XGBoost
%      1,137, Ridge 1,153, Random Forest 1,561); die Expositions-Ablation
%      kostet den Random Forest 29,2 und XGBoost 19,1 RMSE, die gewaehlte
%      Spezifikation ist also die fuer die Vergleichsverfahren guenstigere.
%      Eine expositionsgewichtete Variante wurde NICHT erhoben -- #43 legt
%      eine Spezifikation fuer alle fest, nachtraeglich weitere zu messen und
%      unter ihnen zu waehlen waere ergebnisgetrieben.
%   2. Die Rueckkopplung bei hypothetischem Einsatz (R-23): mehr Praesenz,
%      mehr registrierte Ereignisse, Bestaetigung der Vorhersage. Beim Brand
%      schwaecher als bei der Polizei, bei Fehlalarmen (79 Prozent, dominante
%      Klasse) denkbar. Laut Register gehoert das ohnehin in den AUSBLICK,
%      Kapitel 9, nicht in Kapitel 8.
%   3. Das fehlende lineare Referenzverfahren der Klassifikation (#31): Die
%      Symmetrie linear/Bagging/Boosting besteht nur im Regressionsstrang;
%      der Vergleich in der Klassifikation ist einer zwischen zwei
%      Baumverfahren, die einander konstruktionsbedingt aehnlicher sind.
%      Steht bereits in 6.2 -- in 8.3 nur, wenn Platz ist.
%   4. Die Reflexion der beiden Dokumentationsfehler (B-43/B-44): eine
%      Information lebte an zwei Orten, nur einer wurde gepflegt. Trifft
%      Gutachten-R9 (Reflexion im Verlauf). ENTSCHEIDUNG OFFEN, ob ein selbst
%      berichteter Prozessfehler ueberhaupt in die Arbeit soll.
%
% ---------------------------------------------------------------------------
% Z3. ZAHLENKONFLIKTE -- VOR DEM SCHREIBEN KLAEREN
% ---------------------------------------------------------------------------
% (a) RESET-TEST, ZEILENZAHL. 06_RISIKEN.md (Z. 392), 07_BEFUNDE.md B-41 und
%     der Kommentar in main.tex Z. 2456 nennen 3.828 Zeilen. Richtig sind
%     3.036 -- die Eignungspruefung lief auf den 23 Trainingsstadtteilen von
%     Fold 1 (results/eignungspruefung/eignungspruefung.md, Z. 3). 3.828 ist
%     die TRAININGSGROESSE DES HOLD-OUTS (29 mal 132, holdout.csv). Der
%     Fliesstext in 6.2 (main.tex Z. 2568-2571) ist bereits richtig; die drei
%     Kommentar- und Notizstellen sind nachzuziehen.
% (b) STADTTEIL-MITTELWERT. #29 nennt R2 0,888, die Praezisierung vom
%     04.08.2026 und 06_RISIKEN.md nennen 0,925. Genau diese Zahl traegt
%     Auflage A. Vor der Verwendung festlegen, welche gilt.
% (c) VORSPRUNG DER ENSEMBLES IN DER STRUKTUR. #48 und 07_BEFUNDE B-29 nennen
%     +0,0304/+0,0371 bzw. +0,0296/+0,0362 -- beides Werte vom 06.08. Der
%     BERICHTETE Lauf (#52) liefert +0,0306 und +0,0350
%     (klassifikation/vergleich.csv). In der Arbeit gelten die letzteren.
% (d) DIE GETUNTE REFERENZ 0,314 stammt aus der Messung vom 06.08. und wurde
%     nach #48 bewusst NICHT nachgerechnet. Sie ist damit eine Zahl aus einem
%     anderen Lauf als die uebrigen. Beim Schreiben so kennzeichnen.
% (e) DISPERSIONSINDEX. annahmen.csv nennt 54,2, #42 nennt 62,8. Es gilt
%     54,2 (berichteter Lauf, Fold-1-Trainingsstadtteile).
%
% ---------------------------------------------------------------------------
% Z4. QUELLEN, DIE DIESES KAPITEL BRAUCHT
% ---------------------------------------------------------------------------
% Roberts et al. (2017) ist geprueft, im Literaturverzeichnis und in 6.1 mit
% Seite 925 zitiert -- die einzige der folgenden, die verwendbar ist.
%
% NICHT GEPRUEFT, kein bib-Eintrag, keine Seitenangabe. Nach der geltenden
% Regel darf keine davon vor der Volltextpruefung in den Text:
%   Hurlbert, S. H. (1984): Pseudoreplication and the Design of Ecological
%     Field Experiments. Ecological Monographs 54(2), 187-211.
%     -> braucht 8.3 Absatz 1 fuer "warum n = 10 und nicht 50"
%   Lum, K. & Isaac, W. (2016): To predict and serve? Significance 13(5),
%     14-19.  -> braucht 8.2 Absatz 3
%   Robinson, W. S. (1950): Ecological Correlations and the Behavior of
%     Individuals. American Sociological Review 15(3), 351-357.
%     -> staerkt die Tabellenzeile "Schluss auf Personen"
%   Openshaw, S. (1984): Ecological Fallacies and the Analysis of Areal
%     Census Data. Environment and Planning A 16(1), 17-31.
%     -> staerkt die Tabellenzeile "Gebietszuschnitt"
%   Nadeau, C. & Bengio, Y. (2003): Inference for the Generalization Error.
%     -> staerkt die Tabellenzeile "Abhaengige Wiederholungen"
%
% ENTLASTUNG: Die beiden Tabellenzeilen zum Gebietszuschnitt und zum Schluss
% auf Personen sind Aussagen ueber das EIGENE Design und tragen auch ohne
% Zitat. Nur Hurlbert und Lum & Isaac werden wirklich gebraucht; ohne sie ist
% der jeweilige Satz abzuschwaechen, nicht zu streichen.
%
% ---------------------------------------------------------------------------
% Z5. SCHREIBREIHENFOLGE UND ABHAENGIGKEITEN
% ---------------------------------------------------------------------------
% - Abschnitt 2.8 (Forschungsluecke) wird ZULETZT geschrieben, nach 8.1. Die
%   dort abgeleiteten Erwartungen sind genau die, die 8.1 widerlegt -- sie
%   muessen zusammenpassen (09_SCHREIBPLAN.md Z. 74).
% - AUFLAGE A (Schroeter, 04.08.2026, woertlich): "Wichtig ist, dass Sie die
%   Abweichung von der urspruenglich angekuendigten zeitreihengerechten
%   Kreuzvalidierung transparent erlaeutern." Das begruendete Verwerfen ist
%   VERLANGT, nicht optional. Es steht als letzte Zeile in Tabelle 8.2 -- die
%   Zeile allein erfuellt die Auflage nur, wenn der Beleg mit dranhaengt:
%   92,5 Prozent der Varianz liegen zwischen den Stadtteilen, ein Zeitschnitt
%   wuerde im Kern die zeitliche Stabilitaet messen und die Forschungsfrage
%   nicht pruefen. Pruefen, ob dieser Beleg in 5.4 schon steht -- wenn ja,
%   genuegt in 8.3 der Verweis; wenn nein, muss er in die Tabellenzeile.
% - DARSTELLUNGSREGEL (Schroeter, 10.08.2026): keine grosse iterative
%   Entwicklung im Text; Kapitel 6 beschreibt die finale Spezifikation, die
%   Reflexion ist KONZENTRIERT IN KAPITEL 8. Damit gehoeren #42
%   (Verlustfunktionen) und #43 (Exposition) hierher -- beide ausdruecklich
%   NACH einem Modelllauf getroffen und in der Arbeit als solches zu
%   kennzeichnen. Bei 2,5 Seiten reicht dafuer ein Halbsatz in 8.1 Absatz 2.
% - KI-VERZEICHNIS: Wird eine hier vorformulierte Passage als Geruest
%   genutzt, gehoert ein Eintrag in die Tabelle am Ende des Dokuments, nach
%   dem Muster der vorhandenen Eintraege mit Prompt und Ergebnisverweis.
% - GLOBALE UMBENENNUNG (Schroeter, 24.08.2026): Quelltext/Quellcode wird zu
%   "Code 1: xy". Betrifft auch die \ref-Aufrufe in diesem Kapitel, falls
%   welche dazukommen.
%
% ---------------------------------------------------------------------------
% Z6. PRUEFUNG NACH DEM SCHREIBEN
% ---------------------------------------------------------------------------
% 1. Zweimal kompilieren (Verzeichnisse stabilisieren sich erst im zweiten
%    Lauf), dann die Seitenzahl von \section{Diskussion} bis
%    \section{Fazit und Ausblick} nachmessen. Ziel: hoechstens 2,5.
% 2. Beide Tabellen muessen auf ihrer Seite stehen bleiben. Faellt eine um,
%    ist der Float an den ANFANG seines Blocks zu setzen, nicht ans Ende --
%    [H] allein schiebt sie nur auf die naechste Seite.
% 3. Jede Zahl im Fliesstext gegen die Tabellen pruefen: steht sie zweimal,
%    faellt die im Fliesstext weg.
% 4. Gegen Kapitel 7 lesen: Wird eine Ergebniszahl wiederholt, ohne dass eine
%    neue Folgerung daraus entsteht, streichen.
% ###########################################################################
% ==================================================
% ---------------------------------------------------
% LITERATURVERZEICHNIS
% ---------------------------------------------------
\newpage
% \nocite{*} ENTFERNT am 18.08.2026. Es zog jeden Datenbankeintrag ins
% Literaturverzeichnis, auch ungenutzte und solche, die nur zur Pruefung
% vorgemerkt sind. Gedruckt wird jetzt ausschliesslich, was zitiert wird.
\printbibliography[heading=bibintoc,title={Literaturverzeichnis}]

% ---------------------------------------------------
% KI-VERZEICHNIS
% ---------------------------------------------------
\newpage
\phantomsection
\section*{KI-Verzeichnis}
\addcontentsline{toc}{section}{KI-Verzeichnis}

\begin{table}[H]
    \tabformat
    \caption*{Verzeichnis der verwendeten KI-Hilfsmittel}
    \begin{tabularx}{\textwidth}{@{}L{0.8cm} L{2.8cm} Y@{}}
        \toprule
        \textbf{Nr.} & \textbf{KI-System} & \textbf{Einsatzzweck / Prompt} \\
        \midrule
        1 & Anthropic Claude & \enquote{...} \\
        \addlinespace
        2 & Google Gemini & \enquote{...} \\
        \bottomrule
    \end{tabularx}
    \tabquelle
\end{table}

\end{document}